\documentclass[11pt]{article}
\usepackage[T1]{fontenc}
\usepackage[utf8]{inputenc}
\usepackage{lmodern}
\usepackage[margin=1in]{geometry}
\usepackage{amsmath,amssymb,amsthm,mathtools,bm}
\usepackage{booktabs,array,longtable}
\usepackage{enumitem}
\usepackage{hyperref}
\usepackage{xurl}
\input{release_info.tex}

\hypersetup{colorlinks=true,linkcolor=blue,citecolor=blue,urlcolor=blue}
\setlength{\emergencystretch}{3em}
\setlength{\tabcolsep}{5pt}
\renewcommand{\arraystretch}{1.08}

\newtheorem{theorem}{Theorem}[section]
\newtheorem{corollary}[theorem]{Corollary}
\newtheorem{proposition}[theorem]{Proposition}
\newtheorem{lemma}[theorem]{Lemma}
\newtheorem{definition}[theorem]{Definition}
\newtheorem{axiom}[theorem]{Axiom}
\newtheorem{assumption}[theorem]{Assumption}
\newtheorem{remark}[theorem]{Remark}

\newcommand{\SU}{\mathrm{SU}}
\newcommand{\U}{\mathrm{U}}
\newcommand{\Conf}{\mathrm{Conf}}
\newcommand{\nf}{\operatorname{nf}}
\newcommand{\ellbar}{\bar{\ell}}
\newcommand{\ophid}[1]{\path{#1}}

\title{Recovering Relativity and the Standard Model\\
from the OPH Package Rooted in Observer Consistency}
\author{B.\ M\"uller \and Kai Xue \and Peter Nguyen}
\date{\OPHPaperReleaseDate}

\begin{document}
\maketitle
\OPHPaperReleaseBanner

\begin{abstract}
This paper gives the recovered-core relativity/Standard-Model derivation in Observer-Patch Holography (OPH). Starting from local observer patches on a finite-capacity screen, overlap consistency, and an explicit premise ledger, it reconstructs a schedule-independent overlap normal form, a conditional Lorentz branch, a conditional route to general relativity through a Jacobson-type gravity branch, and a conditional route to the realized Standard Model quotient \(\left(\SU(3)\times\SU(2)\times\U(1)\right)/\mathbb Z_6\). It also derives the exact rational hypercharge lattice and the counting chain \(N_g=3\), \(N_c=3\) for the realized low-energy matter package. The paper separates theorem-grade fixed-cutoff structure from scaling-limit assumptions, calibration sectors, and phenomenological continuations. The same declared local input \(P\) organizes the electroweak/Higgs trunk and, on the declared exact-release extension surface, the gravity coupling \(G\).
\end{abstract}

\section{Introduction}

The standard model of particle physics and general relativity are successful effective descriptions. They leave open a familiar list of questions: why the Lorentz group appears, why the low-energy spacetime description follows a Jacobson-type Einstein branch, why the gauge structure is
\[
\frac{\SU(3)\times\SU(2)\times\U(1)}{\mathbb Z_6},
\]
why the observed three generations and three colors are realized, why hypercharges take the observed rational values, what fixes the electroweak and flavor scales, and what the framework can and cannot say about the cosmological constant and possible dark-sector continuations.

This paper presents a single framework that addresses this list as a conditional reconstruction program. The framework is Observer-Patch Holography (OPH). Its starting point is an observer-centric organization of physical data: each observer has access only to a patch algebra, neighboring observers must agree on overlaps, and the OPH reconstruction program is organized around the resulting consistency conditions. Within the full OPH axiom package and the additional branch premises stated below, the same formal structure that produces schedule-independent overlap consistency then feeds a Lorentzian modular-phase classification, a Jacobson-type Einstein branch conditional on fixed-cap stationarity and the null modular bridge together with the bounded-interval projective branch, compact gauge reconstruction in a bosonic internal-gauge sector, charge quantization, and the low-energy architecture selected by the admissibility conditions used in the gauge branch. Any discrete-horizon or black-hole spectroscopy template remains a Phase-III continuation and is not part of the recovered relativity-plus-Standard-Model core.

The paper separates exact structural theorems, conditional scaling-limit branches, calibration-sector outputs, and phenomenological continuations.\footnote{Broader interpretive questions are deferred here so the claim-tier boundary remains explicit.} Several longer arguments are recorded below as proof sketches.

The framework is intended as a candidate fundamental theory in the precise sense relevant for physics: general relativity and the Standard Model appear as effective sectors derived from a smaller axiom set plus an explicit refinement/scaling-limit package. The paper does not use simulation language, observer continuation stories, or other non-falsifiable additions.

\subsection*{Prioritized achievements}

The results emphasized in this paper are ordered below by a combination of impact and defensibility.

\begin{enumerate}[leftmargin=2em]
\item A conditional Jacobson-type Einstein branch is recovered from observer-overlap modular geometry, an explicit fixed-cap generalized-entropy stationarity premise, the geometric-branch selection burden, and an explicit null-modular scaling branch; in the null modular bridge, quasi-local propagation, endpoint control, half-sided inclusion, the explicit positive null-translation generator, and the half-line generator/charge identification are supplied internally through the derived Borchers stage, while bounded-interval transport through the separate projective branch and the later tensor upgrade remain explicit downstream burdens.
\item The exact Standard Model gauge structure
\[
\frac{\SU(3)\times\SU(2)\times\U(1)}{\mathbb Z_6}
\]
is recovered in the bosonic compact-gauge branch once the MAR admissibility package is imposed.
\item On the realized one-generation chiral matter plus one-Higgs package, anomaly constraints and Yukawa invariance fix the hypercharge lattice, and on the realized MAR-admissible branch the counting chain yields \(N_g=3\) and then, using Witten's \(\SU(2)\) anomaly parity constraint, \(N_c=3\).
\item The quantitative implementation used here uses two external continuous configuration inputs, $P$ and $N_{\mathrm{scr}}$.
\item Downstream matter-sector continuations, including the charged-lepton centered-readback / centered-to-affine frontier, remain visible but explicitly outside this SM/GR derivation paper's recovered-core theorem package.
\item A separate input-dependent capacity corollary ties the cosmological-constant branch to global screen capacity rather than local vacuum energy.
\item Exclusion of gauge-mediated proton decay follows from product-group structure; any stronger proton-lifetime claim remains UV/EFT/QCD-dependent.
\item Dark-sector, baryogenesis, and spectroscopy branches are kept visible only as phenomenological continuations or, where additional ans\"atze are involved, conjectural phenomenology. On the string/worldsheet lane, the fixed-cutoff heat-kernel / Yang--Mills form stays on the declared compact-gauge branch, any worldsheet-like reorganization requires a controlled large-\(N_{\mathrm{edge}}\) regime distinct from that proved chain, and any critical-superstring lift is conjectural.
\item Overlap repair admits a unique schedule-independent normal form on the stated finite patch-net hypotheses.
\item The fixed-cutoff topological UV package closes on the ordinary branch, the central-defect branch, and the genuinely noncentral branch through a compact crossed-module higher-gauge collar theorem; the remaining UV-side burdens are the continuum BW/geometric lift and the refinement-limit compact-gauge premises.
\end{enumerate}

\section{Overview of Results}

This section summarizes the claim tiers used throughout. The labels \(D1\)--\(D12\) are internal
documentary node labels for the OPH dependency ledger, not bibliography references or external
notation. The public paper ledger has no separate D11 row; the other OPH papers should use
only the labels defined in the ledger below. The table below is the local key for this
SM/GR derivation paper, and the other OPH papers reuse the same labels when they point back to this ledger. The
three-phase boundary is
\[
\text{Phase I}=(D1\text{--}D5)\cup(D7\text{--}D9),\qquad
\text{Phase II}=D6\cup D10,\qquad
\text{Phase III}=D12.
\]
Phase I is the recovered-core claim set. Phase II contains the separate input-dependent capacity corollary and the integrated D10 gauge-calibration sector. Phase III collects phenomenological continuations. The canonical five-axiom basis, the two external continuous inputs, the scaling/continuum premises, and the gauge-side assumptions used below are fixed in Section~\ref{sec:premise-ledger}. Each row names a node in the reconstruction program, records its immediate OPH-side parents, separates imported standard mathematics from extra non-axiom physical premises or external inputs, and states the resulting claim tier.

\begingroup
\tiny
\setlength{\tabcolsep}{1pt}
\begin{longtable}{@{}>{\raggedright\arraybackslash}p{0.05\textwidth}>{\raggedright\arraybackslash}p{0.16\textwidth}>{\raggedright\arraybackslash}p{0.14\textwidth}>{\raggedright\arraybackslash}p{0.13\textwidth}>{\raggedright\arraybackslash}p{0.20\textwidth}>{\raggedright\arraybackslash}p{0.16\textwidth}@{}}
\caption{Dependency checklist for the OPH SM/GR reconstruction program.}\\
\toprule
Node & Output & Immediate OPH ingredients & Standard mathematics used & Extra non-axiom physical premises / external inputs & Claim tier \\
\midrule
\endhead
\bottomrule
\endfoot
\textbf{D1} & Unique schedule-independent normal form on the physical quotient, together with observable-level confluence on the declared fixed-cutoff physical algebras for the recovery-derived repair relation & overlap-consistency problem on a finite patch net together with the declared fixed-cutoff collar recovery step, its touched-overlap local-fit acceptance contract, support-local disjoint commutation, restriction-compatible union-collar gluing on the physical quotient, and the fixed-point / quotient-local physical observable algebras carried by that same collar package & local-to-global Lyapunov descent on the inconsistency potential; local-diamond completion from overlap-associative center-sector / higher-gauge union-collar gluing; Newman's lemma & repair completeness; stated Petz support/CPTP control where that branch is used & Phase I structural theorem \\
\textbf{D2} & Collar Markov/entropy split and \(G\)-dictionary setup (Theorem~\ref{thm:markovcollar}, Propositions~\ref{prop:approxrecover}, \ref{prop:splice}, \ref{prop:gensplit}, Definition~\ref{def:newton}) & Axioms~\ref{ax:screen}--\ref{ax:entropy} & HJPW Markov structure theorem; Petz and Fawzi--Renner recovery theory & exact block-factorized identities require exact Markovity or an explicitly stated idealized limit; the \(A/(4G)\) step is the coarse-grained dictionary of Definition~\ref{def:newton} & Phase I structural lemma/proposition layer \\
\textbf{D3} & Lorentz branch \(\Conf^+(S^2)\cong \mathrm{SO}^+(3,1)\) (Theorem~\ref{thm:bw}, Corollary~\ref{cor:lorentz}) & Axioms~\ref{ax:screen}--\ref{ax:entropy} & Bisognano--Wichmann modular geometry; conformal-group classification & T2 together with the theorem-local geometric cap-pair extraction and ordered cut-pair rigidity hypotheses of Theorem~\ref{thm:bw}; Theorem~\ref{thm:geosubnet} fixes the target subnet internally & Phase I conditional scaling-limit theorem \\
\textbf{D4} & Null modular bridge to \(T_{kk}\) (Propositions~\ref{prop:nullec}, \ref{prop:n3}; Corollaries~\ref{cor:n1}, \ref{cor:n2}; Theorem~\ref{lem:density}, Lemma~\ref{lem:translation}; Theorem~\ref{thm:stress}) & D2+D3 and Axioms~\ref{ax:screen}--\ref{ax:entropy} & Borchers--Wiesbrock half-sided modular inclusion; distributional differentiation on half-lines & the theorem-local inherited-strip decomposition and exact-or-controlled Markov hypotheses of Proposition~\ref{prop:nullec} and Corollary~\ref{cor:n1}, plus T2; quasi-local propagation and endpoint control are supplied internally by Axiom~\ref{ax:maxent} on the local finite-constraint branch; the imported continuum input is the standard local modular-Hamiltonian formula for the same half-line family, and bounded-interval formulas still require the separate interval-preserving projective branch & Phase I conditional bridge theorem \\
\textbf{D5} & Jacobson-type Einstein branch: rest-frame relation plus conditional tensor upgrade (Lemma~\ref{lem:smallball}, Theorem~\ref{thm:einstein}, Corollary~\ref{cor:einstein}) & D3+D4 and Axioms~\ref{ax:screen}--\ref{ax:entropy} & Jacobson small-ball first law; null-to-tensor reconstruction modulo a metric term & T3 fixed-cap stationarity; locally Lorentzian \(d=4\) scaling regime; the separate interval-preserving projective branch used to transport the D4 half-line generator/charge identification to bounded intervals; small-ball constancy assumptions; D4 carried remainder negligible at order \(\ell^4\); all local directions/reference states for the tensor upgrade & Phase I conditional scaling-limit theorem/corollary \\
\textbf{D6} & Capacity branch for \(\Lambda\), plus a legacy capacity-level neutrino side estimate (Proposition~\ref{prop:lambda}, Corollary~\ref{cor:lambda}) & D5 leaves a \(+\Lambda g_{ab}\) ambiguity & de Sitter entropy relation and dimensional analysis & external input \(N_{\mathrm{scr}}\) and the screen-capacity identification & Phase II input-dependent corollary / estimate \\
\textbf{D7} & Compact gauge reconstruction (Theorem~\ref{thm:compactgauge}) & Axioms~\ref{ax:screen}--\ref{ax:entropy} & Doplicher--Roberts / Tannaka reconstruction & Assumption~\ref{ass:gaugeprem} on the ordinary or central-defect bosonic branch; the genuinely noncentral fixed-cutoff branch is handled separately by Theorems~\ref{thm:hgdatum}--\ref{thm:hgtransport} and is not itself the ordinary compact-group reconstruction theorem & Phase I conditional structural theorem \\
\textbf{D8} & Product gauge structure up to finite quotient (Lemmas~\ref{lem:mincontent}--\ref{lem:commutant}, Theorem~\ref{thm:sm}) & D7 + Axiom~\ref{ax:mar} & compact Lie representation classification; Schur's lemma & the same ordinary or central-defect bosonic branch as D7, together with a connected positive-dimensional Lie admissible class, one connected abelian factor, and faithful action on the minimal coupled carrier & Phase I realized-branch theorem \\
\textbf{D9} & Realized Standard Model quotient, hypercharges, the counting chain \(N_g=3\) then \(N_c=3\), and product-group corollaries (Theorem~\ref{thm:hypercharge}, Corollaries~\ref{cor:ng}, \ref{cor:nc}, Proposition~\ref{prop:z6}) & D8 & anomaly-cancellation algebra; Witten global-anomaly argument; CKM CP counting & realized one-generation chiral matter plus one Higgs package; asymptotic-freedom and MAR admissibility premises & Phase I realized-branch theorem/corollary chain \\
\textbf{D10} & Forward gauge-coupling closure and target-free source-only electroweak readout of the integrated D10 calibration package & D9 + external input \(P\) & printed RG evolution and matching conventions & pixel constraint and the source-only forward transmutation solve \(\mathcal F(\alpha_U;P)=0\); the fixed-cutoff edge heat-kernel / Casimir theorem on the microphysics surface together with the compact-group / Peter--Weyl lift used on the D10 lane; printed beta-function and threshold conventions & Phase II calibration sector \\
\textbf{D12} & Charged-lepton centered-readback / centered-to-affine frontier, texture, dark-sector, baryogenesis, black-hole spectroscopy, proton-spin, proton-lifetime estimates beyond the gauge-channel exclusion, conditional large-\(N_{\mathrm{edge}}\) string/worldsheet reorganizations beyond the fixed-cutoff heat-kernel / Yang--Mills form on the declared compact-gauge branch, conjectural critical-superstring extensions, and other continuations & various subsets of D6, D9, and D10 & branch-specific EFT and phenomenological manipulations & additional ans\"atze such as the uniform \(\mathbb Z_6\) center-label ensemble, discrete texture choices, dark-sector response assumptions, discrete-horizon assumptions, a controlled large-\(N_{\mathrm{edge}}\) regime distinct from the proved compact-gauge chain, or extra worldsheet/CFT premises & Phase III phenomenological continuation branch \\
\end{longtable}
\endgroup

The recovered core of this SM/GR derivation paper is D1--D5 together with D7--D9. D6 is a separate input-dependent corollary, D10 is the integrated calibration sector, and D12 collects phenomenological continuations. The local/global handoff is exactly the D5\(\to\)D6 handoff: D5 recovers the local Einstein branch only modulo a metric term, and D6 fixes that remaining constant only after the global screen-capacity input \(N_{\mathrm{scr}}\) is supplied. The modular-anomaly dark-sector branch belongs to D12 and does not derive MOND/RAR-style response laws or galaxy/cluster phenomenology. The only external continuous inputs are \(N_{\mathrm{scr}}\) in D6 and \(P\) in D10, together with descendants built from those branches.
\section{Five Axioms, Two External Continuous Inputs in the Current Quantitative Implementation, and Technical Premises}\label{sec:premise-ledger}

This section states the premise ledger and dependency map for this SM/GR derivation paper. The OPH basis used below consists of the five core axioms, the two external continuous inputs listed next, the scaling/continuum premises collected afterward, and the gauge-side assumptions of Assumption~\ref{ass:gaugeprem}. On the fixed-cutoff local finite-constraint MaxEnt branch, the local-Gibbs form, quasi-local dynamics, Lieb--Robinson propagation control, and endpoint-Lipschitz interval control are internal consequences rather than theorem-external regularity selectors. Any Dobrushin/local-Gibbs/mixing language used below names only a sufficient fixed-cutoff collar-recoverability mechanism on that branch. It is not an infinite-volume uniqueness claim for the refinement-limit theory, and it does not by itself constrain whether a transportable refinement-directed edge-sector colimit exists or whether, if it exists, it is trivial or nontrivial. The paper therefore treats any later nontrivial transportable sector colimit as explicit gauge-branch data, not as an entropy-side consequence of the MaxEnt/mixing package. What Axiom~\ref{ax:maxent} fixes under refinement is the realized state-side MaxEnt branch itself. Theorem~\ref{thm:geosubnet} fixes the operational/geometric split at fixed cutoff, and Theorem~\ref{thm:bw} fixes the cap generator and its \(2\pi\) normalization once its theorem-local geometric cap-pair extraction and ordered cut-pair rigidity hypotheses are supplied on that geometric subnet. Corollary~\ref{cor:n2} derives the null half-sided modular pair only after the later null blow-up step; and the remaining canonical theorem-external items are therefore the scaling-limit scope clause, the bounded-interval transport/projective branch where invoked, fixed-cap stationarity, and the bosonic ordinary/central gauge premises recorded in Assumption~\ref{ass:gaugeprem}. The genuinely noncentral fixed-cutoff branch is tracked instead by Theorems~\ref{thm:hgdatum} and \ref{thm:hgtransport}; it is not itself the ordinary compact-group theorem used later.

The formulation uses two external continuous configuration inputs:
\begin{align}
P &\equiv a_{\mathrm{cell}}/\ell_P^2 = 1.63094,\\
N_{\mathrm{scr}} &\equiv \log \dim \mathcal H_{\mathrm{tot}} \sim 10^{122}.
\end{align}
In the quantitative implementation used here, \(P\) is treated as the external particle-physics scale input for the D10 calibration sector and \(N_{\mathrm{scr}}\) is inferred from the observed cosmological constant. The first feeds the forward D10 particle-physics map. The second enters the cosmological-capacity branch. Any retained capacity-level neutrino side estimate on that input is legacy bookkeeping rather than the live weighted-cycle theorem lane.

\begin{axiom}[Screen Net]\label{ax:screen}
Physical data is organized on a horizon screen \(S^2\) carrying a net of local algebras
\[
P \longmapsto \mathcal A(P)
\]
for connected patches \(P\subset S^2\), with isotony
\[
P\subset Q \implies \mathcal A(P)\subset \mathcal A(Q).
\]
\end{axiom}

\begin{axiom}[Overlap Consistency]\label{ax:overlap}
For overlapping patches \(P_1\cap P_2\neq\varnothing\), the local states induced on the shared algebra agree:
\[
\omega_{P_1}|_{\mathcal A(P_1\cap P_2)}
=
\omega_{P_2}|_{\mathcal A(P_1\cap P_2)}.
\]
\end{axiom}

\begin{axiom}[Local MaxEnt and Refinement Stability]\label{ax:maxent}
At the regulator scale \(\ell_{\mathrm{UV}}\), the realized branch is selected by maximizing entropy subject to a finite family of gauge-invariant local constraints
\[
\mathcal C_{\ell_{\mathrm{UV}}}=\{O_a(x)\}_{a=1}^{N_{\mathrm{con}}},
\]
where each \(O_a(x)\) is supported in a ball of radius \(O(\ell_{\mathrm{UV}})\) and the label set \(a\) is independent of the number of regulator cells. Under refinement, the same finite constraint family is preserved, so the realized states at different cutoffs belong to one common finite-dimensional MaxEnt family. The realized low-energy branch is the refinement-stable branch of that family; symmetry-allowed relevant operators are therefore held at zero on that branch only by symmetry or because they lie in the explicitly retained constraint family. This is a statement about branch persistence, not a universal entropy-ordering theorem for arbitrary phases away from that branch.
\end{axiom}

\paragraph{Derived local-Gibbs branch.}
On the finite regulator realization of Axiom~\ref{ax:screen}, the standard finite-dimensional Lagrange-multiplier argument applied to Axiom~\ref{ax:maxent} gives
\[
\omega_{\ell_{\mathrm{UV}}}(\lambda)
=
Z_{\ell_{\mathrm{UV}}}(\lambda)^{-1}\exp\!\bigl(-K_{\ell_{\mathrm{UV}}}(\lambda)\bigr),
\]
with
\[
K_{\ell_{\mathrm{UV}}}(\lambda)
=
\sum_x \sum_{a=1}^{N_{\mathrm{con}}}\lambda_a O_a(x)
+\sum_b \mu_b Q_b,
\]
where the \(Q_b\) are the finitely many optional global conserved-charge constraints. Thus the logarithm of the selected state is itself a quasi-local UV generator built from the same bounded-support densities that define the branch.

\paragraph{Derived quasi-local propagation and interval control.}
Fix the regulator graph metric coming from cell adjacency and let \(\tau_t^{\ell_{\mathrm{UV}}}\) denote the automorphism group generated by \(K_{\ell_{\mathrm{UV}}}(\lambda)\), or more generally by any branch generator lying in the norm-closed algebra generated by the same bounded-support local densities. Standard Lieb--Robinson estimates for finite-range interactions on the finite regulator net then give constants \(C,\xi,v_{\ell_{\mathrm{UV}}}<\infty\) such that for local observables \(A_X,B_Y\), \cite{liebRobinson72}
\[
\bigl\|[\tau_t^{\ell_{\mathrm{UV}}}(A_X),B_Y]\bigr\|
\le
C\,\|A_X\|\,\|B_Y\|\,\min(|X|,|Y|)\,
e^{-(d(X,Y)-v_{\ell_{\mathrm{UV}}}|t|)/\xi}.
\]
So finite-velocity quasi-local propagation is branch-internal: it is the local finite-constraint MaxEnt branch written in propagation form, not an extra selector.

The same locality statement yields the bounded-interval endpoint control used later in the null modular chain. Whenever later sections use a fixed-cutoff interval or collar generator on this branch, it is the corresponding restriction of the same local density list, with the central endpoint term separated off exactly as in the later edge-center bookkeeping. Changing an interval \(I\) to \(I'\) then alters only an \(O(|I'\Delta I|)\) collar of local terms. Hence on any fixed local-energy-bounded domain one has
\[
\left|\langle\psi,(K[I']-K[I])\phi\rangle\right|
\le
C_{\psi,\phi,I_{\max}}\,|I'\Delta I|,
\]
which is precisely the endpoint-Lipschitz / finite-variation control used later in the null modular chain. This paper does \emph{not} separately derive a second independent microscopic generator beyond this quasi-local branch generator; if a UV completion carries one, the only later requirement is that it lie in the same bounded-support algebraic closure so that the same propagation constants control the selected branch.

This local-Gibbs/Lieb--Robinson package, together with any Dobrushin-type estimate used later, should be read only as fixed-cutoff collar-local recoverability, support control, and carried-error bookkeeping on the selected branch. One may impose such estimates on a chosen finite collar model without deciding the refinement-limit gauge phase. The package is therefore not a proof that the refinement limit lands in a trivial thermodynamic phase, not a proof that transportable edge sectors persist under refinement, not a proof that the relevant sector colimit exists, and not a proof that such a colimit is nontrivial. Its theorem-level role is precisely the collar-local mixing and support/error control carried into the later fixed-cutoff and scaling-limit arguments.

\paragraph{Internal refinement notion.}
Choose any family of refinement channels \(\Phi_{\ell\to L}\) compatible with Axiom~\ref{ax:maxent}. Because the same finite local constraint family is preserved, the regulator-scale realized states are parameterized by one common finite-dimensional multiplier space \(\lambda\), and on the realized branch one may write
\[
\Phi_{\ell\to L}\bigl(\omega_{\ell}(\lambda)\bigr)
=
\omega_{L}\!\bigl(R_{\ell\to L}(\lambda)\bigr)
\]
for an induced map \(R_{\ell\to L}\) on that multiplier space. The refinement-stable realized branch is therefore the persistent trajectory or invariant subset selected inside this finite-dimensional space, not an imported theorem-external regularity package. This state-side refinement notion is enough to compare realized states across cutoffs. It does not, by itself, upgrade fixed-cutoff edge labels to transportable refinement-persistent sectors, and it does not show that the directed colimit of such sectors exists or is nontrivial. Any Dobrushin/local-mixing hypothesis used later is logically separate: it supplies a collar estimate on a fixed finite-dimensional model, not information about the refinement-limit gauge phase. Accordingly, whenever later sections speak of a ``refinement-stable directed colimit'' of edge sectors, the entire transportable directed-system/fiber clause is the explicit content of Assumption~\ref{ass:gaugeprem}; Axiom~\ref{ax:maxent} supplies only the realized state branch along which one asks whether such sector transport persists.

\begin{axiom}[Recoverable Generalized Entropy]\label{ax:entropy}
A generalized entropy functional exists on caps,
\[
S_{\mathrm{gen}}(C)=S_{\mathrm{bulk}}(C)+\langle L_C\rangle,
\]
where \(L_C\) is a positive edge-center entropy functional. In the semiclassical scaling branch, its leading coarse-grained contribution is identified with \(A(\partial C)/(4G)\). The functional obeys quantum focusing on null generators, and collar tripartitions have small conditional mutual information with controlled recovery maps \cite{liebRuskai73,petz86,petz88,fawziRenner15}.
\end{axiom}

\begin{axiom}[Minimal Admissible Realization]\label{ax:mar}
Among realized sector packages \(\mathfrak S\) consisting of the connected Lie gauge-sector image relevant in the low-energy EFT, its admissible light chiral matter content, and one Higgs doublet, and which are loop-coherent, anomaly-free, refinement-stable with light chiral matter, single-Higgs Yukawa-completable with one connected abelian charge factor acting nontrivially on the coupled carrier, intrinsically CP-capable, and weak-sector UV-completable, the realized package is the lexicographically minimal one under
\[
C(\mathfrak S)=\bigl(\chi_{\mathrm{cpl}},\,N_{\mathrm{nonab}},\,N_c,\,N_g\bigr).
\]
\(\mathfrak S\) is the sector package on which MAR acts; it is not the bare tensor category alone. MAR is therefore an explicit structural-economy axiom on admissible realized low-energy branches, not a theorem derived from the preceding axioms. \(\chi_{\mathrm{cpl}}\) is the coupled carrier dimension: the smallest unitary carrier containing a common irreducible block on which the admissible pseudoreal and complex nonabelian charge types both act nontrivially. This is stronger than the abstract minimal faithful representation dimension.
\end{axiom}

\begin{assumption}[Scaling and stationarity premises]\label{ass:technical}
This SM/GR derivation paper uses one canonical numbered list for the external scaling and stationarity premises. The fixed-cutoff realized presentation, the local-Gibbs form, quasi-local Lieb--Robinson propagation, endpoint-Lipschitz interval control, the induced finite-dimensional refinement branch, and the operational/geometric split of Theorem~\ref{thm:geosubnet} are not listed here because they are established internally from Axioms~\ref{ax:screen}--\ref{ax:maxent}. The remaining external items outside the gauge branch are:
\begin{itemize}[leftmargin=2.3em]
\item[\textbf{T2}] all Lorentz/null-modular/Einstein statements are claims about a controlled refinement/scaling limit of the regulator net, not literal exactness at fixed finite cutoff;
\item[\textbf{T3}] fixed-cap generalized-entropy stationarity for the admissible first-variation class used in the Jacobson branch;
\end{itemize}
The bosonic ordinary/central gauge premises are stated separately in Assumption~\ref{ass:gaugeprem}.
\end{assumption}

\begin{assumption}[Bosonic gauge reconstruction premises]\label{ass:gaugeprem}
Whenever compact gauge reconstruction is invoked on the ordinary or central-defect branch, the following additional premises are kept explicit.
\begin{enumerate}[label=(G\arabic*),leftmargin=2.3em]
\item loop coherence on that ordinary or central-defect branch, written \([z]=0\);
\item symmetric braiding in the relevant \(3+1\)-dimensional EFT branch;
\item a directed colimit \(\mathsf{Sect}_\infty=\varinjlim_r \mathsf{Sect}_r\) of transportable edge-sector categories with objectwise finite-dimensional fibers;
\item a faithful bosonic fiber functor on \(\mathsf{Sect}_\infty\).
\end{enumerate}
The genuinely noncentral fixed-cutoff branch is handled separately by Theorems~\ref{thm:hgdatum}--\ref{thm:hgtransport} and is not itself the ordinary compact-group reconstruction theorem.
\end{assumption}

\begin{remark}[Log convention]
Unless an explicit base is written, logarithms and entropies in this paper use natural logs (nats). Base-2 logarithms are written as \(\log_2\) and quoted in bits.
\end{remark}

\subsection*{Dependency DAG}

The table below is the dependency map for this SM/GR derivation paper. The ``Immediate OPH ingredients'' column records only parents internal to the OPH program. Imported standard mathematics and extra non-axiom physical premises or external inputs are stated separately rather than hidden inside the word ``derived.''

\begingroup
\scriptsize
\setlength{\tabcolsep}{1pt}
\begin{longtable}{@{}>{\raggedright\arraybackslash}p{0.06\textwidth}>{\raggedright\arraybackslash}p{0.20\textwidth}>{\raggedright\arraybackslash}p{0.14\textwidth}>{\raggedright\arraybackslash}p{0.13\textwidth}>{\raggedright\arraybackslash}p{0.21\textwidth}>{\raggedright\arraybackslash}p{0.12\textwidth}@{}}
\toprule
Node & Theorem labels and output & Immediate OPH ingredients & Standard mathematics used & Extra non-axiom physical premises / external inputs & Status \\
\midrule
\endhead
\bottomrule
\endfoot
\textbf{D1} & Theorem~\ref{thm:confluence}: unique normal form and schedule independence for the declared recovery-derived repair relation & overlap-consistency problem on a finite patch net together with the declared fixed-cutoff collar recovery step, its touched-overlap local-fit acceptance contract, support-local disjoint commutation, and restriction-compatible union-collar gluing on the physical quotient & local-to-global Lyapunov descent on the inconsistency potential; local-diamond completion from overlap-associative center-sector / higher-gauge union-collar gluing; Newman's lemma & repair completeness; stated Petz support/CPTP control where that branch is used & structural theorem \\
\textbf{D2} & Theorem~\ref{thm:markovcollar}, Propositions~\ref{prop:approxrecover}, \ref{prop:splice}, \ref{prop:gensplit}, Definition~\ref{def:newton}: collar Markov/entropy layer & Axioms~\ref{ax:screen}--\ref{ax:entropy} & HJPW Markov structure theorem; Petz and Fawzi--Renner recovery theory & exact Markov identities require exact Markovity or an explicitly stated idealized recoverability limit; the \(A/(4G)\) identification is the coarse-grained dictionary of Definition~\ref{def:newton} & structural lemma/\newline proposition layer \\
\textbf{D3} & Theorem~\ref{thm:bw}, Corollary~\ref{cor:lorentz}: geometric modular flow and Lorentz branch & Axioms~\ref{ax:screen}--\ref{ax:entropy} & Bisognano--Wichmann modular geometry;\newline conformal classification on \(S^2\) & T2 together with the theorem-local geometric cap-pair extraction and ordered cut-pair rigidity hypotheses of Theorem~\ref{thm:bw}; Theorem~\ref{thm:geosubnet} fixes the target subnet internally & conditional scaling-limit\newline theorem \\
\textbf{D4} & Propositions~\ref{prop:nullec}, \ref{prop:n3}; Corollaries~\ref{cor:n1}, \ref{cor:n2}; Theorem~\ref{lem:density}, Lemma~\ref{lem:translation}; Theorem~\ref{thm:stress}: null modular bridge to \(T_{kk}\) & D2+D3 and Axioms~\ref{ax:screen}--\ref{ax:entropy} & Borchers--Wiesbrock half-sided modular inclusion; distributional differentiation on half-lines & the theorem-local inherited-strip decomposition and exact-or-controlled Markov hypotheses of Proposition~\ref{prop:nullec} and Corollary~\ref{cor:n1}, plus T2; quasi-local propagation and endpoint control are supplied internally by Axiom~\ref{ax:maxent} on the local finite-constraint branch; the imported continuum input is the standard local modular-Hamiltonian formula for the same half-line family, and bounded-interval formulas still require the separate interval-preserving projective branch & conditional bridge theorem \\
\textbf{D5} & Lemma~\ref{lem:smallball}, Theorem~\ref{thm:einstein}, Corollary~\ref{cor:einstein}: Jacobson-type Einstein branch via a rest-frame relation plus conditional tensor upgrade & D3+D4 and Axioms~\ref{ax:screen}--\ref{ax:entropy} & Jacobson small-ball first law; Lemma~\ref{lem:nullreconstruct} for the null-to-tensor upgrade modulo a metric term & T3 fixed-cap stationarity; locally Lorentzian \(d=4\) scaling regime; small-ball constancy assumptions; D4 carried remainder negligible at order \(\ell^4\); all local directions/reference states for the tensor upgrade & conditional scaling-limit theorem/corollary \\
\textbf{D6} & Proposition~\ref{prop:lambda}, Corollary~\ref{cor:lambda}: capacity branch for \(\Lambda\), plus a legacy neutrino side estimate & D5, which leaves the \(+\Lambda g_{ab}\) ambiguity & de Sitter entropy relation and dimensional analysis & external input \(N_{\mathrm{scr}}\) and the screen-capacity identification & input-dependent corollary / estimate \\
\textbf{D7} & Theorem~\ref{thm:compactgauge}: compact gauge reconstruction & Axioms~\ref{ax:screen}--\ref{ax:entropy} & Doplicher--Roberts / Tannaka reconstruction & Assumption~\ref{ass:gaugeprem} on the ordinary or central-defect bosonic branch; the genuinely noncentral fixed-cutoff branch is handled separately by Theorems~\ref{thm:hgdatum}--\ref{thm:hgtransport} and is not itself the ordinary compact-group reconstruction theorem & conditional structural theorem \\
\textbf{D8} & Lemmas~\ref{lem:mincontent}--\ref{lem:commutant}, Theorem~\ref{thm:sm}: product gauge structure up to finite quotient & D7 + Axiom~\ref{ax:mar} & compact Lie representation classification; Schur's lemma & the same ordinary or central-defect bosonic branch as D7, together with a connected positive-dimensional Lie admissible class, one connected abelian factor, and faithful action on the minimal coupled carrier & realized-branch theorem \\
\textbf{D9} & Theorem~\ref{thm:hypercharge}, Corollaries~\ref{cor:ng}, \ref{cor:nc}, Proposition~\ref{prop:z6}: realized Standard Model quotient, hypercharges, the counting chain \(N_g=3\) then \(N_c=3\), and product-group corollaries & D8 & anomaly-cancellation algebra; Witten global-anomaly argument; CKM CP counting & realized one-generation chiral matter plus one Higgs package; asymptotic-freedom and MAR admissibility premises & realized-branch theorem/corollary chain \\
\textbf{D10} & forward gauge-coupling closure and target-free source-only electroweak readout of Section~7 & D9 + external input \(P\) & printed RG evolution and matching conventions & pixel constraint and the source-only forward transmutation solve \(\mathcal F(\alpha_U;P)=0\); the fixed-cutoff edge heat-kernel / Casimir theorem on the microphysics surface together with the compact-group / Peter--Weyl lift used on the D10 lane; printed beta-function and threshold conventions & calibration sector \\
\textbf{D12} & Sections~8--10 continuations:\newline charged leptons,\newline proton spin/lifetime,\newline and the string/worldsheet lane beyond the fixed-cutoff heat-kernel / Yang--Mills form on the declared compact-gauge branch: conditional large-\(N_{\mathrm{edge}}\) worldsheet-like reorganization and conjectural critical-superstring extension & various subsets of D6, D9, and D10 & branch-specific EFT\newline and phenomeno\-logical\newline manipulations & additional ans\"atze such as the uniform\newline \(\mathbb Z_6\) center-label ensemble, texture choices,\newline dark-sector response assumptions,\newline discrete-horizon assumptions, a controlled\newline large-\(N_{\mathrm{edge}}\) regime distinct from the proved compact-gauge chain, or extra\newline worldsheet/CFT premises for the conjectural critical-superstring lift & phenomeno\-logical\newline continuation \\
\end{longtable}
\endgroup

The dependency DAG records the OPH parents, imported mathematics, and extra premises for each node.

\begin{theorem}[Summary theorem for the recovered relativity-plus-Standard-Model core]\label{thm:master}
This theorem packages the recovered-core nodes D1--D5 and D7--D9 of the dependency DAG above. Node D6 is a separate input-dependent corollary, while D10 and D12 remain outside this theorem's scope. Assume Axioms~\ref{ax:screen}--\ref{ax:entropy}, Axiom~\ref{ax:mar}, Assumptions~\ref{ass:technical} and \ref{ass:gaugeprem}, and the hypotheses of Theorems~\ref{thm:confluence}, \ref{thm:bw}, \ref{thm:stress}, \ref{thm:einstein}, \ref{thm:compactgauge}, \ref{thm:sm}, and \ref{thm:hypercharge}, together with Corollaries~\ref{cor:ng}, \ref{cor:nc}, and Proposition~\ref{prop:z6}. Then, in the controlled refinement/scaling regime singled out by Assumption~\ref{ass:technical}:
\begin{enumerate}[label=(\roman*),leftmargin=2em]
\item overlap repair admits a unique schedule-independent normal form;
\item cap modular flow on the extracted prime geometric cap pair is geometric and yields the connected Lorentz group
\[
\Conf^+(S^2)\cong \mathrm{SO}^+(3,1);
\]
\item admissible fixed-cap stationarity together with the null modular bridge and the bounded-interval projective branch yield the Jacobson-type rest-frame relation
\[
\delta\!\left(G_{00}+\Lambda g_{00}\right)=8\pi G\,\delta\langle T_{00}\rangle
\]
at the cap center in the diamond rest frame; if that rest-frame relation holds for all local directions and reference states in the scaling regime, then the null-to-tensor upgrade gives
\[
G_{ab}+\Lambda g_{ab}=8\pi G\,\langle T_{ab}\rangle;
\]
\item under Assumption~\ref{ass:gaugeprem}, the edge-sector category reconstructs a compact gauge group, and the realized connected gauge structure has the form
\[
\frac{\SU(3)\times\SU(2)\times\U(1)}{\Gamma}
\]
for some finite central subgroup \(\Gamma\).
\end{enumerate}
After the hypercharge lattice, generation-count, color-count, and trivial-action quotient steps, the finite quotient is fixed to \(\Gamma=\mathbb Z_6\), so the realized gauge structure on the realized MAR-admissible branch is
\[
\frac{\SU(3)\times\SU(2)\times\U(1)}{\mathbb Z_6},
\qquad
N_g=3,
\qquad
N_c=3,
\]
with the exact Standard Model hypercharge lattice on the realized matter package. Here the realized-branch counting chain is sequential: the asymptotic-freedom/CP-capability window and minimal admissibility first yield \(N_g=3\), and then Witten's \(\SU(2)\) anomaly parity constraint together with the same realized-branch selector yield \(N_c=3\).
\end{theorem}

\begin{proof}
Items (i)--(iii) are collected from Theorem~\ref{thm:confluence}, Corollary~\ref{cor:lorentz}, Theorem~\ref{thm:einstein}, and Corollary~\ref{cor:einstein}. The compact gauge reconstruction is Theorem~\ref{thm:compactgauge}, applied only after the transportable sector colimit and its faithful bosonic fiber functor are supplied by Assumption~\ref{ass:gaugeprem}; the preceding MaxEnt/local-Gibbs/collar-mixing package is not used there to prove existence of that colimit or to decide whether it is trivial or nontrivial. The product gauge structure up to finite quotient is Theorem~\ref{thm:sm}; and the final identification of the exact Standard Model quotient is Theorem~\ref{thm:hypercharge} together with Corollaries~\ref{cor:ng}, \ref{cor:nc}, and Proposition~\ref{prop:z6}.
\end{proof}

Theorem~\ref{thm:master} therefore packages only the recovered relativity-plus-Standard-Model core: D1--D5 together with D7--D9. The D6 capacity relation is a separate input-dependent corollary; the calibration sector and phenomenological continuations remain outside its scope.

\subsection*{Usually Postulated Elsewhere, Recovered or Continued Here with Explicit Status}

\begingroup
\footnotesize
\begin{tabular}{@{}>{\raggedright\arraybackslash}p{0.28\textwidth}>{\raggedright\arraybackslash}p{0.26\textwidth}>{\raggedright\arraybackslash}p{0.36\textwidth}@{}}
\toprule
Structure & Common treatment & OPH treatment \\
\midrule
Lorentz kinematics & background symmetry or starting axiom & conditional scaling-limit branch from geometric modular flow on screen caps \\
Einstein dynamics & fundamental field equation & conditional first-variation relation with tensor upgrade from generalized entropy, null modular data, and entanglement equilibrium \\
Gauge group & model input & conditional compact group reconstruction from edge sectors; exact SM quotient selected on the realized MAR-admissible branch \\
Hypercharge lattice & matter-assignment input & solved from anomaly cancellation and Yukawa invariance on the realized one-generation chiral matter plus one-Higgs package \\
Color and generation count & empirical input & fixed on the realized MAR-admissible branch by anomaly, CP capability, asymptotic freedom, and minimality \\
Flavor hierarchy unit & model-dependent small parameter & phenomenological ansatz \(\varepsilon=1/6\) motivated by a uniform \(\mathbb Z_6\) center-label ensemble \\
Charged-lepton continuation & Koide-style phenomenological fit & exact centered readback plus an open centered-to-affine frontier above \(\widehat C_e^{\mathrm{cand}}\); any \(\delta=2/9\) phase relation is continuation-only \\
Cosmological constant & local vacuum-energy puzzle & global screen-capacity term after the input-dependent capacity identification \\
\bottomrule
\end{tabular}
\endgroup
\section{Observer-Overlap Consistency as a Fixed-Point Problem}

The discrete overlap problem is the mathematical skeleton of the framework.

\begin{definition}[Patch Net]\label{def:patchnet}
Let $G=(V,E)$ be a finite connected graph. Each vertex $i\in V$ carries a finite local state space $S_i$. For each edge $e=\{i,j\}$ let $I_e$ be an interface alphabet and let
\[
\pi_{i,e}:S_i\to I_e,
\qquad
\pi_{j,e}:S_j\to I_e
\]
be interface projections. The global state space is
\[
\Sigma=\prod_{i\in V} S_i,
\]
and the consistency set is
\[
C=\bigl\{s\in\Sigma:\pi_{i,e}(s_i)=\pi_{j,e}(s_j)\ \text{for all }e=\{i,j\}\in E\bigr\}.
\]
\end{definition}

This finiteness is a regulator-level assumption for the discrete patch-net normal-form theorem; it is not the continuum limit statement.

Define the inconsistency potential
\[
\Phi(s)=\sum_{e=\{i,j\}\in E}w_e\, d_e\!\bigl(\pi_{i,e}(s_i),\pi_{j,e}(s_j)\bigr),
\]
with $w_e>0$ and $d_e(a,b)=0$ iff $a=b$.

\begin{definition}[Recovery-derived local repair law]
A law $\lambda$ is a family of local repair maps
\[
T_i^\lambda:\Sigma\to\Sigma
\]
changing only patch $i$ or a bounded neighborhood of $i$. On the fixed-cutoff collar branch these
are not free rewrite primitives: each accepted local update is read from exact Markov splice or a
declared Petz/Fawzi--Renner recovery move on a collar chart around \(i\), then lifted back to the
finite patch presentation, and committed only when it strictly lowers the touched-overlap
contribution to \(\Phi\) on the overlaps it can modify. Disjoint accepted repairs therefore
commute directly by support locality, and when two accepted repairs touch a common overlap they are read inside one
finite union collar whose physical glued state is parenthesization-independent on the quotient and
whose local decoders are restriction-compatible on nested collars.
\end{definition}

\begin{theorem}[Asynchronous Fixed-Point Law]\label{thm:confluence}
Assume normal forms of the repair relation are exactly the globally consistent states $C$.
Then every $s\in\Sigma$ has a unique normal form
\[
\nf_\lambda(s)\in C,
\]
and every maximal asynchronous repair sequence terminates at that same state.
\end{theorem}

\begin{proof}
Each accepted local update changes only its touched-overlap set and is committed only when the sum
of the corresponding \(\Phi\)-terms strictly decreases. Therefore every accepted repair strictly
lowers the global inconsistency potential \(\Phi\). Because \(\Sigma\) is finite, no infinite
repair sequence exists, so the repair relation is terminating. If two accepted repairs are
disjoint, support locality makes them commute directly. If they overlap, the same definition places
them inside one parenthesization-independent union-collar gluing state on the physical quotient,
so the two one-step branches have a common descendant. Thus the repair relation is locally
confluent. Newman's lemma~\cite{newman42} now gives confluence. The resulting normal form is
unique. By assumption it is a globally consistent state.
\end{proof}

\begin{corollary}[Schedule Independence]
If observables are evaluated on $\nf_\lambda(s)$, the resulting physical law is independent of update order.
\end{corollary}

\begin{theorem}[Cycle Obstruction / Holonomy Criterion]\label{thm:holonomy}
Let $A$ be an abelian group. For each oriented edge $e:u\to v$ assign $b_e\in A$ and consider the affine overlap equations
\[
x_v-x_u=b_e.
\]
A global solution exists if and only if the signed sum of $b_e$ vanishes on every cycle.
\end{theorem}

\begin{proof}
Summing the edge equations around a cycle gives necessity. For sufficiency, fix a root and define each $x_v$ by summing labels along a path from the root to $v$. Vanishing cycle sums make this independent of the chosen path.
\end{proof}

\begin{definition}[Physical repair law and representative lift]\label{def:physicalrepair}
Suppose a local gauge group $\Gamma=\prod_i \Gamma_i$ acts on $\Sigma$ while leaving all interface
data invariant. Write
\[
q:\Sigma\to\Sigma/\Gamma,
\qquad
q(s)=[s].
\]
A physical repair law is the family of local quotient maps induced by the
recovery-derived collar updates just described:
\[
\overline T_i^\lambda:\Sigma/\Gamma\to\Sigma/\Gamma
\]
on the overlap-invariant quotient. A representative repair family is any family
\[
T_i^\lambda:\Sigma\to\Sigma
\]
with
\[
q\circ T_i^\lambda=\overline T_i^\lambda\circ q.
\]
\end{definition}

\begin{proposition}[Representative lifts descend to the quotient]\label{prop:repairquot}
For any representative repair family of Definition~\ref{def:physicalrepair},
\[
q\bigl(T_i^\lambda(\gamma\cdot s)\bigr)=q\bigl(T_i^\lambda(s)\bigr)
\qquad
\forall\, i,\ \forall\, \gamma\in\Gamma,\ \forall\, s\in\Sigma.
\]
In particular the repair relation descends to \(\Sigma/\Gamma\) without any extra
gauge-covariance axiom.
\end{proposition}

\begin{proof}
Because \(q(\gamma\cdot s)=q(s)\),
\[
q\bigl(T_i^\lambda(\gamma\cdot s)\bigr)
=
\overline T_i^\lambda\bigl(q(\gamma\cdot s)\bigr)
=
\overline T_i^\lambda\bigl(q(s)\bigr)
=
q\bigl(T_i^\lambda(s)\bigr).
\]
So gauge-equivalent inputs induce the same repaired physical state.
\end{proof}

\begin{proposition}[Gauge Quotient]\label{prop:gauge}
Under Definition~\ref{def:physicalrepair}, Proposition~\ref{prop:repairquot}, and
Theorem~\ref{thm:confluence}, the normal-form map descends to the quotient:
\[
\overline{\nf}_\lambda:\Sigma/\Gamma\to\Sigma/\Gamma,
\qquad
[s]\mapsto [\nf_\lambda(s)].
\]
Hence gauge-invariant observables are unique on the quotient, and the fixed-cutoff physical
observable algebra has representative-independent terminal expectations on that carrier.
\end{proposition}

\begin{proof}
By Proposition~\ref{prop:repairquot}, every repair step has quotient image determined only by the
current orbit. Iterating, the quotient image of any repair sequence depends only on the initial
orbit. Since every maximal repair sequence from \(s\) terminates at the unique normal form
\(\nf_\lambda(s)\) by Theorem~\ref{thm:confluence}, the terminal orbit \([\nf_\lambda(s)]\)
depends only on \([s]\). This makes \(\overline{\nf}_\lambda\) well-defined.
\end{proof}

On the fixed-cutoff quantum lift, the same quotient-local carrier does more than fix the terminal
orbit: for every declared fixed-cutoff physical algebra, the terminal expectation functional is
the same on any two representative lifts whose regional collar data lie in one quotient-local
glued state, even when the microscopic representatives differ by gauge or sector relabelings on
that carrier.

This removes the extra gauge-covariance axiom and keeps repair on its actual quotient-local
carrier: recovery dynamics on fixed-cutoff collars. On the declared fixed-cutoff branch, the
remaining theorem-local burden is repair completeness, together with the Petz support/CPTP clause
where that branch is used. The touched-overlap local-fit contract supplies Lyapunov
\(\Phi\)-descent on accepted moves, and the support-local commutation and union-collar
compatibility package belongs to the declared repair law itself. Stability of that package under
refinement or branch change is a separate question.

\begin{definition}[Ancilla stabilization]\label{def:ancillastab}
Fix a finite-cutoff OPH realization
\[
\mathfrak U=
\bigl(\{\mathcal H_P,\mathcal A(P),\omega_P\}_{P\in\mathcal P},\,\Gamma,\,T^\lambda\bigr).
\]
Choose finite-dimensional ancillary factors \(K_P\) with product state \(\eta=\bigotimes_{P\in\mathcal P}\eta_P\). The associated ancilla stabilization is the realization
\[
\mathfrak U^\eta=
\bigl(\{\mathcal H_P^\eta,\mathcal A^\eta(P),\omega_P^\eta\}_{P\in\mathcal P},\,\Gamma,\,(T^\lambda)^\eta\bigr),
\]
with
\[
\mathcal H_P^\eta:=\mathcal H_P\otimes K_P,\qquad
\mathcal A^\eta(P):=\mathcal A(P)\otimes \mathbf 1_{K_P},\qquad
\omega_P^\eta:=\omega_P\otimes \eta_P,
\]
and lifted repair dynamics acting trivially on the ancillas.
\end{definition}

\begin{proposition}[Ancilla-stable UV underdetermination]\label{prop:ancillauv}
Let \(\mathfrak U^\eta\) be an ancilla stabilization of a finite-cutoff OPH realization \(\mathfrak U\). Then:
\begin{enumerate}[leftmargin=2em]
\item observable expectations on the physical subalgebras are unchanged:
\[
\omega_P^\eta(a\otimes \mathbf 1_{K_P})=\omega_P(a)
\qquad
\forall\, a\in\mathcal A(P);
\]
\item the interacting local MaxEnt branch on the physical subalgebra is unchanged, since
\[
\omega_{\ell_{\mathrm{UV}}}^\eta(\lambda)=\omega_{\ell_{\mathrm{UV}}}(\lambda)\otimes \eta;
\]
\item for every collar split \(A:B:D\),
\[
I(A:D\mid B)_{\rho\otimes \eta}=I(A:D\mid B)_\rho,
\]
so the Fawzi--Renner remainder \(r_{\mathrm{FR}}(\varepsilon)\) and the collar Markov modulus \(\delta^{\mathrm M}_{A:B:D}(\varepsilon)\) are unchanged;
\item if the ancillas are inert under repair, then
\[
\nf_\lambda^\eta(s\otimes k)=\nf_\lambda(s)\otimes k,
\]
so the quotient normal form on physical observables is identical;
\item if the ancillas carry only trivial neutral sectors, the MAR-selected realized sector package is unchanged.
\end{enumerate}
Hence \(\mathfrak U\) and \(\mathfrak U^\eta\) are OPH-indistinguishable although they are different microscopic regulator realizations.
\end{proposition}

\begin{proof}
Item 1 is immediate from the definition of \(\mathcal A^\eta(P)\) and \(\omega_P^\eta\). Item 2 is the product-state form of the same fixed-cutoff MaxEnt branch. Item 3 follows from additivity of entropy for product ancillas, which cancels in the conditional mutual-information combination and therefore leaves both \(r_{\mathrm{FR}}\) and \(\delta^{\mathrm M}\) unchanged. Item 4 holds because the lifted repair maps do not act on the ancillary factors. Item 5 is immediate when the ancillas are neutral and carry no extra realized sector data.
\end{proof}

\begin{definition}[OPH-stable UV equivalence]\label{def:ophuveq}
Two finite-cutoff realizations \(\mathfrak U\) and \(\mathfrak U'\) are \emph{OPH-stably equivalent}, written
\[
\mathfrak U\sim_{\mathrm{OPH}}\mathfrak U',
\]
if after finite ancilla stabilizations they are related by local gauge-covariant \(^*\)-isomorphisms intertwining the observable patch net, overlap maps, local states, and repair dynamics on the physical subalgebras.
\end{definition}

\begin{corollary}[Unique physical UV branch only modulo OPH-stable equivalence]\label{cor:ophuveq}
Under the OPH axiom language, the UV invariant determined by the theory is the class \([\mathfrak U]_{\mathrm{OPH}}\), not a unique microscopic regulator presentation. Literal microscopic UV uniqueness is therefore not an OPH invariant.
\end{corollary}

\begin{proof}
Propositions~\ref{prop:repairquot} and \ref{prop:gauge} fix the schedule-independent physical
branch on the quotient, while Proposition~\ref{prop:ancillauv} shows that inert ancillary
refinements leave every OPH observable invariant. So the physical branch is fixed only modulo
\(\sim_{\mathrm{OPH}}\), not at the level of one microscopic representative.
\end{proof}

These three results encode much of the eventual physical interpretation. Objectivity becomes
confluence, gauge symmetry becomes quotient invariance induced by the physical overlap algebra, and
stable defects are represented by nontrivial overlap holonomy classes. The physically relevant
uniqueness statement is therefore quotient uniqueness together with ancilla-stable equivalence: OPH
fixes a unique gauge-invariant physical branch modulo boundary redundancy, implementation hiding,
and inert ancillary stabilization, not a unique microscopic regulator presentation.

\section{Collars, Edge Centers, and Generalized Entropy}

The fixed-point picture explains why overlap consistency produces gauge quotients: once repair is
read on the physical overlap algebra, descent to the quotient is automatic. The gravity and
consensus arguments later use a more specific collar fact: after edge-center completion, interior
observables become insensitive to compatible exterior substitutions, and modular additivity becomes
exact in the Markov normal form. The point of this section is therefore twofold. First, at fixed
regulator scale, we derive the collar block decomposition from overlap consistency itself on the
ordinary or central-defect branch and then state its genuinely noncentral higher-gauge replacement.
Second, we state precisely when the approximate recoverability premise of Axiom~\ref{ax:entropy} is
allowed to converge to the exact HJPW normal form used by the later spatial and null collar
theorems.

The sharp claim boundary is as follows. Small conditional mutual information always gives a constructive recovered comparison state with \(O(\varepsilon^{1/2})\) observable error. It does \emph{not} by itself give a universal one-shot trace-norm bound to an exact Markov state. The exact Markov normal form is recovered either when \(I(A:D\mid B)=0\) holds literally, or for a controlled family on one fixed finite-dimensional collar model, or after pullback to such a model, where \(\varepsilon\to0\). On that fixed collar one gets a collar-dependent modulus \(\delta^{\mathrm M}_{A:B:D}(\varepsilon)\to0\) measuring distance to the exact Markov set, and this is the error that must be carried into later exact splice or modular-additivity identities.

\begin{proposition}[Derived regulator gluing datum]\label{prop:regulatedrealization}
Fix a regulator-scale finite patch cover of a cap neighborhood and let \(R\) be any finite union of regulator cells. Assume only the finite patch-net presentation implicit in Definition~\ref{def:patchnet} together with overlap consistency on common interfaces. Then:
\begin{enumerate}[leftmargin=2em]
\item each regulator cell \(i\) with finite local state space \(S_i\) can be Hilbertized as \(\tilde{\mathcal H}_i\cong \mathbb C^{|S_i|}\), so the extended algebra before overlap quotienting is the finite type-I algebra
\[
\widetilde{\mathcal A}(R)=\mathcal B(\tilde{\mathcal H}_R),
\qquad
\tilde{\mathcal H}_R:=\bigotimes_{i\subset R}\tilde{\mathcal H}_i;
\]
\item for each connected boundary component \(\Sigma\subset\partial R\), after choosing a reference cut presentation \(\tilde{\mathcal H}_\Sigma\), every overlap-consistent change of local chart along \(\Sigma\) is implemented by a unitary on \(\tilde{\mathcal H}_\Sigma\), and the compact closure of the subgroup generated by all such recharting unitaries is a compact boundary redundancy group
\[
K_\Sigma\subset U(\tilde{\mathcal H}_\Sigma);
\]
before the choice of reference chart, the same data form a compact unitary groupoid of overlap-preserving transitions;
\item if triple-overlap defects are central, the projective composition law of item 2 lifts to an direct action of a compact central extension \(\widehat K_\Sigma\); a genuinely noncentral defect is the only obstruction to reducing the transition system to an ordinary compact group action;
\item in the invariant-state realization of the quotient, the chart-independent endomorphisms of the lifted presentation form the boundary-invariant algebra
\[
\mathcal A_{\mathrm{inv}}(R)=\widetilde{\mathcal A}(R)^{\widehat K_{\partial R}},
\qquad
\widehat K_{\partial R}:=\prod_{\Sigma\subset\partial R}\widehat K_\Sigma,
\]
\noindent with the convention \(\widehat K_\Sigma=K_\Sigma\) whenever no central extension is needed; the physical state space is the corresponding invariant subspace, and on collars the sector-preserving algebra induced on that subspace is the block-diagonal algebra used in Theorem~\ref{thm:markovcollar}.
\end{enumerate}
Items 1 and the lifted fixed-point presentation of item 4 are the fixed-cutoff realized data used later; they are not extra axioms beyond the finite regulator presentation plus overlap redundancy.
\end{proposition}

\begin{proof}[Proof sketch]
Item 1 is the Hilbertization of finite sets. For item 2, each overlap-consistent recharting acts on a finite-dimensional matrix algebra, and every \(^*\)-automorphism of such an algebra is inner. After transporting all local overlap presentations to a chosen reference chart, each recharting is therefore implemented by a unitary on the cut Hilbert space. The subgroup generated by those unitaries has compact closure inside a finite-dimensional unitary group, which yields \(K_\Sigma\); without fixing the reference chart one has the corresponding compact unitary groupoid. Item 3 is the usual projective-versus-central-extension lift for a central 2-cocycle. Item 4 records the lifted fixed-point presentation: chart-independent endomorphisms of the unreduced boundary data are exactly the fixed points of the derived boundary action, while the later collar theorem uses the invariant-state realization and the sector-preserving algebra induced on the matched-cut subspace. If the defect is genuinely noncentral, the quotient is still well defined, but the correct bookkeeping object is the crossed-module or higher-gauge data rather than an ordinary compact group.
\end{proof}

\begin{theorem}[Derived edge-center collar decomposition and exact Markov normal form]\label{thm:markovcollar}
Let \(A\)-\(B\)-\(D\) be a collar tripartition realized at fixed regulator scale on the ordinary or central-defect branch of Proposition~\ref{prop:regulatedrealization}. Write \(B=B_L\cup B_R\) with common interface \(\Sigma=\partial C\), and let \(\widehat K_\Sigma\) be the derived compact boundary action attached to that cut, with \(\widehat K_\Sigma=K_\Sigma\) on the ordinary branch. In the invariant-state realization,
\[
\mathcal H_B=(\tilde{\mathcal H}_{B_L}\otimes \tilde{\mathcal H}_{B_R})^{\widehat K_\Sigma}.
\]
Decompose the left boundary data as
\[
\tilde{\mathcal H}_{B_L}\cong \bigoplus_\alpha W_\alpha\otimes \mathcal H_{b_L^\alpha},
\]
with \(W_\alpha\) irreducible \(\widehat K_\Sigma\)-modules. Because the right half-collar carries the inverse overlap transport across the same cut, it decomposes contragrediently:
\[
\tilde{\mathcal H}_{B_R}\cong \bigoplus_\beta W_\beta^*\otimes \mathcal H_{b_R^\beta}.
\]
Then
\[
\mathcal H_B\cong\bigoplus_\alpha \mathcal H_{b_L^\alpha}\otimes \mathcal H_{b_R^\alpha},
\]
and the sector-preserving collar algebra
\[
\mathcal A_{\mathrm{EC}}(B):=\bigoplus_\alpha \mathcal B(\mathcal H_{b_L^\alpha})\otimes \mathcal B(\mathcal H_{b_R^\alpha})
\]
has center
\[
Z(\mathcal A_{\mathrm{EC}}(B))=\bigoplus_\alpha \mathbb C\,\mathbf 1_\alpha.
\]
If, in addition, the reference state on \(A\cup B\cup D\) is exact Markov,
\[
I(A:D|B)_\rho=0,
\]
then the state decomposes as
\[
\rho_{ABD}
=
\bigoplus_\alpha p_\alpha\,
\rho^{(\alpha)}_{A b_L^\alpha}\otimes \rho^{(\alpha)}_{b_R^\alpha D}.
\]
The same formula holds in any limit state that is exact Markov on this fixed collar model.
Thus the collar-center block decomposition is forced by overlap consistency plus the derived regulator package on the ordinary or central-defect branch; exact Markov factorization is an additional state condition, not a consequence of the block decomposition alone.
\end{theorem}

\begin{proof}
By Proposition~\ref{prop:regulatedrealization}, one may realize the physical collar states as the \(\widehat K_\Sigma\)-invariant subspace of \(\tilde{\mathcal H}_{B_L}\otimes \tilde{\mathcal H}_{B_R}\). Complete reducibility of finite-dimensional unitary representations gives the displayed decomposition of \(\tilde{\mathcal H}_{B_L}\), and the right half-collar carries the inverse transport law across the same cut, hence the contragredient decomposition of \(\tilde{\mathcal H}_{B_R}\). Therefore
\[
\tilde{\mathcal H}_{B_L}\otimes \tilde{\mathcal H}_{B_R}
\cong
\bigoplus_{\alpha,\beta}
(W_\alpha\otimes W_\beta^*)\otimes
(\mathcal H_{b_L^\alpha}\otimes \mathcal H_{b_R^\beta}).
\]
Taking \(\widehat K_\Sigma\)-invariants gives
\[
\mathcal H_B
\cong
\bigoplus_{\alpha,\beta}
(W_\alpha\otimes W_\beta^*)^{\widehat K_\Sigma}\otimes
(\mathcal H_{b_L^\alpha}\otimes \mathcal H_{b_R^\beta}).
\]
By Schur's lemma,
\[
(W_\alpha\otimes W_\beta^*)^{\widehat K_\Sigma}
\cong
\begin{cases}
\mathbb C, & \alpha=\beta,\\
0, & \alpha\neq \beta,
\end{cases}
\]
which yields the displayed block decomposition. The sector-preserving collar algebra therefore has the displayed block-diagonal form, so its center is generated by the block projectors. Within each block, observables from \(A\cup B\) act only on the left factor and observables from \(B\cup D\) act only on the right factor. If the state is exact Markov, the standard HJPW structure theorem gives the blockwise factorized normal form \cite{hjpw04}. The same formula is used in the idealized recoverability limit when exact identities are taken literally.
\end{proof}

\begin{theorem}[Derived higher-gauge cut datum]\label{thm:hgdatum}
Fix a connected cut \(\Sigma\) and a finite regulator chart \(\{P_i\}_{i\in I_\Sigma}\) meeting along \(\Sigma\), with finite nerve \(N_\Sigma\). On the genuinely noncentral branch, the weak gluing data are implemented by a compact crossed module
\[
\mathbb K_\Sigma=(H_\Sigma\xrightarrow{\partial_\Sigma}G_\Sigma,\triangleright)
\]
together with maps
\[
g_{ij}:P_{ij}\to G_\Sigma,
\qquad
h_{ijk}:P_{ijk}\to H_\Sigma,
\]
obeying
\[
g_{ij}g_{jk}=\partial_\Sigma(h_{ijk})\,g_{ik},
\qquad
h_{jkl}h_{ijl}=(g_{ij}\triangleright h_{ikl})\,h_{ijk}.
\]
The corresponding higher-gauge change system on the cut is the compact semidirect product
\[
\mathcal T_\Sigma:=C^1(N_\Sigma,H_\Sigma)\rtimes C^0(N_\Sigma,G_\Sigma),
\]
with multiplication
\[
(\eta,u)\cdot(\eta',u')
=
\bigl(\eta\,(u\triangleright \eta'),\,uu'\bigr),
\]
and standard coboundary action
\[
g_{ij}\mapsto u_i\,\partial_\Sigma(\eta_{ij})\,g_{ij}\,u_j^{-1},
\]
\[
h_{ijk}\mapsto
(u_i\triangleright h_{ijk})\,
\eta_{ij}\,
(g_{ij}\triangleright \eta_{jk})\,
\eta_{ik}^{-1}.
\]
The fixed-cutoff physical sector on the genuinely noncentral branch is therefore the orbit class
\[
q_\Sigma=[(g,h)]\in \check H^2(N_\Sigma,H_\Sigma\to G_\Sigma),
\]
and no external continuous input \(P\) or \(N_{\mathrm{scr}}\) enters this theorem package.
\end{theorem}

\begin{proof}[Proof sketch]
On pair overlaps, every regulator-scale recharting acts on a finite-dimensional matrix algebra and is therefore inner, so it is implemented by a unitary on the cut Hilbert space. On triple overlaps, genuinely noncentral failure of strict composition is recorded by unitary associators rather than scalars. These \(1\)- and \(2\)-morphisms form a compact unitary \(2\)-group of rechartings. Skeletal strictification of a compact \(2\)-group is equivalent to crossed-module data, which yields the displayed compact \(\mathbb K_\Sigma\). Because the nerve is finite, \(C^1(N_\Sigma,H_\Sigma)\) and \(C^0(N_\Sigma,G_\Sigma)\) are finite products of compact groups, hence \(\mathcal T_\Sigma\) is compact. The displayed formulas are the standard crossed-module coboundary action, and the orbit class is exactly the nonabelian \v{C}ech \(2\)-class of the defect data; see \cite{bullivant17} for an explicit higher-lattice-gauge realization of the same crossed-module structures.
\end{proof}

\begin{theorem}[Higher-gauge edge-center completion]\label{thm:hgedgecenter}
Let \(B=B_L\cup B_R\) be a fixed-cutoff collar around a connected cut \(\Sigma\) on the genuinely noncentral branch of Theorem~\ref{thm:hgdatum}. Then, as unitary \(\mathcal T_\Sigma\)-modules,
\[
\widetilde{\mathcal H}_{B_L}
\cong
\bigoplus_{\lambda\in\Lambda_\Sigma}
W_\lambda\otimes \mathcal H_{b_L^\lambda},
\qquad
\widetilde{\mathcal H}_{B_R}
\cong
\bigoplus_{\mu\in\Lambda_\Sigma}
W_\mu^*\otimes \mathcal H_{b_R^\mu},
\]
for a finite set \(\Lambda_\Sigma\) of irreducible unitary representations of \(\mathcal T_\Sigma\) that actually occur. Hence the physical higher-gauge collar space is
\[
\mathcal H_B^{2g}
:=
\bigl(\widetilde{\mathcal H}_{B_L}\otimes \widetilde{\mathcal H}_{B_R}\bigr)^{\mathcal T_\Sigma}
\cong
\bigoplus_{\lambda\in\Lambda_\Sigma}
\mathcal H_{b_L^\lambda}\otimes \mathcal H_{b_R^\lambda},
\]
and the collar algebra and its center are
\[
\mathcal A_{2g}(B)
\cong
\bigoplus_{\lambda\in\Lambda_\Sigma}
\mathcal B(\mathcal H_{b_L^\lambda})\otimes
\mathcal B(\mathcal H_{b_R^\lambda}),
\qquad
Z\bigl(\mathcal A_{2g}(B)\bigr)
=
\bigoplus_{\lambda\in\Lambda_\Sigma}\mathbb C\,\mathbf 1_\lambda.
\]
\end{theorem}

\begin{proof}[Proof sketch]
Because \(\mathcal T_\Sigma\) is compact and the half-collar spaces are finite-dimensional, both \(\widetilde{\mathcal H}_{B_L}\) and \(\widetilde{\mathcal H}_{B_R}\) split orthogonally into irreducible unitary \(\mathcal T_\Sigma\)-modules with finite multiplicity. The right side sees inverse transport across the same cut, hence the contragredient module. Therefore
\[
\widetilde{\mathcal H}_{B_L}\otimes \widetilde{\mathcal H}_{B_R}
\cong
\bigoplus_{\lambda,\mu}
(W_\lambda\otimes W_\mu^*)\otimes
(\mathcal H_{b_L^\lambda}\otimes \mathcal H_{b_R^\mu}).
\]
Taking \(\mathcal T_\Sigma\)-invariants leaves
\[
(W_\lambda\otimes W_\mu^*)^{\mathcal T_\Sigma}
\cong
\mathrm{Hom}_{\mathcal T_\Sigma}(W_\mu,W_\lambda)
\cong
\begin{cases}
\mathbb C, & \lambda=\mu,\\
0, & \lambda\neq\mu,
\end{cases}
\]
by Schur's lemma. This yields the matched block decomposition and the center formula.
\end{proof}

\begin{theorem}[Higher-gauge Markov collar and carried errors]\label{thm:hgmarkov}
Let \(\rho_{ABD}\) be a faithful state on a collar whose middle region \(B\) carries the higher-gauge block decomposition of Theorem~\ref{thm:hgedgecenter}. If
\[
I(A:D\mid B)_\rho=0,
\]
then
\[
\rho_{ABD}
=
\bigoplus_{\lambda\in\Lambda_\Sigma}
p_\lambda\,
\rho_{A b_L^\lambda}\otimes \rho_{b_R^\lambda D}.
\]
If instead \(I(A:D\mid B)_\rho\le \varepsilon\) on one fixed faithful collar model, then the same carried Fawzi--Renner and fixed-collar Markov errors as in the ordinary branch continue to hold:
\[
r_{\mathrm{FR}}(\varepsilon)=2\sqrt{1-e^{-\varepsilon}}\le 2\sqrt{\varepsilon},
\qquad
4\lambda_\ast^{-1}\,\delta^{\mathrm M}_{A:B:D}(\varepsilon).
\]
\end{theorem}

\begin{proof}[Proof sketch]
Theorem~\ref{thm:hgedgecenter} shows that the higher-gauge collar algebra is still a finite direct sum of type-I tensor blocks. HJPW depends only on that finite-dimensional algebraic structure, not on whether the block label arose from an ordinary compact group, a central extension, or a crossed-module gauge system. Likewise, the Fawzi--Renner remainder and the fixed-collar Markov modulus are finite-dimensional entropy and recovery statements and are blind to the origin of the block label.
\end{proof}

\begin{theorem}[Higher-gauge defect transportability]\label{thm:hgtransport}
With the genuinely noncentral cut data of Theorem~\ref{thm:hgdatum}, define
\[
q_\Sigma:=\bigl[(g,h)\bigr]\in \check H^2(N_\Sigma,H_\Sigma\to G_\Sigma).
\]
Then:
\begin{enumerate}[leftmargin=2em]
\item \(q_\Sigma\) is invariant under all local rechartings and higher-gauge changes \((\eta,u)\in \mathcal T_\Sigma\).
\item Two genuinely noncentral representatives describe the same fixed-cutoff physical cut sector if and only if they determine the same class \(q_\Sigma\).
\item The defect is removable if and only if \(q_\Sigma=0\).
\end{enumerate}
Thus the genuinely noncentral branch carries the same transportability logic as the ordinary or central branch, but in crossed-module \v{C}ech degree \(2\) rather than ordinary group cohomology.
\end{theorem}

\begin{proof}[Proof sketch]
Items 1--3 are the crossed-module version of ordinary \v{C}ech transportability: coboundaries
change representatives while preserving classes, the trivial class strictifies the cocycle, and a
nontrivial class cannot be gauged away. In any explicit crossed-module lattice realization, the
same sectors are equivalently the connected components of the \(2\)-flat configuration groupoid,
or homotopy classes of maps into the classifying space \(B\mathbb K_\Sigma\).
\end{proof}

\begin{proposition}[Higher-gauge interacting compatibility]\label{prop:hginteracting}
In a fixed-cutoff crossed-module lattice realization of Theorem~\ref{thm:hgdatum}, one may add local collar terms
\[
K_{\mathrm{collar}}^{2g}
=
\sum_{v\in\Sigma_0}\beta_v(1-A_v)
+
\sum_{e\in\Sigma_1}\beta_e(1-B_e)
+
\sum_{c\in\mathcal C_\Sigma}\beta_c(1-C_c),
\]
where \(A_v\), \(B_e\), and \(C_c\) are the local gauge and fake-flat projectors of the chosen higher-gauge realization. These terms are bounded-support and commute in the standard higher-gauge projector realization \cite{bullivant17}, so the full fixed-cutoff generator remains quasi-local and inside the Axiom~\ref{ax:maxent} class.
\end{proposition}

\begin{proof}[Proof sketch]
The additional terms are local projector constraints supported on finitely many collar cells, so they preserve bounded support and quasi-locality. In the standard higher-gauge projector realization they commute, so the same finite-constraint MaxEnt/Lieb--Robinson bookkeeping used on the ordinary branch continues to apply.
\end{proof}

\begin{remark}[Status of the regulator package and the precise Markov boundary]
Proposition~\ref{prop:regulatedrealization}, Theorem~\ref{thm:markovcollar}, and Theorems~\ref{thm:hgdatum}--\ref{thm:hgtransport} together with Proposition~\ref{prop:hginteracting} show that edge-center completion is a theorem of overlap consistency plus the finite regulator presentation on the ordinary, central-defect, and genuinely noncentral higher-gauge branches. No preferred quantum-link realization is assumed on the ordinary or central branch, and the genuinely noncentral branch is handled by the compact crossed-module replacement rather than by one more ordinary-group lemma. The fixed-cutoff topological UV package is closed on all three branches. The continuum modular/geometric lift on the geometric subnet under \textbf{T2} and, whenever refinement-limit compact gauge reconstruction is invoked, the bosonic ordinary/central gauge premises of Assumption~\ref{ass:gaugeprem}, form the current boundary. The distinct issue addressed next is \emph{state} control: exact splice and modular-additivity identities require either literal exact Markovity or a controlled family that converges to the exact Markov set on one fixed finite-dimensional collar.
\end{remark}

\begin{definition}[Exact Markov set and collar-distance modulus]\label{def:markovmodulus}
Fix the finite-dimensional collar Hilbert space of Theorem~\ref{thm:markovcollar}. Let
\[
\mathfrak M_{A:B:D}
:=
\left\{
\sigma_{ABD}:\ I(A:D\mid B)_\sigma=0
\right\}
\]
be the set of exact Markov states on that fixed collar. Equivalently, by HJPW and Theorem~\ref{thm:markovcollar}, \(\sigma\in\mathfrak M_{A:B:D}\) iff
\[
\sigma_{ABD}
=
\bigoplus_\alpha q_\alpha\,
\sigma^{(\alpha)}_{A b_L^\alpha}\otimes \sigma^{(\alpha)}_{b_R^\alpha D}
\]
for the fixed edge-center decomposition of \(B\).

For a general state \(\rho_{ABD}\), define its distance to the exact Markov set by
\[
d_{\mathrm M}(\rho):=\inf_{\sigma\in\mathfrak M_{A:B:D}}\|\rho-\sigma\|_1,
\]
and the fixed-collar exact-Markov modulus by
\[
\delta^{\mathrm M}_{A:B:D}(\varepsilon)
:=
\sup\left\{
d_{\mathrm M}(\rho):
I(A:D\mid B)_\rho\le \varepsilon
\right\}.
\]
\end{definition}

\begin{proposition}[Controlled exact Markov-collar limit]\label{prop:controlledmarkov}
On a fixed finite-dimensional collar, the exact Markov set \(\mathfrak M_{A:B:D}\) is compact and
\[
\delta^{\mathrm M}_{A:B:D}(\varepsilon)\longrightarrow 0
\qquad
\text{as }\varepsilon\downarrow 0.
\]
Consequently, if \(\rho^{(n)}_{ABD}\) is any sequence of states on that same collar with
\[
I(A:D\mid B)_{\rho^{(n)}}\le \varepsilon_n,
\qquad
\varepsilon_n\to0,
\]
then there exist exact Markov states \(\sigma^{(n)}\in\mathfrak M_{A:B:D}\) such that
\[
\|\rho^{(n)}-\sigma^{(n)}\|_1
\le
\delta^{\mathrm M}_{A:B:D}(\varepsilon_n)
\longrightarrow 0.
\]
Thus approximate recoverability converges to the exact HJPW normal form used later.
\end{proposition}

\begin{proof}
The full state space on a fixed finite-dimensional Hilbert space is compact in trace norm. Conditional mutual information is continuous there because von Neumann entropy is continuous in finite dimension. Hence \(\mathfrak M_{A:B:D}\) is closed, so it is compact as a closed subset of a compact space.

Suppose \(\delta^{\mathrm M}_{A:B:D}(\varepsilon)\not\to 0\). Then there exist \(\eta>0\), a sequence \(\varepsilon_n\downarrow 0\), and states \(\rho^{(n)}\) with \(I(A:D\mid B)_{\rho^{(n)}}\le \varepsilon_n\) but \(d_{\mathrm M}(\rho^{(n)})\ge \eta\) for all \(n\). By compactness, pass to a trace-norm convergent subsequence \(\rho^{(n_k)}\to \rho^\ast\). Continuity of conditional mutual information gives
\[
I(A:D\mid B)_{\rho^\ast}
=
\lim_{k\to\infty} I(A:D\mid B)_{\rho^{(n_k)}}
=
0,
\]
so \(\rho^\ast\in \mathfrak M_{A:B:D}\). But then
\[
d_{\mathrm M}(\rho^{(n_k)})
\le
\|\rho^{(n_k)}-\rho^\ast\|_1
\longrightarrow 0,
\]
contradicting \(d_{\mathrm M}(\rho^{(n_k)})\ge \eta\). Thus \(\delta^{\mathrm M}_{A:B:D}(\varepsilon)\to 0\). The final statement follows by choosing \(\sigma^{(n)}\in\mathfrak M_{A:B:D}\) within \(\delta^{\mathrm M}_{A:B:D}(\varepsilon_n)+1/n\) of \(\rho^{(n)}\).
\end{proof}

\begin{proposition}[Approximate collar recovery]\label{prop:approxrecover}
Let \(A\)-\(B\)-\(D\) be a collar tripartition satisfying
\[
I(A:D|B)_\rho\le \varepsilon,
\]
and let $\mathcal R_{B\to BD}$ be a recovery map such that
\[
F\!\left(\rho_{ABD},(\mathrm{id}_A\otimes \mathcal R_{B\to BD})(\rho_{AB})\right)\ge e^{-\varepsilon/2}.
\]
Define the recovered comparison state
\[
\rho^{\mathrm{rec}}_{ABD}:=(\mathrm{id}_A\otimes \mathcal R_{B\to BD})(\rho_{AB}),
\qquad
r_{\mathrm{FR}}(\varepsilon):=2\sqrt{1-e^{-\varepsilon}}\le 2\sqrt{\varepsilon}.
\]
Then
\[
\|\rho_{ABD}-\rho^{\mathrm{rec}}_{ABD}\|_1\le r_{\mathrm{FR}}(\varepsilon).
\]
Consequently, for every CPTP map \(\Lambda_{BD\to B'D'}\) and every bounded observable \(X\) supported on \(A\cup B'\),
\[
\left|\mathrm{Tr}\!\left[X(\mathrm{id}_A\otimes \Lambda)(\rho_{ABD}-\rho^{\mathrm{rec}}_{ABD})\right]\right|
\le
\|X\|_\infty\,\|(\mathrm{id}_A\otimes \Lambda)(\rho_{ABD}-\rho^{\mathrm{rec}}_{ABD})\|_1
\le
\|X\|_\infty\,r_{\mathrm{FR}}(\varepsilon).
\]
\end{proposition}

\begin{proof}
The fidelity lower bound and the Fuchs--van de Graaf inequality give
\[
\|\rho_{ABD}-\rho^{\mathrm{rec}}_{ABD}\|_1
\le
2\sqrt{1-F(\rho_{ABD},\rho^{\mathrm{rec}}_{ABD})^2}
\le
2\sqrt{1-e^{-\varepsilon}}
=
r_{\mathrm{FR}}(\varepsilon).
\]
The estimate \(r_{\mathrm{FR}}(\varepsilon)\le 2\sqrt{\varepsilon}\) follows from \(1-e^{-x}\le x\). Contractivity of trace norm under CPTP maps gives the transported bound, and the observable estimate is the duality between trace norm and operator norm \cite{fawziRenner15}. The bound is constructive and dimension-free. What it controls is closeness to a recovered comparison state; no claim is made that \(\rho^{\mathrm{rec}}_{ABD}\) is itself exact Markov.
\end{proof}

\begin{proposition}[Exact splice theorem and controlled collar approximation]\label{prop:splice}
Under Theorem~\ref{thm:markovcollar}, let \(\sigma_{ABD}\in \mathfrak M_{A:B:D}\) have exact Markov form
\[
\sigma_{ABD}
=
\bigoplus_\alpha p_\alpha\,
\sigma^{(\alpha)}_{A b_L^\alpha}\otimes \sigma^{(\alpha)}_{b_R^\alpha D}.
\]
Let \(\tau^{(\alpha)}_{b_R^\alpha D'}\) be any normalized family of states compatible with the same right-boundary sectors, and define
\[
\sigma'_{AB D'}
=
\bigoplus_\alpha p_\alpha\,
\sigma^{(\alpha)}_{A b_L^\alpha}\otimes \tau^{(\alpha)}_{b_R^\alpha D'}.
\]
Define the common left algebra
\[
\mathcal A_{A b_L}:=
\bigoplus_\alpha \mathcal B(\mathcal H_{A b_L^\alpha})\otimes \mathbf 1_{b_R^\alpha},
\]
canonically represented on both direct-sum Hilbert spaces. Then:
\begin{enumerate}[leftmargin=2em]
\item for every \(X\in\mathcal A_{A b_L}\),
\[
\mathrm{Tr}(X\sigma'_{AB D'})=\mathrm{Tr}(X\sigma_{ABD});
\]
\item if \(\rho_{ABD}\) is any state on the same fixed collar with \(I(A:D\mid B)_\rho\le \varepsilon\), and if \(\sigma_\varepsilon\in\mathfrak M_{A:B:D}\) is chosen so that
\[
\|\rho_{ABD}-\sigma_\varepsilon\|_1
\le
\delta^{\mathrm M}_{A:B:D}(\varepsilon),
\]
then the corresponding exact splice \(\sigma'_{\varepsilon}\) satisfies
\[
\left|
\mathrm{Tr}(X\rho_{ABD})-\mathrm{Tr}(X\sigma'_{\varepsilon})
\right|
\le
\|X\|_\infty\,\delta^{\mathrm M}_{A:B:D}(\varepsilon)
\]
for all \(X\in\mathcal A_{A b_L}\).
\end{enumerate}
Hence the exact splice identity used later is justified either at exact Markovity or along any controlled collar family for which \(\delta^{\mathrm M}_{A:B:D}(\varepsilon_\delta)\to 0\).
\end{proposition}

\begin{proof}
For item 1, blockwise factorization gives
\[
\mathrm{Tr}(X\sigma'_{AB D'})
=
\sum_\alpha p_\alpha
\mathrm{Tr}\!\left(X_\alpha\sigma^{(\alpha)}_{A b_L^\alpha}\right)
\mathrm{Tr}\!\left(\tau^{(\alpha)}_{b_R^\alpha D'}\right).
\]
Here \(X_\alpha\) denotes the \(\alpha\)-block of \(X\) in \(\mathcal A_{A b_L}\). Each \(\tau^{(\alpha)}\) is normalized, so the second factor is \(1\). The same computation for \(\sigma_{ABD}\) gives the same value because the original right factor is likewise normalized.

For item 2, apply item 1 to \(\sigma_\varepsilon\):
\[
\mathrm{Tr}(X\sigma'_{\varepsilon})=\mathrm{Tr}(X\sigma_{\varepsilon}).
\]
Therefore
\[
\left|
\mathrm{Tr}(X\rho_{ABD})-\mathrm{Tr}(X\sigma'_{\varepsilon})
\right|
=
\left|
\mathrm{Tr}\!\left[X(\rho_{ABD}-\sigma_{\varepsilon})\right]
\right|
\le
\|X\|_\infty\,\|\rho_{ABD}-\sigma_{\varepsilon}\|_1
\le
\|X\|_\infty\,\delta^{\mathrm M}_{A:B:D}(\varepsilon).
\]
\end{proof}

\begin{proposition}[Transport of modular-additivity errors from the exact Markov reference]\label{prop:modulartransport}
Let \(K_X(\omega):=-\log \omega_X\) for a faithful reduced state on region \(X\), and define the modular defect
\[
\Delta K(\omega):=
K_{ABD}(\omega)-K_{AB}(\omega)-K_{BD}(\omega)+K_B(\omega).
\]
Fix a collar and let \(\rho\) be a faithful state with \(I(A:D\mid B)_\rho\le \varepsilon\). Choose \(\sigma_\varepsilon\in \mathfrak M_{A:B:D}\) with
\[
\|\rho-\sigma_\varepsilon\|_1\le \delta^{\mathrm M}_{A:B:D}(\varepsilon).
\]
Assume that the four marginals \(\rho_X\) and \((\sigma_\varepsilon)_X\) for \(X\in\{ABD,AB,BD,B\}\) all satisfy
\[
\rho_X\ge \lambda_\ast \mathbf 1,
\qquad
(\sigma_\varepsilon)_X\ge \lambda_\ast \mathbf 1
\]
on their supports for some \(\lambda_\ast>0\). Then
\[
\|\Delta K(\rho)-\Delta K(\sigma_\varepsilon)\|_\infty
\le
4\lambda_\ast^{-1}\,\delta^{\mathrm M}_{A:B:D}(\varepsilon).
\]
Since \(\Delta K(\sigma_\varepsilon)\) is central or blockwise constant, every matrix element of the modular defect of \(\rho\) is within \(4\lambda_\ast^{-1}\delta^{\mathrm M}_{A:B:D}(\varepsilon)\) of the exact Markov value.
\end{proposition}

\begin{proof}
By monotonicity of trace distance under partial trace,
\[
\|\rho_X-(\sigma_\varepsilon)_X\|_1
\le
\|\rho-\sigma_\varepsilon\|_1
\le
\delta^{\mathrm M}_{A:B:D}(\varepsilon)
\]
for each \(X\in\{ABD,AB,BD,B\}\). On the interval \([\lambda_\ast,1]\), the function \(\log x\) has derivative bounded by \(\lambda_\ast^{-1}\). By the standard integral representation of the operator logarithm, this implies the operator-norm Lipschitz bound
\[
\|K_X(\rho)-K_X(\sigma_\varepsilon)\|_\infty
=
\|\log (\sigma_\varepsilon)_X-\log \rho_X\|_\infty
\le
\lambda_\ast^{-1}\,\|\rho_X-(\sigma_\varepsilon)_X\|_\infty
\le
\lambda_\ast^{-1}\,\delta^{\mathrm M}_{A:B:D}(\varepsilon).
\]
Summing the four region contributions gives
\[
\|\Delta K(\rho)-\Delta K(\sigma_\varepsilon)\|_\infty
\le
4\lambda_\ast^{-1}\,\delta^{\mathrm M}_{A:B:D}(\varepsilon).
\]
For \(\sigma_\varepsilon\in \mathfrak M_{A:B:D}\), the exact HJPW block factorization implies that \(\Delta K(\sigma_\varepsilon)\) is central or blockwise constant, so the same bound controls all matrix elements relative to the exact Markov modular-additivity value.
\end{proof}

\begin{remark}[What is and is not carried forward]
There are therefore two distinct quantitative controls, and this paper keeps them separate:
\begin{enumerate}[leftmargin=2em]
\item the constructive one-shot recoverability bound \(r_{\mathrm{FR}}(\varepsilon)=O(\varepsilon^{1/2})\), which compares the physical state to a recovered comparison state and is enough for bounded-observable statements at one regulator stage;
\item the fixed-collar exact-Markov modulus \(\delta^{\mathrm M}_{A:B:D}(\varepsilon)\to 0\), which compares the physical state to the exact Markov set and is the correct quantity for justifying exact splice or modular-additivity identities in a controlled limit, with an additional factor \(4\lambda_\ast^{-1}\) when logarithms are involved.
\end{enumerate}
Later spatial and null collar theorems therefore use exact identities only in two regimes: literal exact Markovity, or a controlled refinement or scaling family in which the relevant \(\delta^{\mathrm M}(\varepsilon_\delta)\) tends to zero after transport to one fixed finite-dimensional collar model. The manuscript does \emph{not} claim a universal dimension-free one-shot trace-norm estimate from small conditional mutual information directly to an exact Markov state.
\end{remark}

\begin{remark}[Floor transfer to the exact-Markov reference]
On one fixed finite-dimensional collar model, Proposition~\ref{prop:modulartransport} does not require a second independent faithfulness input on the exact-Markov comparison family. Suppose the transported physical marginals satisfy
\[
\rho_X\ge \bar\lambda_{m,\delta}\mathbf 1
\]
eventually for each \(X\in\{ABD,AB,BD,B\}\), and choose exact-Markov replacements \(\sigma_{n,m,\delta}\) with
\[
\|\rho_X-(\sigma_{n,m,\delta})_X\|_1
\le
\delta^{\mathrm M}_{m,\delta}(\varepsilon_{n,m,\delta})
\longrightarrow 0.
\]
Then \(\|\rho_X-(\sigma_{n,m,\delta})_X\|_\infty\le \|\rho_X-(\sigma_{n,m,\delta})_X\|_1\), so for all sufficiently large \(n\),
\[
(\sigma_{n,m,\delta})_X
\ge
\tfrac12\bar\lambda_{m,\delta}\mathbf 1.
\]
Thus the only nonlatent faithfulness-side input below the modular-defect term is the eventual lower spectral bound on the transported marginals themselves; the exact-Markov reference inherits a common floor once the exact-Markov modulus tends to zero.
\end{remark}

\begin{proposition}[Generalized-entropy split]\label{prop:gensplit}
If the reduced cap state takes the block form
\[
\rho_C=
\bigoplus_\alpha p_\alpha
\left(
\rho^{(\alpha)}_{\mathrm{bulk},C}\otimes
\frac{\mathbf 1^{(\alpha)}_{\mathrm{edge}}}{d_\alpha}
\right),
\]
then
\[
S(\rho_C)=S_{\mathrm{bulk}}(C)+\mathrm{Tr}(\rho_C L_C),
\]
where
\[
S_{\mathrm{bulk}}(C)=H(p_\alpha)+\sum_\alpha p_\alpha S(\rho^{(\alpha)}_{\mathrm{bulk},C}),
\qquad
L_C=\sum_\alpha (\log d_\alpha)P_\alpha.
\]
\end{proposition}

\begin{proof}
The entropy of a direct sum is
\[
S\!\left(\bigoplus_\alpha p_\alpha \sigma_\alpha\right)
=
H(p_\alpha)+\sum_\alpha p_\alpha S(\sigma_\alpha).
\]
Apply this identity with
\[
\sigma_\alpha=
\rho^{(\alpha)}_{\mathrm{bulk},C}\otimes
\frac{\mathbf 1^{(\alpha)}_{\mathrm{edge}}}{d_\alpha}.
\]
Since
\[
S(\sigma_\alpha)=S(\rho^{(\alpha)}_{\mathrm{bulk},C})+\log d_\alpha,
\]
the first statement follows.
\end{proof}

\begin{definition}[Effective Newton constant in the refinement dictionary]\label{def:newton}
In the refinement-scaling regime, if
\[
\mathrm{Tr}(\rho_C L_C)\approx N_\Sigma \bar\ell(t),
\qquad
A(\partial C)\approx N_\Sigma a_{\mathrm{cell}},
\]
then matching $\mathrm{Tr}(\rho_C L_C)$ to the area term $A(\partial C)/(4G)$ defines the effective Newton constant
\[
G:=\frac{a_{\mathrm{cell}}}{4\bar\ell(t)}.
\]
This is a continuum dictionary identification, not part of the exact algebraic proposition.
\end{definition}

Proposition~\ref{prop:splice} is the algebraic version of interior invariance under compatible exterior substitutions. Proposition~\ref{prop:modulartransport} is the precise bridge from small collar conditional mutual information to the exact modular identities later used by the spatial and null branches: at finite stage one carries a remainder, and exact identities appear only when that remainder is zero or tends to zero in the controlled collar limit. Proposition~\ref{prop:gensplit} isolates the exact entropy split, while Definition~\ref{def:newton} supplies the continuum identification that turns the edge center into a gravitational coupling.

\section{Modular Geometry, Relativity, and Einstein Dynamics}\label{sec:gravity}

\subsection{Geometric modular flow}

Before stating the continuum modular theorem, fix the algebraic target. At fixed cutoff, the full operational patch algebra contains geometric overlap data, declared record summaries, and compare/write/verify pointer layers. The BW question should be asked only on the overlap-generated geometric subnet, not on the entire operational algebra.

\begin{definition}[Operational total algebra, geometric subnet, and auxiliary/record sector]\label{def:geosubnet}
At fixed cutoff \(n\) and patch \(P\), let \(\mathcal A_n^{\mathrm{tot}}(P)\) denote the operational patch algebra consisting of the physical patch algebra together with the declared overlap-readable summaries used by the compare/write/verify surface. Let \(B_{IJ}\subset P\) range over overlap collars incident on \(P\), let \(\Pi_\alpha^{(IJ)}\) denote the overlap sector projectors, and let \(Q_a^{(IJ)}\) denote the boundary observables carrying the geometric cut data. Define the geometric subnet by
\[
\mathcal A_n^{\mathrm{geo}}(P)
:=
W^\ast\!\Big\langle
\bigcup_{B_{IJ}\subset P}
\iota_{IJ\to P}\!\bigl(
\{\Pi_\alpha^{(IJ)}\}_\alpha\cup\{Q_a^{(IJ)}\}_a
\bigr)
\Big\rangle_{\mathrm{repair,isotony}}
\Big/
\mathcal N_n(P),
\]
where \(\mathcal N_n(P)\) is the maximal repair-invariant overlap-trivial kernel, namely the maximal part whose restriction to every overlap algebra inside \(P\) is scalar. The complementary operational sector generated by record summaries, pointer registers, and other interface-inert observables is denoted \(\mathcal A_n^{\mathrm{aux/rec}}(P)\).
\end{definition}

\begin{theorem}[Geometric subnet and operational primeness]\label{thm:geosubnet}
For each fixed-cutoff patch \(P\), the operational algebra admits a canonical quotient
\[
\pi_{n,P}^{\mathrm{geo}}:\mathcal A_n^{\mathrm{tot}}(P)\twoheadrightarrow \mathcal A_n^{\mathrm{geo}}(P)
\]
with the following properties:
\begin{enumerate}[label=\textup{(\roman*)},leftmargin=2em]
\item \(\mathcal A_n^{\mathrm{geo}}(P)\) is generated by the overlap sector projectors and geometric boundary observables, and is closed under repair and isotony;
\item the record-summary block, pointer algebra, and interface-inert auxiliary observables lie in the complementary operational sector \(\mathcal A_n^{\mathrm{aux/rec}}(P)\) and do not enlarge \(\mathcal A_n^{\mathrm{geo}}(P)\);
\item if \(N\subset \mathcal A_n^{\mathrm{geo}}(P)\) is a repair-invariant subalgebra whose restriction to every overlap algebra inside \(P\) is scalar, then \(N=\mathbb C\,\mathbf 1\).
\end{enumerate}
\end{theorem}

\begin{proof}
Items \textup{(i)} and \textup{(ii)} are the content of Definition~\ref{def:geosubnet}. For item \textup{(iii)}, if such an \(N\) were nontrivial, then its pullback under \(\pi_{n,P}^{\mathrm{geo}}\) would define a nontrivial repair-invariant overlap-trivial component of \(\mathcal A_n^{\mathrm{tot}}(P)\), contradicting the maximality of \(\mathcal N_n(P)\). Thus no nontrivial overlap-trivial factor survives inside \(\mathcal A_n^{\mathrm{geo}}(P)\).
\end{proof}

The collar analysis proves a fixed-cutoff statement on the finite type-I regulator net: the reduced cap state has a literal density matrix, its modular Hamiltonian exists, and its nonadditive part is confined to a shrinking collar up to carried errors. The Lorentz claim is therefore not a literal statement about the finite regulator matrices. Under the controlled scaling clause \textbf{T2}, the target object is a realized scaling-limit geometric cap pair
\[
(\mathcal A_\infty^{\mathrm{geo}}(C),\omega_\infty^{\mathrm{geo},C})
\]
for each round cap \(C\subset S^2\). Axiom~\ref{ax:maxent} controls the state-side realized branch across refinement through one common finite-dimensional MaxEnt family, and Theorem~\ref{thm:geosubnet} removes the overlap-trivial auxiliary/record layers from the BW target. What still remains theorem-local on the present UV surface is not a separate global branch selector but two explicit objects on that geometric subnet: a realized geometric cap-pair extraction from transported marginals, and ordered cut-pair rigidity on the extracted prime cap pair. The theorem below makes that claim boundary explicit, fixes the \(2\pi\) normalization once those two objects are supplied, and leaves the remaining UV burden exposed rather than hidden.

Fix a cap \(C\subset S^2\) and a shrinking collar family \((A_\delta,B_\delta,D_\delta)\) around \(\partial C\), with \(\delta\downarrow 0\) and \(A_\delta\cup B_\delta\cup D_\delta\) covering a neighborhood of the cut. Write
\[
\varepsilon_\delta:=I(A_\delta:D_\delta\mid B_\delta)_\omega,
\qquad
r_{\mathrm{FR}}(\varepsilon_\delta):=
2\sqrt{1-e^{-\varepsilon_\delta}}
\le 2\sqrt{\varepsilon_\delta},
\]
and, on each fixed faithful collar model,
\[
\eta_\delta^{\mathrm M}
:=
4\lambda_\ast^{-1}\,
\delta^{\mathrm M}_{A_\delta:B_\delta:D_\delta}(\varepsilon_\delta),
\]
where \(\lambda_\ast>0\) is the lower spectral bound used to compare modular Hamiltonians. Every exact collar identity below is therefore to be read in one of two ways:
\begin{enumerate}[leftmargin=2em]
\item literal exactness when the reference collar state is exact Markov; or
\item a controlled collar family for which
\[
\delta^{\mathrm M}_{A_\delta:B_\delta:D_\delta}(\varepsilon_\delta)\to 0
\qquad(\delta\downarrow 0),
\]
with the finite-stage errors \(r_{\mathrm{FR}}(\varepsilon_\delta)\) and \(\eta_\delta^{\mathrm M}\) carried explicitly.
\end{enumerate}

For each cap \(C\), let \(\lambda_C(s)\subset \Conf^+(S^2)\) denote the standard cap-preserving conformal one-parameter subgroup, normalized so that the null blow-up near a smooth cut acts by \(v\mapsto e^{-s}v\), and let \(\alpha_{\lambda_C(s)}\) denote the induced automorphism of the scaling-limit cap net.

\begin{theorem}[Geometric modular flow on caps on the extracted prime geometric subnet]\label{thm:bw}
Assume Axioms~\ref{ax:screen}--\ref{ax:entropy}, Theorem~\ref{thm:geosubnet}, the derived fixed-cutoff regulator/collar package established above, and the controlled scaling-limit scope clause \textbf{T2} of Assumption~\ref{ass:technical}. Assume in addition that for each round cap \(C\subset S^2\):
\begin{enumerate}[label=\textup{(\alph*)},leftmargin=2em]
\item the transported cap marginals on \(\mathcal A_n^{\mathrm{geo}}(C_n)\) admit a realized scaling-limit geometric cap pair
\[
(\mathcal A_\infty^{\mathrm{geo}}(C),\omega_\infty^{\mathrm{geo},C})
\]
through the carried-collar local weak-\(*\) / GNS extraction package; and
\item ordered cut-pair rigidity holds on that extracted prime geometric pair, so the cap-preserving conformal subgroup singled out by the cut data is unique up to normalization.
\end{enumerate}
Then:
\begin{enumerate}[label=\textup{(\roman*)},leftmargin=2em]
\item At each finite regulator stage, the cap algebra is type I, the reduced cap state has a density matrix \(\rho_C^{(\delta)}\), and its modular Hamiltonian
\[
K_C^{(\delta)}:=-\log \rho_C^{(\delta)}
\]
has nonadditive part confined to the shrinking collar up to the carried errors measured by \(r_{\mathrm{FR}}(\varepsilon_\delta)\) and \(\eta_\delta^{\mathrm M}\). In particular, on each fixed collar model the physical collar state differs from its constructive recovered comparison state by \(r_{\mathrm{FR}}(\varepsilon_\delta)\) in trace norm, and its modular defect differs from the exact-Markov reference defect by at most \(\eta_\delta^{\mathrm M}\).

\item For every realized scaling-limit geometric cap pair \((\mathcal A_\infty^{\mathrm{geo}}(C),\omega_\infty^{\mathrm{geo},C})\) furnished by hypothesis \textup{(a)}, the modular automorphism group is geometric:
\[
\sigma_t^{\omega_\infty^{\mathrm{geo},C}}
=
\alpha_{\lambda_C(2\pi t)}.
\]
Equivalently, the modular parameter \(t\) and the geometric cap-dilation parameter \(s\) are related by
\[
s=2\pi t.
\]

\item No separate cap-isotropy/\(\mathrm{SO}(2)\)-equivariance selector or Euclidean-regularity premise is used in this theorem. Once ordered cut-pair rigidity has fixed the geometric cap subgroup on the extracted prime geometric pair and \(\lambda_C(s)\) is normalized by the null blow-up \(v\mapsto e^{-s}v\), the modular KMS normalization fixes the coefficient \(2\pi\).

\item If the scaling-limit cap algebra happens to remain type I, item \textup{(ii)} may be written as the operator identity
\[
K_C=2\pi B_C.
\]
In the generic continuum case, where the scaling-limit cap algebra has left the regulator class and is expected in QFT examples to be non-type-I, the correct theorem statement is item \textup{(ii)} itself, and the geometric modular action is generally outer rather than inner.
\end{enumerate}

Consequently, for every controlled collar family with
\[
\delta^{\mathrm M}_{A_\delta:B_\delta:D_\delta}(\varepsilon_\delta)\to 0
\]
and surviving faithfulness bound, replacing the finite-stage modular action by the geometric cap-dilation action incurs only the carried errors \(r_{\mathrm{FR}}(\varepsilon_\delta)\) and \(\eta_\delta^{\mathrm M}\), and these vanish in the refinement limit.
\end{theorem}

\begin{proof}
Step 1: fixed-cutoff collar control. On the finite type-I regulator net, Propositions~\ref{prop:splice}, \ref{prop:modulartransport}, and \ref{prop:gensplit} localize the modular defect to the shrinking collar and quantify the discrepancy between the physical collar state, its constructive recovered comparison state, and its exact-Markov reference by the carried errors \(r_{\mathrm{FR}}(\varepsilon_\delta)\) and \(\eta_\delta^{\mathrm M}\).

Step 2: scaling-limit claim tier. Theorem~\ref{thm:bw} is a scaling-limit statement by \textbf{T2}, so the theorem concerns only scaling-limit cap pairs furnished by the controlled refinement/scaling package, not the finite regulator matrices themselves. Nothing in this step assumes that type-I structure survives the limit.

Step 3: geometric identification on the prime subnet. By hypothesis \textup{(a)} and Theorem~\ref{thm:geosubnet}, the realized limit pair is extracted on the overlap-generated prime geometric subnet, with overlap-trivial auxiliary/record factors removed. By hypothesis \textup{(b)}, the cap-preserving conformal action singled out by the ordered cut data is unique up to normalization. Therefore there is a constant \(\kappa_C>0\) such that
\[
\sigma_t^{\omega_\infty^{\mathrm{geo},C}}=\alpha_{\lambda_C(\kappa_C t)}.
\]

Step 4: normalization. Because ordered cut-pair rigidity has identified the standard cap-preserving conformal subgroup on the extracted geometric pair, no separate cap-isotropy or Euclidean-regularity selector enters. The modular group is KMS at inverse temperature \(1\) by definition. Since \(\lambda_C(s)\) is normalized by unit null dilation under the blow-up \(v\mapsto e^{-s}v\), the Bisognano--Wichmann normalization fixes
\[
\kappa_C=2\pi,
\]
hence
\[
\sigma_t^{\omega_\infty^{\mathrm{geo},C}}
=
\alpha_{\lambda_C(2\pi t)}.
\]

Step 5: operator versus automorphism form. If the limit cap algebra remains type I, one can represent the modular automorphism group by a modular Hamiltonian \(K_C=-\log\rho_C\) and recover \(K_C=2\pi B_C\). If the limit algebra leaves the regulator class, the modular automorphism group remains well defined while the inner operator representative need not exist inside \(\mathcal A_\infty^{\mathrm{geo}}(C)\). In that case the automorphism identity is the full theorem statement. The final sentence is exactly the carried-error conclusion from Step 1 transported along the controlled scaling family of \textbf{T2}.
\end{proof}

\begin{remark}[BW-side UV scaffold after the geometric-subnet split]
Theorem~\ref{thm:geosubnet} removes the old ambiguity about the target algebra: the continuum Lorentz claim is asked only on the extracted prime geometric subnet, not on the full operational algebra. The UV-side boundary is therefore a two-step chain on \(\mathcal A_n^{\mathrm{geo}}\), and neither step is an independent cap-isotropy selector or a Euclidean-regularity premise.

First, one needs a \emph{realized transported geometric cap-local system} for caps and their conformal images. The relevant package is the cap-local test family on \(\mathcal A_n^{\mathrm{geo}}(C_n)\), the projectively compatible transported marginal family, and the asymptotic transport-equivalence certificate. On each fixed local collar model \((m,\delta)\), let \(\varepsilon_{n,m,\delta}\) be the transported collar conditional mutual information. If the transported physical geometric marginals satisfy an eventual common floor
\[
\rho_X^{(n,m,\delta)}\ge \bar\lambda_{m,\delta}\mathbf 1
\qquad
\bigl(X\in\{ABD,AB,BD,B\}\bigr)
\]
and
\[
\delta^{\mathrm M}_{A_{n,m,\delta}:B_{n,m,\delta}:D_{n,m,\delta}}
\bigl(\varepsilon_{n,m,\delta}\bigr)\to 0,
\]
then the exact-Markov comparison marginals inherit the same floor for all sufficiently large \(n\). The carried-collar schedule
\[
\eta_{n,m,\delta}^{\mathrm{cap}}
:=
r_{\mathrm{FR}}(\varepsilon_{n,m,\delta})
+
4\bar\lambda_{m,\delta}^{-1}
\delta^{\mathrm M}_{A_{n,m,\delta}:B_{n,m,\delta}:D_{n,m,\delta}}
\bigl(\varepsilon_{n,m,\delta}\bigr)
\]
is therefore \emph{derived rather than primitive}. The only nonlatent lower input on the faithfulness side is the eventual fixed-local-collar modular-transport common floor on the transported physical marginals themselves; no second comparison-state faithfulness clause is missing.

Second, once local weak-\(*\) extraction and GNS gluing have emitted the realized scaling-limit geometric cap pair, one needs \emph{ordered cut-pair rigidity} on that realized pair to collapse the residual cap-preserving conformal freedom to the standard hyperbolic subgroup. Only after that second object is supplied does the geometric modular statement become fully internal on \(\mathcal A_\infty^{\mathrm{geo}}(C)\). The present paper proves neither object and therefore does not claim that MaxEnt alone completes the BW lift.
\end{remark}

\begin{corollary}[Lorentz Kinematics]\label{cor:lorentz}
Under the hypotheses of Theorem~\ref{thm:bw},
\[
\Conf^+(S^2)\cong \mathrm{PSL}(2,\mathbb C)\cong \mathrm{SO}^+(3,1).
\]
The cap modular flows are therefore the standard one-parameter Lorentz boost/dilation subgroups in the celestial-sphere realization, and the induced local kinematic group is the connected Lorentz group.
\end{corollary}

\begin{proof}
Orientation-preserving conformal maps of \(S^2\) are exactly the M\"obius transformations, so
\[
\Conf^+(S^2)\cong \mathrm{PSL}(2,\mathbb C)\cong \mathrm{SO}^+(3,1).
\]
By Theorem~\ref{thm:bw}, on the extracted prime geometric subnet the realized scaling-limit cap modular automorphism groups are the standard cap-preserving conformal dilations with the \(2\pi\) normalization fixed internally. Hence the local kinematics induced by modular flow is the connected Lorentz group.
\end{proof}
\subsection{The null modular bridge}

The fixed-cutoff strip analysis begins from the transferred cut-center data and the inherited-strip Markov structure. At fixed regulator scale one obtains exact or controlled four-term strip additivity together with endpoint-Lipschitz control of the renormalized half-line family; on the same OPH geometric scaling branch used in Theorem~\ref{thm:bw}, the null half-line blow-up net then inherits geometric dilation and therefore half-sided modular inclusion. Standard Borchers--Wiesbrock theory~\cite{wiesbrock93} consequently applies to the resulting OPH half-sided modular pair, yielding positivity and unitary implementation of null translations. What remains downstream is narrower: the separate interval-preserving projective branch is still needed to transport the half-line result to bounded null intervals, and the later tensor upgrade still carries the standard null-invisible metric ambiguity. The half-line generator itself is identified below with the effective local null-stress charge.
\par
The present subsection therefore establishes the following chain:
\begin{enumerate}[leftmargin=2em]
\item null-strip center transfer and the inherited left/right split hypothesis;
\item exact strip additivity on the exact inherited Markov model, and a carried defect operator on controlled exact-Markov replacements on one fixed strip model;
\item endpoint-Lipschitz matrix elements for the renormalized half-line family and the resulting weak tail generator;
\item derived half-sided modular inclusion on the null half-line blow-up net from OPH scaling-limit geometry, and the resulting Borchers positive translation generator.
\end{enumerate}
The only explicit downstream boundary retained below is the bounded-interval transport/projective branch together with the later tensor upgrade; the half-line generator/charge identification is proved inside the bridge itself. Lemma~\ref{lem:translation} states the positive null-translation generator explicitly: the Borchers unitary group is constructed on its Stone domain, positivity is theorem-level rather than implicit, and the half-line modular Hamiltonians are tied to it by explicit affine covariance and quadratic-form identities. The density-upgrade template of Theorem~\ref{lem:density} is recorded separately and does not carry the operator-identification burden.

\begin{proposition}[Null-strip center transfer and inherited split]\label{prop:nullec}
Fix a null generator \(\Omega\) and a regulated tripartition
\[
I_-=(v_1,v_2),\qquad
J=(v_2,v_3),\qquad
I_+=(v_3,v_4),
\]
with cuts
\[
\Gamma_-:=\{v=v_2\},
\qquad
\Gamma_+:=\{v=v_3\}.
\]
Assume the derived fixed-cutoff regulator/collar package established above, and assume that the two null cuts inherit the ordinary or central-defect boundary-redundancy data of Proposition~\ref{prop:regulatedrealization} in the following precise sense:
\begin{enumerate}[leftmargin=2em,label=(\roman*)]
\item each cut \(\Gamma_\pm\) carries a compact derived boundary action \(\widehat K_\pm\) on a reference cut Hilbert space, with \(\widehat K_\pm=K_\pm\) on the ordinary branch;
\item in a compatible type-I regulator presentation, the adjacent regions carry inverse transport across each cut, so for irreducible \(\widehat K_\pm\)-modules \(W_{\alpha_\pm}\) one has decompositions
\[
\tilde{\mathcal H}_{I_-}
\cong
\bigoplus_{\alpha_-}
W_{\alpha_-}\otimes \mathcal H_{i_-^{\alpha_-}},
\]
\[
\tilde{\mathcal H}_{J}
\cong
\bigoplus_{\alpha_-,\alpha_+}
W_{\alpha_-}^{*}\otimes
\mathcal H_{j^{\alpha_-,\alpha_+}}\otimes
W_{\alpha_+},
\]
\[
\tilde{\mathcal H}_{I_+}
\cong
\bigoplus_{\alpha_+}
W_{\alpha_+}^{*}\otimes \mathcal H_{i_+^{\alpha_+}};
\]
\item the gauge-invariant strip algebra is the commutant of the transported boundary action on \(\tilde{\mathcal H}_J\), so that the pair \((\alpha_-,\alpha_+)\) is a central cut label;
\item for each \((\alpha_-,\alpha_+)\), the multiplicity space \(\mathcal H_{j^{\alpha_-,\alpha_+}}\) admits a chosen factorization
\[
\mathcal H_{j^{\alpha_-,\alpha_+}}
\cong
\mathcal H_{j_L^{\alpha_-,\alpha_+}}\otimes
\mathcal H_{j_R^{\alpha_-,\alpha_+}}
\]
such that, blockwise, \(\mathcal A(I_-\cup J)\) acts only on
\(\mathcal H_{i_-^{\alpha_-}}\otimes \mathcal H_{j_L^{\alpha_-,\alpha_+}}\)
and \(\mathcal A(J\cup I_+)\) acts only on
\(\mathcal H_{j_R^{\alpha_-,\alpha_+}}\otimes \mathcal H_{i_+^{\alpha_+}}\).
\end{enumerate}
Then the strip algebra has the inherited central decomposition
\[
\mathcal A(J)
\cong
\bigoplus_{\alpha_-,\alpha_+}
\mathcal B(\mathcal H_{j^{\alpha_-,\alpha_+}}),
\qquad
Z(\mathcal A(J))
=
\bigoplus_{\alpha_-,\alpha_+}\mathbb C\,P_{\alpha_-,\alpha_+},
\]
and, under item (iv),
\[
\mathcal A(J)
\cong
\bigoplus_{\alpha_-,\alpha_+}
\mathcal B(\mathcal H_{j_L^{\alpha_-,\alpha_+}})
\otimes
\mathcal B(\mathcal H_{j_R^{\alpha_-,\alpha_+}}).
\]
Moreover the glued tripartition carries the block decomposition
\[
\mathcal H_{I_-\cup J\cup I_+}
\cong
\bigoplus_{\alpha_-,\alpha_+}
\mathcal H_{i_-^{\alpha_-}}
\otimes
\mathcal H_{j_L^{\alpha_-,\alpha_+}}
\otimes
\mathcal H_{j_R^{\alpha_-,\alpha_+}}
\otimes
\mathcal H_{i_+^{\alpha_+}}.
\]
Thus the two null cuts transfer the same center data as the spatial collar branch, while the left/right split of the strip multiplicity spaces is exactly the extra decomposition-inheritance hypothesis and is not forced by center transfer alone.
\end{proposition}

\begin{proof}
By complete reducibility of the finite-dimensional unitary actions \(\widehat K_\pm\), the displayed decompositions of \(\tilde{\mathcal H}_{I_-}\), \(\tilde{\mathcal H}_J\), and \(\tilde{\mathcal H}_{I_+}\) exist. Under item (iii), the strip algebra is the commutant of the transported \(\widehat K_-\times \widehat K_+\) action on
\[
\tilde{\mathcal H}_J
=
\bigoplus_{\alpha_-,\alpha_+}
W_{\alpha_-}^{*}\otimes
\mathcal H_{j^{\alpha_-,\alpha_+}}\otimes
W_{\alpha_+}.
\]
Write an operator \(X\in \mathcal B(\tilde{\mathcal H}_J)\) in matrix form relative to that direct sum. The block
\[
X_{(\alpha_-,\alpha_+),(\beta_-,\beta_+)}
\]
intertwines
\[
W_{\alpha_-}^{*}\otimes W_{\alpha_+}
\longrightarrow
W_{\beta_-}^{*}\otimes W_{\beta_+}
\]
for the product action of \(\widehat K_-\times \widehat K_+\). By Schur's lemma, such an intertwiner is zero unless \((\alpha_-,\alpha_+)=(\beta_-,\beta_+)\), and when the sector pair agrees it is scalar on the representation factors. Therefore the commutant is exactly
\[
\mathcal A(J)
\cong
\bigoplus_{\alpha_-,\alpha_+}
\mathcal B(\mathcal H_{j^{\alpha_-,\alpha_+}}),
\]
and the center is generated by the corresponding block projectors \(P_{\alpha_-,\alpha_+}\).

For the glued tripartition, tensor the three regulator Hilbert spaces and decompose:
\[
\tilde{\mathcal H}_{I_-}\otimes \tilde{\mathcal H}_{J}\otimes \tilde{\mathcal H}_{I_+}
\cong
\bigoplus_{\alpha_-,\beta_-,\alpha_+,\beta_+}
(W_{\alpha_-}\otimes W_{\beta_-}^{*})
\otimes
\mathcal H_{i_-^{\alpha_-}}
\otimes
\mathcal H_{j^{\beta_-,\beta_+}}
\otimes
(W_{\beta_+}\otimes W_{\alpha_+}^{*})
\otimes
\mathcal H_{i_+^{\alpha_+}}.
\]
Taking \(\widehat K_-\times\widehat K_+\)-invariants and applying Schur's lemma to each cut gives
\[
(W_{\alpha_-}\otimes W_{\beta_-}^{*})^{\widehat K_-}
\cong
\begin{cases}
\mathbb C, & \alpha_-=\beta_-,\\
0, & \alpha_-\neq \beta_-,
\end{cases}
\qquad
(W_{\beta_+}\otimes W_{\alpha_+}^{*})^{\widehat K_+}
\cong
\begin{cases}
\mathbb C, & \beta_+=\alpha_+,\\
0, & \beta_+\neq \alpha_+,
\end{cases}
\]
so only matching cut sectors survive and one obtains
\[
\mathcal H_{I_-\cup J\cup I_+}
\cong
\bigoplus_{\alpha_-,\alpha_+}
\mathcal H_{i_-^{\alpha_-}}\otimes
\mathcal H_{j^{\alpha_-,\alpha_+}}\otimes
\mathcal H_{i_+^{\alpha_+}}.
\]
If item (iv) holds, substitute the chosen factorization
\[
\mathcal H_{j^{\alpha_-,\alpha_+}}
\cong
\mathcal H_{j_L^{\alpha_-,\alpha_+}}\otimes
\mathcal H_{j_R^{\alpha_-,\alpha_+}}
\]
to obtain the displayed inherited split and the stated action of the left and right union algebras. Nothing in the Schur-lemma argument itself forces item (iv); it is exactly the extra structural condition required to reproduce the same blockwise tensor pattern as Theorem~\ref{thm:markovcollar}.
\end{proof}

\begin{corollary}[Exact or controlled four-term null modular relation on an inherited strip model]\label{cor:n1}
Fix one finite-dimensional regulated strip model satisfying Proposition~\ref{prop:nullec}, and define
\[
\mathfrak M_{I_-:J:I_+}
:=
\left\{
\sigma:\ I(I_-:I_+\mid J)_\sigma=0
\right\},
\]
\[
\delta^{\mathrm M}_{I_-:J:I_+}(\varepsilon)
:=
\sup\left\{
\inf_{\sigma\in\mathfrak M_{I_-:J:I_+}}\|\rho-\sigma\|_1:
I(I_-:I_+\mid J)_\rho\le \varepsilon
\right\}.
\]
For any state \(\eta\) on this strip model, write
\[
\Delta K_J(\eta)
:=
K_{I_-\cup J\cup I_+}(\eta)-K_{I_-\cup J}(\eta)-K_{J\cup I_+}(\eta)+K_J(\eta).
\]

If the reference state \(\omega\) is exact Markov on this inherited decomposition, then
\[
\Delta K_J(\omega)\in Z(\mathcal A(J)),
\]
and in the canonical blockwise identification of Proposition~\ref{prop:nullec} one may take
\[
\Delta K_J(\omega)=0.
\]
Equivalently, without fixing that normalization there exists a central operator
\[
K_{\partial,J}(\omega)\in Z(\mathcal A(J))
\]
such that
\[
K_{I_-\cup J\cup I_+}(\omega)
=
K_{I_-\cup J}(\omega)+K_{J\cup I_+}(\omega)-K_J(\omega)+K_{\partial,J}(\omega).
\]

If instead \(I(I_-:I_+\mid J)_\omega\le \varepsilon\), then on this fixed strip model one may choose an exact Markov state \(\widetilde\omega_J\in \mathfrak M_{I_-:J:I_+}\) with
\[
\|\omega-\widetilde\omega_J\|_1
\le
\delta^{\mathrm M}_{I_-:J:I_+}(\varepsilon).
\]
Define the carried defect operator
\[
\mathfrak D_J(\omega,\widetilde\omega_J)
:=
\Delta K_J(\omega)-\Delta K_J(\widetilde\omega_J).
\]
Then
\[
K_{I_-\cup J\cup I_+}(\omega)
=
K_{I_-\cup J}(\omega)+K_{J\cup I_+}(\omega)-K_J(\omega)
+K_{\partial,J}(\widetilde\omega_J)
+\mathfrak D_J(\omega,\widetilde\omega_J),
\]
where
\[
K_{\partial,J}(\widetilde\omega_J):=\Delta K_J(\widetilde\omega_J)\in Z(\mathcal A(J)).
\]
For every bounded observable \(O\) in \(\mathcal A(I_-\cup J)\) or \(\mathcal A(J\cup I_+)\),
\[
\bigl|\omega(O)-\widetilde\omega_J(O)\bigr|
\le
\|O\|_\infty\,
\delta^{\mathrm M}_{I_-:J:I_+}(\varepsilon).
\]
If the marginals entering \(\Delta K_J\) for \(\omega\) and \(\widetilde\omega_J\) are uniformly faithful with lower spectral bound \(\lambda_\ast>0\), then
\[
\|\mathfrak D_J(\omega,\widetilde\omega_J)\|_\infty
\le
4\lambda_\ast^{-1}\,
\delta^{\mathrm M}_{I_-:J:I_+}(\varepsilon).
\]
Independently, the Fawzi--Renner recovery map supplies a constructive comparison state with trace-norm error
\[
r_{\mathrm{FR}}(\varepsilon)
=
2\sqrt{1-e^{-\varepsilon}}
\le
2\sqrt{\varepsilon}.
\]
Hence the exact four-term strip relation is available only at exact Markovity, or along a controlled strip family on one fixed inherited strip model for which
\[
\delta^{\mathrm M}_{I_-:J:I_+}(\varepsilon_J)\to 0.
\]
In the controlled case, \(\mathfrak D_J\) is carried explicitly at each finite stage.
\end{corollary}

\begin{proof}
Assume first that \(\omega\) is exact Markov on the inherited split of Proposition~\ref{prop:nullec}. Then the HJPW structure theorem gives, on each sector pair \(\alpha=(\alpha_-,\alpha_+)\),
\[
\omega_{I_-JI_+}|_\alpha
=
p_\alpha\,
\omega_{I_-\,j_L}^{(\alpha)}
\otimes
\omega_{j_R\,I_+}^{(\alpha)}.
\]
Accordingly,
\[
K_{I_-\cup J\cup I_+}(\omega)\big|_\alpha
=
(-\log p_\alpha)\mathbf 1
+
K_{I_-\,j_L}^{(\alpha)}\otimes \mathbf 1_{j_R I_+}
+
\mathbf 1_{I_- j_L}\otimes K_{j_R I_+}^{(\alpha)}.
\]
Tracing over \(I_+\), \(I_-\), or both gives the marginal block formulas
\[
K_{I_-\cup J}(\omega)\big|_\alpha
=
(-\log p_\alpha)\mathbf 1
+
K_{I_-\,j_L}^{(\alpha)}\otimes \mathbf 1_{j_R}
+
\mathbf 1_{I_- j_L}\otimes K_{j_R}^{(\alpha)},
\]
\[
K_{J\cup I_+}(\omega)\big|_\alpha
=
(-\log p_\alpha)\mathbf 1
+
K_{j_L}^{(\alpha)}\otimes \mathbf 1_{j_R I_+}
+
\mathbf 1_{j_L}\otimes K_{j_R I_+}^{(\alpha)},
\]
\[
K_J(\omega)\big|_\alpha
=
(-\log p_\alpha)\mathbf 1
+
K_{j_L}^{(\alpha)}\otimes \mathbf 1_{j_R}
+
\mathbf 1_{j_L}\otimes K_{j_R}^{(\alpha)}.
\]
Subtracting the last three expressions from the first gives zero on every block. Thus, in the canonical blockwise identification supplied by Proposition~\ref{prop:nullec},
\[
\Delta K_J(\omega)=0.
\]
If one chooses to keep the blockwise endpoint-label bookkeeping explicit rather than absorbing those constants into the canonical identification, the remainder is a direct sum of block constants and is therefore central. This yields the stated central form with \(K_{\partial,J}(\omega)\in Z(\mathcal A(J))\).

For the controlled statement, apply Proposition~\ref{prop:controlledmarkov} on this fixed finite-dimensional strip model, relabeling \(A,B,D\) there as \(I_-,J,I_+\). This yields an exact Markov replacement \(\widetilde\omega_J\in \mathfrak M_{I_-:J:I_+}\) with the displayed trace-norm bound. The identity
\[
\Delta K_J(\omega)=\Delta K_J(\widetilde\omega_J)+\mathfrak D_J(\omega,\widetilde\omega_J)
\]
is tautological, so the displayed four-term formula follows from the exact-Markov case applied to \(\widetilde\omega_J\). The observable estimate is immediate from
\[
\bigl|\omega(O)-\widetilde\omega_J(O)\bigr|
\le
\|O\|_\infty\,
\|\omega-\widetilde\omega_J\|_1.
\]
Finally, Proposition~\ref{prop:modulartransport}, again relabeled \(A,B,D\mapsto I_-,J,I_+\), yields the operator-norm bound on \(\mathfrak D_J\) whenever the relevant marginals have the stated lower spectral bound. The Fawzi--Renner estimate is the same one-shot recoverability bound used elsewhere and is independent of the exact-Markov replacement modulus.
\end{proof}

\begin{definition}[Renormalized null modular functional]
For a null interval \(I\) on generator \(\Omega\), let \(K_\partial(I,\Omega)\) denote the central endpoint-label term singled out by Corollary~\ref{cor:n1} on the inherited strip model, or by its controlled exact-Markov replacement when that model is used as reference. Define
\[
\widetilde K[I,\Omega]
:=
K[I,\Omega]-K_\partial(I,\Omega).
\]
For a null half-line \(H_a:=(a,\infty)\), abbreviate
\[
\widetilde K_a(\Omega):=\widetilde K[H_a,\Omega].
\]
\end{definition}

\begin{proposition}[Endpoint-Lipschitz null modular families and weak tail generator]\label{prop:n3}
Under the derived quasi-local propagation and bounded-interval endpoint control internal to Axiom~\ref{ax:maxent}, the renormalized interval family obeys the following matrix-element bounds on every fixed local-energy-bounded domain. For bounded intervals
\[
I=(a,b),\qquad
I'=(a',b'),
\]
contained in one compact endpoint window,
\[
\bigl|\langle \psi,(\widetilde K[I',\Omega]-\widetilde K[I,\Omega])\phi\rangle\bigr|
\le
C_{\psi,\phi,\Omega}\,\bigl(|a'-a|+|b'-b|\bigr).
\]
In particular the matrix elements of \(\widetilde K[I,\Omega]\) are jointly Lipschitz in the two endpoints and therefore have finite variation there.

For null half-lines \(H_a=(a,\infty)\), one likewise has
\[
\bigl|\langle \psi,(\widetilde K_{a'}(\Omega)-\widetilde K_a(\Omega))\phi\rangle\bigr|
\le
C_{\psi,\phi,\Omega}\,|a'-a|
\]
for \(a,a'\) in every compact window. Hence, for fixed \(\psi,\phi\),
\[
f_{\psi,\phi}(a):=\langle\psi,\widetilde K_a(\Omega)\phi\rangle
\]
is locally Lipschitz and therefore absolutely continuous. Its distributional derivative defines a locally \(L^\infty\) sesquilinear-form density
\[
\langle\psi,q(a,\Omega)\phi\rangle
:=
-\frac{1}{2\pi}\,\partial_a f_{\psi,\phi}(a),
\]
and for \(a<b\),
\[
\langle\psi,(\widetilde K_b(\Omega)-\widetilde K_a(\Omega))\phi\rangle
=
-2\pi\int_a^b \langle\psi,q(v,\Omega)\phi\rangle\,dv.
\]
The object \(q(a,\Omega)\) is the weak tail generator. At this stage it is not yet asserted to be represented by a local operator-valued density; its identification with the positive self-adjoint Borchers generator occurs only after the derived half-sided modular inclusion of Corollary~\ref{cor:n2} and Lemma~\ref{lem:translation}.
\end{proposition}

\begin{proof}
The bounded-interval endpoint-control statement derived from Axiom~\ref{ax:maxent} applies precisely to the renormalized family obtained after removal of the central endpoint term. Therefore, for intervals contained in one fixed compact endpoint window,
\[
\bigl|\langle\psi,(\widetilde K[I',\Omega]-\widetilde K[I,\Omega])\phi\rangle\bigr|
\le
C_{\psi,\phi,\Omega}\,|I'\Delta I|,
\]
where \(I'\Delta I\) is the symmetric difference of the two intervals. If
\[
I=(a,b),\qquad I'=(a',b'),
\]
then
\[
|I'\Delta I|
\le
|a'-a|+|b'-b|,
\]
which gives the displayed joint endpoint-Lipschitz bound.

For half-lines \(H_a=(a,\infty)\) and \(H_{a'}=(a',\infty)\), the symmetric difference is the bounded interval between the two endpoints, so
\[
|H_{a'}\Delta H_a|=|a'-a|.
\]
Because the central endpoint-label term has been removed in the definition of \(\widetilde K_a\), the same endpoint-control estimate yields
\[
\bigl|\langle\psi,(\widetilde K_{a'}(\Omega)-\widetilde K_a(\Omega))\phi\rangle\bigr|
\le
C_{\psi,\phi,\Omega}\,|a'-a|.
\]

Hence, on every compact \(a\)-interval, the function
\[
f_{\psi,\phi}(a):=\langle\psi,\widetilde K_a(\Omega)\phi\rangle
\]
is Lipschitz. A Lipschitz function on a compact interval is absolutely continuous, so it has an almost-everywhere derivative in \(L^\infty\) and obeys the fundamental theorem of calculus:
\[
f_{\psi,\phi}(b)-f_{\psi,\phi}(a)=\int_a^b \partial_v f_{\psi,\phi}(v)\,dv.
\]
Define
\[
\langle\psi,q(v,\Omega)\phi\rangle
:=
-\frac{1}{2\pi}\,\partial_v f_{\psi,\phi}(v)
\]
distributionally. Substituting this definition into the previous identity gives
\[
\langle\psi,(\widetilde K_b(\Omega)-\widetilde K_a(\Omega))\phi\rangle
=
-2\pi\int_a^b \langle\psi,q(v,\Omega)\phi\rangle\,dv.
\]
This is exactly the claimed weak-tail formula.
\end{proof}

\begin{theorem}[Downstream density-upgrade template]\label{lem:density}
Fix a null generator \(\Omega\) and the renormalized half-line family \(\widetilde K_a(\Omega)\) of Proposition~\ref{prop:n3}. For fixed vectors \(\psi,\phi\) in a common dense domain, let
\[
f_{\psi,\phi}(a):=\langle\psi,\widetilde K_a(\Omega)\phi\rangle,
\qquad
q_{\psi,\phi}(a):=\langle\psi,q(a,\Omega)\phi\rangle
=
-\frac{1}{2\pi}\,\partial_a f_{\psi,\phi}(a).
\]
Assume in addition that \(q_{\psi,\phi}\) is weakly differentiable in \(a\), that
\[
\lim_{a\to+\infty} f_{\psi,\phi}(a)=0,
\qquad
\lim_{a\to+\infty} q_{\psi,\phi}(a)=0,
\]
and define distributionally
\[
\langle\psi,p(a,\Omega)\phi\rangle
:=
-\partial_a q_{\psi,\phi}(a)
=
\frac{1}{2\pi}\,\partial_a^2 f_{\psi,\phi}(a).
\]
Then
\[
q_{\psi,\phi}(a)
=
\int_a^\infty \langle\psi,p(v,\Omega)\phi\rangle\,dv,
\]
and
\[
\langle\psi,\widetilde K_a(\Omega)\phi\rangle
=
2\pi\int_a^\infty (v-a)\,\langle\psi,p(v,\Omega)\phi\rangle\,dv.
\]
This statement is not part of the present fixed-cutoff bridge; it records the exact additional input/output package carried into the downstream density-upgrade task.
\end{theorem}

\paragraph{Null half-line blow-up net.}
Fix a smooth cut point and choose affine coordinate \(v\) on generator \(\Omega\) so that the cut sits at \(v=0\). For \(a\ge 0\), write
\[
H_a:=(a,\infty),
\qquad
\mathcal M_a(\Omega)
:=
\overline{\bigvee_{a<c<d<\infty}\mathcal A((c,d),\Omega)}.
\]
Thus \(\mathcal M_a(\Omega)\) is the half-line blow-up algebra generated by bounded null interval algebras inside \(H_a\). By isotony of the null interval net,
\[
a\le b
\quad\Longrightarrow\quad
\mathcal M_b(\Omega)\subseteq \mathcal M_a(\Omega).
\]

\begin{corollary}[Derived half-sided modular inclusion on null half-lines]\label{cor:n2}
On the OPH geometric scaling branch of Theorem~\ref{thm:bw}, the blow-up modular action near a smooth entangling cut acts on the null coordinate by
\[
v\mapsto e^{-2\pi t}v.
\]
Therefore, for every \(a\ge 0\),
\[
\sigma_t^\omega\bigl(\mathcal M_a(\Omega)\bigr)
=
\mathcal M_{e^{-2\pi t}a}(\Omega).
\]
Hence, for every \(a>0\),
\[
\sigma_t^\omega\bigl(\mathcal M_a(\Omega)\bigr)\subseteq \mathcal M_a(\Omega)
\qquad (t\le 0),
\]
so the inclusion
\[
\mathcal M_a(\Omega)\subset \mathcal M_0(\Omega)
\]
is half-sided modular. Equivalently, after the harmless convention change \(t\mapsto -t\), one recovers the standard positive-time half-sided-inclusion form. The former half-sided-inclusion premise is therefore derived from OPH null-interval structure, isotony, and the scaling-limit geometric action rather than imported as an additional modular-QFT ingredient.
\end{corollary}

\begin{proof}
By Theorem~\ref{thm:bw}, on every controlled collar family whose carried remainder vanishes in the refinement limit, the cap modular flow converges to the exact geometric cap dilation with \(2\pi\) normalization. Blow up the cap near a smooth cut and restrict to the chosen null generator. In that tangent limit, the cap-preserving flow becomes the null dilation \(v\mapsto e^{-2\pi t}v\). For every bounded interval \((c,d)\subset H_a\), the blow-up action therefore sends
\[
\mathcal A((c,d),\Omega)\longmapsto \mathcal A((e^{-2\pi t}c,e^{-2\pi t}d),\Omega).
\]
Taking the von Neumann closure of the algebra generated by all such intervals gives
\[
\sigma_t^\omega(\mathcal M_a(\Omega))
=
\mathcal M_{e^{-2\pi t}a}(\Omega).
\]
If \(t\le 0\), then \(e^{-2\pi t}a\ge a\), hence
\[
H_{e^{-2\pi t}a}\subseteq H_a.
\]
Isotony of the null interval net therefore implies
\[
\mathcal M_{e^{-2\pi t}a}(\Omega)\subseteq \mathcal M_a(\Omega).
\]
This is exactly the half-sided modular-inclusion property for the pair \(\mathcal M_a(\Omega)\subset \mathcal M_0(\Omega)\). Reparametrizing the modular parameter by \(t\mapsto -t\) gives the more common positive-time convention if desired.
\end{proof}

\begin{lemma}[Positive null-translation generator]\label{lem:translation}
For the derived half-sided modular inclusion
\[
\mathcal M_a(\Omega)\subset \mathcal M_0(\Omega)
\qquad (a>0)
\]
of Corollary~\ref{cor:n2}, let \(\Delta_0(\Omega)\) be the modular operator of the standard pair \((\mathcal M_0(\Omega),\omega)\) and define
\[
K_0(\Omega):=-\log \Delta_0(\Omega).
\]
Then Borchers--Wiesbrock yields a unique strongly continuous one-parameter unitary group
\[
U_\Omega(a)=e^{iaP_\Omega},
\qquad a\in\mathbb R,
\]
such that
\[
U_\Omega(a)\omega=\omega,
\qquad
U_\Omega(a)\,\mathcal M_b(\Omega)\,U_\Omega(a)^*=\mathcal M_{a+b}(\Omega)
\quad (a,b\ge 0),
\]
and whose generator \(P_\Omega\) is positive and self-adjoint on the Stone domain
\[
D(P_\Omega)
:=
\left\{
\psi\in\mathcal H:
\lim_{a\to 0}\frac{U_\Omega(a)\psi-\psi}{ia}\ \text{exists}
\right\}.
\]
Moreover,
\[
\Delta_0(\Omega)^{it}U_\Omega(a)\Delta_0(\Omega)^{-it}
=
U_\Omega(e^{-2\pi t}a),
\qquad
\Delta_0(\Omega)^{it}P_\Omega\Delta_0(\Omega)^{-it}
=
e^{-2\pi t}P_\Omega,
\]
and on the common invariant analytic core of \(K_0(\Omega)\) and \(P_\Omega\),
\[
[K_0(\Omega),P_\Omega]=-\,i\,2\pi P_\Omega.
\]

For each \(a\ge 0\), let \(\Delta_a(\Omega)\) be the modular operator of the translated standard pair \((\mathcal M_a(\Omega),\omega)\) and define
\[
K_a(\Omega):=-\log \Delta_a(\Omega).
\]
Then
\[
\Delta_a(\Omega)=U_\Omega(a)\Delta_0(\Omega)U_\Omega(a)^*,
\qquad
K_a(\Omega)=U_\Omega(a)K_0(\Omega)U_\Omega(a)^*,
\]
hence
\[
K_a(\Omega)=K_0(\Omega)-2\pi a\,P_\Omega
\]
as a quadratic-form identity on \(D(K_0(\Omega))\cap D(P_\Omega)\). Equivalently,
\[
\langle\psi,(K_b(\Omega)-K_a(\Omega))\phi\rangle
=
-2\pi(b-a)\,\langle\psi,P_\Omega\phi\rangle
\]
for all \(\psi,\phi\in D(K_0(\Omega))\cap D(P_\Omega)\).

In the canonical blockwise normalization of Corollary~\ref{cor:n1}, where the half-line endpoint term has been absorbed into the central part, the weak endpoint derivative of Proposition~\ref{prop:n3} agrees with the Borchers generator:
\[
\langle\psi,P_\Omega\phi\rangle
=
-\frac{1}{2\pi}\frac{d}{da}\Big|_{a=0^+}
\langle\psi,\widetilde K_a(\Omega)\phi\rangle
\]
for the same domain vectors. The half-sided-inclusion step proves only this affine half-line pair \((K_0(\Omega),P_\Omega)\); bounded-interval modular-Hamiltonian formulas remain downstream consequences of the separate interval-preserving projective branch recorded in Theorem~\ref{thm:stress}.
\end{lemma}

\begin{proof}
Corollary~\ref{cor:n2} gives the required half-sided modular inclusion. The standard Borchers--Wiesbrock theorem for a standard half-sided inclusion~\cite{wiesbrock93} then yields a unique strongly continuous unitary group \(U_\Omega(a)\) with \(U_\Omega(a)\omega=\omega\), the translation law for the nested half-line net, and a positive self-adjoint generator \(P_\Omega\); Stone's theorem gives the displayed domain formula.

The affine covariance relation
\[
\Delta_0(\Omega)^{it}U_\Omega(a)\Delta_0(\Omega)^{-it}
=
U_\Omega(e^{-2\pi t}a)
\]
is the Borchers relation. Differentiating at \(a=0\) gives
\[
\Delta_0(\Omega)^{it}P_\Omega\Delta_0(\Omega)^{-it}
=
e^{-2\pi t}P_\Omega.
\]
Since \(\Delta_0(\Omega)^{it}=e^{-itK_0(\Omega)}\), differentiating at \(t=0\) on the common analytic core yields
\[
-i[K_0(\Omega),P_\Omega]=-2\pi P_\Omega,
\]
hence
\[
[K_0(\Omega),P_\Omega]=-\,i\,2\pi P_\Omega.
\]

Because \(U_\Omega(a)\omega=\omega\) and
\[
\mathcal M_a(\Omega)=U_\Omega(a)\mathcal M_0(\Omega)U_\Omega(a)^*,
\]
the modular operator of the translated standard pair is
\[
\Delta_a(\Omega)=U_\Omega(a)\Delta_0(\Omega)U_\Omega(a)^*,
\]
so
\[
K_a(\Omega)=U_\Omega(a)K_0(\Omega)U_\Omega(a)^*.
\]
Differentiating this identity in \(a\) on \(D(K_0(\Omega))\cap D(P_\Omega)\) gives
\[
\frac{d}{da}K_a(\Omega)
=
i[P_\Omega,K_a(\Omega)]
=
-2\pi P_\Omega,
\]
because the commutator relation just proved is invariant under conjugation by \(U_\Omega(a)\). Integrating from \(0\) to \(a\) yields the displayed quadratic-form identity, and taking matrix elements gives the equivalent form. Proposition~\ref{prop:n3} had produced the weak endpoint derivative of the renormalized half-line family; in the canonical normalization of Corollary~\ref{cor:n1} the central term has been removed, so that weak derivative is exactly the Borchers generator. The final sentence is the explicit claim boundary: bounded intervals require the additional projective interval action not supplied by half-sided inclusion alone.
\end{proof}

\begin{theorem}[The OPH half-line generator is the local null-stress charge]\label{thm:stress}
Fix a null generator \(\Omega\) and work in the canonical blockwise normalization of Corollary~\ref{cor:n1}. Let \(\widetilde K_a(\Omega)\) denote the renormalized modular Hamiltonian of the half-line \(H_a\), and let \(P_\Omega\) be the positive Borchers generator of Lemma~\ref{lem:translation}. Then on
\[
D_\Omega:=D(K_0(\Omega))\cap D(P_\Omega)
\]
one has
\[
\langle\psi,P_\Omega\phi\rangle
=
-\frac{1}{2\pi}\frac{d}{da}\Big|_{a=0^+}
\langle\psi,\widetilde K_a(\Omega)\phi\rangle
\qquad (\psi,\phi\in D_\Omega),
\]
and this operator is exactly the local null-stress charge of the effective spacetime description on that same half-line family.

Equivalently, if the effective spacetime description writes the same renormalized family as
\[
\widetilde K^{\mathrm{eff}}_a(\Omega)
=
2\pi\,Q^{\mathrm{mod}}_a(\Omega)+K_{\partial}(a,\Omega),
\]
with \(Q^{\mathrm{mod}}_a(\Omega)\) the standard local null modular-energy term and \(K_{\partial}(a,\Omega)\) central and endpoint-supported, then
\[
Q_{kk}(\Omega)
:=
-\frac{1}{2\pi}\frac{d}{da}\Big|_{a=0^+}\widetilde K^{\mathrm{eff}}_a(\Omega)
\]
satisfies
\[
P_\Omega = Q_{kk}(\Omega)
\]
as a quadratic-form identity on \(D_\Omega\).

Hence the renormalized OPH null generator is not a separate EFT-side premise: it is the same operator that the effective spacetime description calls the local null-stress charge of the null half-line.
\end{theorem}

\begin{proof}
Lemma~\ref{lem:translation} proves the first displayed identity: in the canonical normalization of Corollary~\ref{cor:n1}, the weak endpoint derivative of the renormalized half-line family is exactly the Borchers generator. Passing to the effective spacetime description does not change the operator family; it rewrites the same \(\widetilde K_a(\Omega)\) in local continuum variables appropriate to the local Lorentzian regime. In that description, the local null-stress charge is precisely the right endpoint derivative
\[
-\frac{1}{2\pi}\frac{d}{da}\Big|_{a=0^+}\widetilde K^{\mathrm{eff}}_a(\Omega)
\]
of the same half-line modular family. Because \(\widetilde K^{\mathrm{eff}}_a(\Omega)\) and \(\widetilde K_a(\Omega)\) represent the same renormalized modular Hamiltonians, their right derivatives agree on \(D_\Omega\). Therefore the OPH generator and the effective null-stress charge coincide:
\[
P_\Omega = Q_{kk}(\Omega).
\]
The endpoint-supported central term does not alter the noncentral generator; in the canonical blockwise normalization it is absorbed into the central part fixed in Corollary~\ref{cor:n1}. This proves the identification.
\end{proof}

\paragraph{Imported input and remaining boundary.}
The only imported continuum input used above is the standard local null modular-Hamiltonian formula for the effective spacetime description of the same half-line family. No separate ``null-stress identification premise'' remains. What is still downstream is only:
\begin{enumerate}[label=(\roman*),leftmargin=2.2em]
\item transport of the half-line statement to bounded null intervals, which still requires the separate interval-preserving projective branch and its affine-covariant kernel \(g_I(v)\); and
\item reconstruction of a full local symmetric tensor from the directional charges, which is determined only up to the null-invisible metric term \(\phi g_{ab}\).
\end{enumerate}
So the gap closed here is the generator/charge identification itself: inside the null bridge and on the declared half-line family, \(P_\Omega\) is exactly the local null-stress charge.

\begin{lemma}[Null data determine the stress tensor modulo a metric term]\label{lem:nullreconstruct}
Let \(X_{ab}\) be a symmetric tensor. If
\[
X_{ab}k^ak^b=0
\]
for every null vector \(k\) at a point, then
\[
X_{ab}=\phi g_{ab}
\]
for some scalar \(\phi\).
\end{lemma}

\begin{proof}
Write
\[
X_{ab}=Y_{ab}+\phi g_{ab}
\]
with \(Y_{ab}\) traceless. Since \(g_{ab}k^ak^b=0\) on null vectors, the hypothesis implies \(Y_{ab}k^ak^b=0\) for all null \(k\). In local inertial coordinates let \(k=(1,\hat n)\) with \(|\hat n|=1\). Then
\[
Y_{ab}k^ak^b
=
Y_{00}+2\hat n^iY_{0i}+\hat n^i\hat n^jY_{ij}.
\]
Oddness under \(\hat n\mapsto-\hat n\) forces \(Y_{0i}=0\). Averaging the remaining relation over \(S^2\) gives
\[
Y_{00}+\frac13\delta^{ij}Y_{ij}=0.
\]
Tracelessness of \(Y_{ab}\) gives
\[
-Y_{00}+\delta^{ij}Y_{ij}=0,
\]
so both \(Y_{00}\) and the spatial trace \(\delta^{ij}Y_{ij}\) vanish. The remaining condition is therefore
\[
\hat n^i\hat n^j Y_{ij}^{\mathrm{TF}}=0
\qquad
\text{for all }\hat n\in S^2,
\]
where \(Y_{ij}^{\mathrm{TF}}\) is the traceless spatial part. The \(\ell=2\) spherical harmonics are linearly independent, so this implies \(Y_{ij}^{\mathrm{TF}}=0\). Hence \(Y_{ab}=0\) and therefore \(X_{ab}=\phi g_{ab}\).
\end{proof}

\begin{remark}[Claim boundary in the null bridge]
The c.12-local bridge ends with explicit theorem-level objects: the positive self-adjoint null-translation generator \(P_\Omega\) on its Stone domain, its Borchers affine covariance relation, the half-line modular identities above, and the exact half-line generator/charge identification of Theorem~\ref{thm:stress}. What remains downstream is only transport to bounded null intervals through the separate interval-preserving projective branch and the later tensor upgrade modulo the null-invisible metric term. The density-upgrade template of Theorem~\ref{lem:density} is recorded separately and does not carry the operator-identification burden.
\end{remark}
\subsection{Generalized entropy and the Einstein equation}

From this point through Corollary~\ref{cor:einstein} we specialize to the physical $d=4$ scaling branch. The claim boundary is explicit. The only standard geometric input used below, beyond the OPH outputs D3 and D4, is the fixed-volume area-variation identity for a small geodesic ball. The small-ball modular kernel used below is \emph{not} imported from a separate EFT law: it is the local expression of the OPH geometric cap generator $B_C$ from Theorem~\ref{thm:bw}, evaluated using the half-line null-stress identification of Theorem~\ref{thm:stress} together with the separate interval-preserving projective branch that transports it to bounded null intervals.

For a reference state $\omega$ and a small cap $C$,
\[
\delta S_C=\delta \langle K_C\rangle.
\]
By Theorem~\ref{thm:bw},
\[
\delta S_C = 2\pi\,\delta\langle B_C\rangle.
\]
By Proposition~\ref{prop:gensplit},
\[
S_{\mathrm{gen}}(C)=\mathrm{Tr}(\rho L_C)+S_{\mathrm{bulk}}(C).
\]
Together with Definition~\ref{def:newton}, this identifies the area term as the coarse-grained form of the edge entropy density rather than an independent postulate. After this edge-center split, the first-law piece entering the small-ball argument is the bulk contribution:
\[
\delta S_{\mathrm{bulk}}(C)=2\pi\,\delta\langle B_C\rangle.
\]

\begin{definition}[Admissible equilibrium variation]
Throughout this subsection, $\delta$ denotes a first-order state variation about a reference state $\omega$, within the cap-label-preserving MaxEnt class, at fixed cap $C$, fixed cap size/volume, and fixed conserved charges. Shape or null-cut deformations are denoted separately by $\delta_{\mathrm{shape}}$ or $\partial_\lambda$.
\end{definition}

\begin{assumption}[Fixed-cap generalized-entropy stationarity]\label{ass:equilibrium}
For admissible variations preserving cap labels, cap size/volume, and relevant conserved charges, the realized reference state is stationary for generalized entropy:
\[
\delta S_{\mathrm{gen}}(C)=0.
\]
\end{assumption}

This is the Jacobson-type fixed-cap stationarity premise used in this paper; it is not
derived from MaxEnt alone.

\begin{lemma}[Internal d=4 small-ball bridge]\label{lem:smallball}
Assume the hypotheses of Theorem~\ref{thm:stress}, together with the separate interval-preserving projective branch invoked there for bounded intervals. Let $p$ be the cap center and let $u^a$ be the future-directed unit tangent of the local diamond rest frame at $p$. In local inertial coordinates $(t,x^i)$ adapted to $u^a$, let $D_\ell$ be the causal diamond whose $t=0$ slice is the Euclidean ball
\[
B_\ell=\{t=0,\ r<\ell\},
\qquad
r^2=\delta_{ij}x^ix^j.
\]
Then
\[
\delta S_{\mathrm{bulk}}(C)
=
2\pi\int_{B_\ell}\frac{\ell^2-r^2}{2\ell}\,\delta\langle T_{00}\rangle\,d^3x
+
\delta\langle E^{(\eta)}_{C,\ell}\rangle.
\]
If, moreover, $\delta\langle T_{00}\rangle$ is approximately constant across $B_\ell$ and
\[
\delta\langle E^{(\eta)}_{C,\ell}\rangle=o(\ell^4)
\]
along the same scaling family, then in $d=4$,
\[
\delta S_{\mathrm{bulk}}(C)
=
\frac{8\pi^2\ell^4}{15}\,\delta\langle T_{00}\rangle
+O(\ell^5\partial T)+o(\ell^4).
\]
\end{lemma}

\begin{proof}
Work in the local Lorentzian scaling regime around $p$. On the geometric-subnet branch of Theorem~\ref{thm:bw}, the cap modular flow is the geometric flow of the diamond-preserving conformal Killing field. In the tangent diamond $D_\ell$, that field is
\[
\xi_{D_\ell}
=
\frac{1}{2\ell}
\Bigl((\ell^2-r^2-t^2)\,\partial_t-2t\,x^i\partial_i\Bigr),
\]
so on the $t=0$ slice one has
\[
\xi_{D_\ell}\!\cdot n=\frac{\ell^2-r^2}{2\ell},
\]
with $n^a=u^a$ the slice normal at the center.

The same kernel is the bounded-interval kernel supplied by that separate projective branch. Along each null generator of the diamond, Theorem~\ref{thm:stress} identifies the OPH half-line generator with the local null-stress charge, and the interval-preserving projective branch transports that identification to the affine-covariant interval weight that vanishes at the two diamond endpoints. On the $t=0$ slice that weight is exactly $(\ell^2-r^2)/(2\ell)$. Therefore the null-stress bridge reconstructs the bulk modular charge of the geometric generator as
\[
\delta\langle B_C\rangle
=
\int_{B_\ell}\frac{\ell^2-r^2}{2\ell}\,\delta\langle T_{00}\rangle\,d^3x
+\frac{1}{2\pi}\,\delta\langle E^{(\eta)}_{C,\ell}\rangle.
\]
Multiplying by the first-law factor $2\pi$ gives the displayed formula for $\delta S_{\mathrm{bulk}}(C)$.

If $\delta\langle T_{00}\rangle$ is approximately constant across $B_\ell$, then
\[
\int_{B_\ell}\frac{\ell^2-r^2}{2\ell}\,\delta\langle T_{00}\rangle\,d^3x
=
\frac{4\pi\ell^4}{15}\,\delta\langle T_{00}\rangle
+O(\ell^5\partial T).
\]
\noindent
The additional hypothesis $\delta\langle E^{(\eta)}_{C,\ell}\rangle=o(\ell^4)$ then yields
\[
\delta S_{\mathrm{bulk}}(C)
=
2\pi\cdot\frac{4\pi\ell^4}{15}\,\delta\langle T_{00}\rangle
+O(\ell^5\partial T)+o(\ell^4),
\]
which is exactly the stated coefficient. No separate EFT small-ball first law has been used: the kernel comes from the OPH geometric generator together with the D4 null-stress bridge.
\end{proof}

\begin{theorem}[Conditional Jacobson-type rest-frame relation from entanglement equilibrium]\label{thm:einstein}
Assume Axioms~\ref{ax:screen}--\ref{ax:entropy}, Assumption~\ref{ass:technical}, Assumption~\ref{ass:equilibrium}, and the hypotheses of Theorem~\ref{thm:stress} and Lemma~\ref{lem:smallball}. Then, for every sufficiently small cap centered at $p$ and every admissible first variation $\delta$ about a reference state $\omega$ in the locally Lorentzian $d=4$ scaling regime, if $u^a$ is the future-directed unit tangent of the local diamond rest frame at $p$, one has
\[
\delta\!\Bigl[(G_{ab}+\Lambda g_{ab})u^a u^b\Bigr]
=
8\pi G\,u^a u^b\,\delta\langle T_{ab}\rangle
\]
at $p$. In the adapted rest frame $u^a=(1,0,0,0)$, this is
\[
\delta\!\left(G_{00}+\Lambda g_{00}\right)=8\pi G\,\delta\langle T_{00}\rangle.
\]
\end{theorem}

\begin{proof}
Assumption~\ref{ass:equilibrium} gives the fixed-cap stationarity condition
\[
0=\delta S_{\mathrm{gen}}(C)
=
\delta S_{\mathrm{bulk}}(C)+\frac{\delta A}{4G}.
\]
By Lemma~\ref{lem:smallball},
\[
\delta S_{\mathrm{bulk}}(C)
=
\frac{8\pi^2\ell^4}{15}\,u^a u^b\,\delta\langle T_{ab}\rangle
+O(\ell^5\partial T)+o(\ell^4),
\]
where in the adapted rest frame the contraction is $\delta\langle T_{00}\rangle$.

For a fixed-volume small geodesic ball in $d=4$, the standard area-variation identity gives
\[
\delta A\big|_{V,\Lambda}
=
-\frac{4\pi\ell^4}{15}\,
\delta\!\Bigl[(G_{ab}+\Lambda g_{ab})u^a u^b\Bigr].
\]
Therefore
\[
0
=
\frac{8\pi^2\ell^4}{15}\,u^a u^b\,\delta\langle T_{ab}\rangle
-\frac{\pi\ell^4}{15G}\,
\delta\!\Bigl[(G_{ab}+\Lambda g_{ab})u^a u^b\Bigr]
+O(\ell^5\partial T)+o(\ell^4).
\]
Dividing by $\ell^4$ and taking the small-ball limit removes the $O(\ell^5\partial T)+o(\ell^4)$ terms and yields
\[
\delta\!\Bigl[(G_{ab}+\Lambda g_{ab})u^a u^b\Bigr]
=
8\pi G\,u^a u^b\,\delta\langle T_{ab}\rangle.
\]
This is precisely the desired rest-frame scalar relation.
\end{proof}

\begin{corollary}[Internal tensor upgrade in the scaling regime]\label{cor:einstein}
If the first-variation relation of Theorem~\ref{thm:einstein} holds for all local observer four-velocities and all reference states in a connected scaling branch, then
\[
G_{ab}+\Lambda g_{ab}=8\pi G\,\langle T_{ab}\rangle
\]
throughout that branch. The only local ambiguity is the metric term isolated by Lemma~\ref{lem:nullreconstruct}, and that ambiguity is exactly the same branch constant $\Lambda$.
\end{corollary}

\begin{proof}
Define
\[
Y_{ab}:=G_{ab}+\Lambda g_{ab}-8\pi G\,\langle T_{ab}\rangle.
\]
At any point $p$, Theorem~\ref{thm:einstein} gives
\[
u^a u^b\,\delta Y_{ab}=0
\]
for the unit future-directed four-velocity $u^a$ of the local diamond rest frame. By the Lorentz branch of Theorem~\ref{thm:bw} and overlap consistency across observers through the same point, these local rest frames exhaust all timelike directions. Hence
\[
u^a u^b\,\delta Y_{ab}=0
\qquad
\text{for every unit timelike }u^a.
\]

Choose local inertial coordinates at $p$ and write
\[
u^a=\gamma(1,v^i),
\qquad
|\vec v|<1,
\qquad
\gamma=(1-|\vec v|^2)^{-1/2}.
\]
Then
\[
0=u^a u^b\,\delta Y_{ab}
=
\gamma^2\Bigl(\delta Y_{00}+2v^i\delta Y_{0i}+v^iv^j\delta Y_{ij}\Bigr)
\]
for every $|\vec v|<1$. The quadratic polynomial in $v^i$ therefore vanishes on an open ball, so all of its coefficients vanish:
\[
\delta Y_{00}=0,\qquad
\delta Y_{0i}=0,\qquad
\delta Y_{ij}=0.
\]
Thus
\[
\delta Y_{ab}=0
\]
as a full tensor.

Along any path of admissible first variations inside the connected scaling branch, $Y_{ab}$ is therefore constant. The null bridge isolated the only purely local freedom as a metric term via Lemma~\ref{lem:nullreconstruct}; fixing that term on one maximally symmetric reference state determines the same branch constant $\Lambda$. Equivalently,
\[
Y_{ab}=0,
\]
so
\[
G_{ab}+\Lambda g_{ab}=8\pi G\,\langle T_{ab}\rangle
\]
throughout the branch.
\end{proof}

\subsection{The cosmological constant}

\begin{lemma}[Vacuum-Energy Blindness]\label{lem:vac}
For any null vector $k$ and any vacuum-energy contribution $T^{\mathrm{vac}}_{ab}=-\rho_{\mathrm{vac}}g_{ab}$,
\[
T^{\mathrm{vac}}_{kk}=0.
\]
\end{lemma}

\begin{proposition}[Structural Separation of $\Lambda$]\label{prop:lambda}
Local null-modular data fix $T_{ab}$ only up to $\phi g_{ab}$. Therefore $\Lambda$ is not determined by local overlap consistency and requires a global input.
\end{proposition}

\begin{corollary}[Cosmological Constant from Capacity]\label{cor:lambda}
If the global screen capacity is identified with the de Sitter static-patch entropy,
\[
N_{\mathrm{scr}}=S_{\mathrm{dS}}=\frac{A_{\mathrm{dS}}}{4G}=\frac{3\pi}{G\Lambda},
\]
then
\[
\Lambda=\frac{3\pi}{G N_{\mathrm{scr}}}.
\]
\end{corollary}

Proposition~\ref{prop:lambda} and Corollary~\ref{cor:lambda} are the intended local/global D6 package: local null-modular reconstruction leaves the Einstein equation determined only up to a metric term, and the global screen-capacity identification fixes the remaining constant. This is a structural aspect of the framework. The vacuum-energy problem is reorganized by separating local null data from the global capacity term.

\section{Gauge Reconstruction and Standard Model Structure}

\subsection{Compact gauge reconstruction}

At any fixed UV cutoff, edge-center completion equips collars with finitely many sector labels and fusion rules. These are fixed-cutoff collar data. The local MaxEnt / collar-mixing package established earlier controls only fixed-cutoff recoverability, modular-support localization, and carried error terms on the realized branch. It is logically separate from the question whether transportable edge sectors survive refinement, and it supplies no theorem that the refinement-limit sector category is trivial or nontrivial. A later nontrivial transportable colimit is therefore treated here as compatible additional gauge data rather than as something decided by the mixing estimate itself. Whenever compact gauge reconstruction is invoked, the entire transportable refinement-directed sector system is part of Assumption~\ref{ass:gaugeprem}.

Assume therefore that in the EFT regime there is a directed system \((\mathsf{Sect}_r)_r\) of transportable edge-sector categories with refinement functors, and write
\[
\mathsf{Sect}_\infty := \varinjlim_r \mathsf{Sect}_r
\]
for the directed colimit retaining the transportable sectors and intertwiners that persist in that system, with objectwise finite-dimensional fibers. This is the category on which Doplicher--Roberts / Tannaka reconstruction is applied. The theorem below is neutral about whether \(\mathsf{Sect}_\infty\) is trivial or nontrivial: if the transportable colimit of Assumption~\ref{ass:gaugeprem} furnished only the tensor unit, the reconstructed compact group would be the trivial group.

\begin{theorem}[Conditional compact gauge reconstruction in the bosonic branch]\label{thm:compactgauge}
Assume Axioms~\ref{ax:screen}--\ref{ax:entropy} and Assumption~\ref{ass:gaugeprem}. Let
\[
\mathcal F:\mathsf{Sect}_\infty\to\mathsf{Hilb}_{\mathrm{fd}}
\]
be the faithful bosonic fiber functor on that transportable colimit in the EFT regime. If \(\mathsf{Sect}_\infty\) is a rigid symmetric \(C^*\)-tensor category, then
\[
G:=\mathrm{Aut}_\otimes(\mathcal F)
\]
is a compact group and
\[
\mathsf{Sect}_\infty\simeq \mathrm{Rep}(G).
\]
In particular \(G\) is uniquely determined up to isomorphism.
\end{theorem}

\begin{proof}
The theorem begins only once the directed system of transportable sectors and its objectwise finite-dimensional fiber functor are given; those data are part of Assumption~\ref{ass:gaugeprem}, not consequences of the local-Gibbs or Dobrushin-type collar estimates. Fix a small skeleton of \(\mathsf{Sect}_\infty\), which does not change the tensor category or the fiber functor up to equivalence. For every object \(X\), any monoidal natural automorphism \(\eta\in\mathrm{Aut}_\otimes(\mathcal F)\) has a unitary component \(\eta_X\in U(\mathcal F(X))\) because \(\mathcal F\) is a \(^*\)-functor to finite-dimensional Hilbert spaces. Hence there is an injective map
\[
\mathrm{Aut}_\otimes(\mathcal F)\hookrightarrow \prod_{X}U(\mathcal F(X)),\qquad
\eta\mapsto (\eta_X)_X.
\]
For each intertwiner \(f:X\to Y\), naturality imposes the closed relation \(\mathcal F(f)\eta_X=\eta_Y\mathcal F(f)\). For each pair \(X,Y\), the monoidal structure isomorphism \(J_{X,Y}:\mathcal F(X)\otimes\mathcal F(Y)\xrightarrow{\sim}\mathcal F(X\otimes Y)\) imposes the closed relation
\[
\eta_{X\otimes Y}=J_{X,Y}\,(\eta_X\otimes\eta_Y)\,J_{X,Y}^{-1},
\qquad
\eta_{\mathbf 1}=\mathrm{id}_{\mathbb C}.
\]
Therefore \(\mathrm{Aut}_\otimes(\mathcal F)\) is a closed subgroup of the compact product \(\prod_X U(\mathcal F(X))\), so \(G=\mathrm{Aut}_\otimes(\mathcal F)\) is compact.

The remaining hypotheses required by Doplicher--Roberts / Tannaka reconstruction are now exactly checked on the declared sector category \(\mathsf{Sect}_\infty\): \(\mathsf{Sect}_\infty\) is rigid, symmetric, and \(C^*\) by assumption, and \(\mathcal F\) is a faithful bosonic fiber functor with finite-dimensional values by Assumption~\ref{ass:gaugeprem}. The standard reconstruction theorem therefore applies to the pair \((\mathsf{Sect}_\infty,\mathcal F)\) and yields a symmetric \(C^*\)-tensor equivalence
\[
\mathsf{Sect}_\infty \simeq \mathrm{Rep}(G),
\]
where an object \(X\) is sent to the representation \(g\mapsto g_X\) on \(\mathcal F(X)\) and morphisms are sent to the corresponding intertwiners \cite{dr89,dr90}. If the same category were also equivalent to \(\mathrm{Rep}(G')\) through another faithful bosonic fiber functor, then composing the two reconstruction equivalences would give a symmetric tensor equivalence \(\mathrm{Rep}(G)\simeq\mathrm{Rep}(G')\) compatible with the forgetful fiber functors. Applying the same reconstruction theorem to either side identifies the two automorphism groups of the fiber functor, so \(G\cong G'\). Thus the reconstructed compact group is unique up to isomorphism.
\end{proof}

If a fermionic sign object is present, the correct reconstruction statement is super-Tannakian rather than purely Tannakian. This paper does not prove that fermionic/chiral extension; it works only in the bosonic internal-gauge branch.

\subsection{Minimal admissibility and the Standard Model quotient}

Throughout this subsection, $\chi_{\mathrm{cpl}}$ means the dimension of the minimal coupled carrier supporting the required weak-type and color-type nonabelian charges on a common block. It is not the abstract minimal faithful representation dimension of the final gauge group.

\begin{definition}[Weak-type and color-type nonabelian roles]
A weak-type nonabelian charge is a pseudoreal nonabelian doublet role. A color-type nonabelian charge is an intrinsically complex nonabelian role, meaning that its nonabelian factor itself acts by an irreducible representation not equivalent to its conjugate. Abelian twisting of a pseudoreal nonabelian irreducible representation does not count as a color-type role.
\end{definition}

The derivation of the Standard Model quotient separates into five steps:
\begin{enumerate}[leftmargin=2em]
\item existence of a compact reconstructed group;
\item identification of the minimal nonabelian sector content;
\item identification of the minimal coupled carrier;
\item uniqueness of the connected abelian factor once admitted on that carrier;
\item determination of product gauge structure up to finite quotient, then of the global finite quotient from the realized matter spectrum.
\end{enumerate}

\begin{lemma}[Minimal nonabelian sector content]\label{lem:mincontent}
Any admissible low-energy sector must contain both a weak-type pseudoreal nonabelian charge role and an intrinsically complex color-type nonabelian charge role.
\end{lemma}

\begin{proof}
Light chiral matter requires left-handed multiplets and right-handed singlets to carry inequivalent gauge data. A pseudoreal doublet structure is the minimal way to realize the weak sector without immediately producing vectorlike masses. A genuinely complex nonabelian representation is required to distinguish quarks from antiquarks. A single nonabelian simple factor cannot simultaneously supply both a minimal weak-type role and an intrinsically complex color-type role; abelian twisting does not count toward the color-type requirement.
\end{proof}

\begin{lemma}[Connected 2D pseudoreal image classification]\label{lem:su2class}
Let $H$ be a connected compact group with a faithful irreducible $2$-dimensional pseudoreal unitary representation. Then the connected derived subgroup of its image is conjugate to $\SU(2)$. Equivalently, the nonabelian factor acting in that role is $\SU(2)$ up to finite central quotient.
\end{lemma}

\begin{proof}[Proof sketch]
Irreducibility places the image inside $\U(2)$ with scalar commutant. Its connected derived subgroup is therefore a connected compact subgroup of $\SU(2)$. The connected compact subgroups of $\SU(2)$ are either tori or $\SU(2)$ itself. The representation is nonabelian and pseudoreal, so the torus case is excluded. Hence the derived subgroup is conjugate to $\SU(2)$, with only a finite central kernel left invisible in the representation.
\end{proof}

\begin{lemma}[Connected 3D intrinsically complex image classification]\label{lem:su3class}
Let $H$ be a connected compact group with a faithful irreducible intrinsically complex $3$-dimensional unitary representation. Then the connected derived subgroup of its image is conjugate to $\SU(3)$. Equivalently, the nonabelian factor acting in that role is $\SU(3)$ up to finite central quotient.
\end{lemma}

\begin{proof}[Proof sketch]
Because the representation is irreducible on $\mathbb C^3$, the semisimple part of the Lie algebra of the image acts irreducibly on $\mathbb C^3$. A product of two nontrivial simple factors would force a tensor-product decomposition of dimension at least $4$, so only one simple factor can occur. Among compact simple Lie algebras, the only ones with nontrivial irreducible representations of dimension at most $3$ are $\mathfrak{su}(2)$ and $\mathfrak{su}(3)$. The $3$-dimensional $\mathfrak{su}(2)$ representation is real, not intrinsically complex, whereas the fundamental $\mathfrak{su}(3)$ representation is intrinsically complex. Therefore the connected derived subgroup is conjugate to $\SU(3)$, again up to finite central kernel.
\end{proof}

\begin{lemma}[Minimal coupled carrier under MAR]\label{lem:carrier}
Within the connected positive-dimensional Lie admissible class used by MAR, and under the weak-type / intrinsically complex color-type criteria above, the minimal coupled carrier supporting both roles on a common block has the form
\[
V=\mathbb C^3\otimes \mathbb C^2,
\qquad
\chi_{\mathrm{cpl}}=6.
\]
\end{lemma}

\begin{proof}
By Lemma~\ref{lem:su2class}, the minimal weak-type nonabelian role is the pseudoreal doublet of $\SU(2)$, hence has dimension $2$. By Lemma~\ref{lem:su3class}, the minimal intrinsically complex color-type nonabelian role is the triplet of $\SU(3)$, hence has dimension $3$.

By definition of $\chi_{\mathrm{cpl}}$, the weak-type and color-type roles must act nontrivially on one common irreducible nonabelian block rather than on disconnected direct-sum summands. Commuting irreducible actions on such a common block therefore require dimension at least the product of the minimal weak and color dimensions:
\[
\chi_{\mathrm{cpl}}\ge 2\cdot 3=6.
\]
Equality is achieved by the tensor-product carrier $\mathbb C^3\otimes\mathbb C^2$. The block-diagonal representation $\mathbb C^3\oplus \mathbb C^2$ of $S(U(3)\times U(2))$ is faithful of dimension $5$, but it is not coupled and hence is not the MAR minimizer. Connected non-semisimple counterexamples of $U(2)$ type are excluded by the explicit color-type criterion rather than being silently ignored.
\end{proof}

\begin{remark}\label{rem:connectedlie}
The dimension comparisons used in Lemma~\ref{lem:carrier} apply to the positive-dimensional connected Lie image carrying the admissible nonabelian charges, equivalently to the identity component of the reconstructed group on that sector. Finite, disconnected, or connected non-semisimple counterexamples exist, but they are outside the admissible class fixed by the explicit weak-type / color-type criteria.
\end{remark}

MAR supplies existence of a connected abelian factor; the next lemma proves only uniqueness and identification on the minimal coupled carrier.

\begin{lemma}[Uniqueness of the connected abelian factor on the minimal coupled carrier]\label{lem:commutant}
Inside $\U(6)$, the commutant of $\SU(3)\times\SU(2)$ acting on $\mathbb C^3\otimes \mathbb C^2$ is exactly $\U(1)$.
\end{lemma}

\begin{proof}
The $\SU(3)$ action is irreducible on the first tensor factor and trivial on the second. The $\SU(2)$ action is irreducible on the second tensor factor and trivial on the first. By Schur's lemma, any operator commuting with both actions is scalar on each irreducible factor. Hence the full commutant is a single $\U(1)$.
\end{proof}

\begin{theorem}[Product Gauge Structure up to Finite Quotient]\label{thm:sm}
Assume Axiom~\ref{ax:mar}, the hypotheses of Theorem~\ref{thm:compactgauge} on the ordinary or central-defect realized low-energy branch, the weak-type / color-type MAR premises above, and that the realized connected gauge image acts faithfully on the minimal coupled carrier. Then the realized connected gauge structure has the form
\[
G_{\mathrm{phys}}=
\frac{\SU(3)\times\SU(2)\times\U(1)}{\Gamma}
\]
for some finite central subgroup $\Gamma$.
\end{theorem}

\begin{proof}
Theorem~\ref{thm:compactgauge} yields some compact group $G$. Lemma~\ref{lem:mincontent} identifies the minimal admissible sector content, Lemmas~\ref{lem:su2class} and \ref{lem:su3class} identify the connected nonabelian factors realizing the minimal weak-type and color-type roles, and Lemma~\ref{lem:carrier} fixes the minimal coupled carrier. The tensor-product structure of that carrier implies commuting weak and color actions, so the connected semisimple part of the realized image is locally $\SU(3)\times\SU(2)$. Axiom~\ref{ax:mar} supplies the existence of one connected abelian charge factor acting nontrivially on the coupled carrier, and Lemma~\ref{lem:commutant} shows that any such connected abelian factor is necessarily a single $\U(1)$. A compact connected Lie group with Lie algebra $\mathfrak{su}(3)\oplus\mathfrak{su}(2)\oplus\mathfrak{u}(1)$ is a quotient of $\SU(3)\times\SU(2)\times\U(1)$ by a finite central subgroup. Faithfulness rules out an invisible extra torus, so the only remaining ambiguity is the finite central quotient $\Gamma$.
\end{proof}

\begin{theorem}[Hypercharge lattice on the realized matter package]\label{thm:hypercharge}
Assume gauge group $\SU(N_c)\times\SU(2)\times\U(1)_Y$, one generation of chiral matter $(Q,u^c,d^c,L,e^c)$, one Higgs doublet $H$, and Yukawa terms
\[
QHu^c,\qquad QH^\dagger d^c,\qquad LH^\dagger e^c.
\]
Then anomaly cancellation and Yukawa invariance determine the hypercharge ratios up to an overall \(\U(1)_Y\) normalization:
\[
Y_L=-N_cY_Q,\quad
Y_H=N_cY_Q,\quad
Y_u=-(N_c+1)Y_Q,\quad
Y_d=(N_c-1)Y_Q,\quad
Y_e=2N_cY_Q.
\]
Fixing the normalization by $Q=T_3+Y$ and $Q(\nu_L)=0$ gives
\[
Y_Q=\frac{1}{2N_c}.
\]
For $N_c=3$ one recovers the exact Standard Model lattice
\[
Y_Q=\frac16,\quad
Y_L=-\frac12,\quad
Y_u=-\frac23,\quad
Y_d=\frac13,\quad
Y_e=1,\quad
Y_H=\frac12.
\]
\end{theorem}

\begin{proof}
Yukawa invariance gives
\[
Y_u=-(Y_Q+Y_H),\qquad
Y_d=-Y_Q+Y_H,\qquad
Y_e=-Y_L+Y_H.
\]
The mixed anomalies yield
\[
N_cY_Q+Y_L=0,
\]
and
\[
2N_cY_Q+N_cY_u+N_cY_d+2Y_L+Y_e=0.
\]
Solving the first anomaly gives
\[
Y_L=-N_cY_Q.
\]
Substituting the Yukawa relations and $Y_L=-N_cY_Q$ into the mixed gravitational anomaly then yields
\[
Y_H=N_cY_Q,\qquad
Y_u=-(N_c+1)Y_Q,\qquad
Y_d=(N_c-1)Y_Q,\qquad
Y_e=2N_cY_Q.
\]
With these relations in place, the $\SU(N_c)^2\U(1)$ and $\U(1)^3$ anomalies cancel identically. The normalization of $Y_Q$ is fixed by the electric charge operator $Q=T_3+Y$ together with $Q(\nu_L)=0$, which implies
\[
Y_L=-\frac12=-N_cY_Q,
\qquad\text{hence}\qquad
Y_Q=\frac{1}{2N_c}.
\]
\end{proof}

\begin{corollary}[Three Generations on the realized MAR-admissible branch]\label{cor:ng}
Under the hypotheses of Theorem~\ref{thm:sm}, the asymptotic-freedom/CP-capability window together with the realized MAR selector gives
\[
N_g=3.
\]
\end{corollary}

\begin{proof}
Intrinsic CKM CP violation requires
\[
\frac{(N_g-1)(N_g-2)}{2}>0,
\]
hence $N_g\ge 3$. For one Higgs doublet, the one-loop weak-sector coefficient is
\[
b_2=\frac{11}{3}C_A-\frac{2}{3}\sum_f T(R_f)-\frac{1}{3}\sum_s T(R_s)
=
\frac{22}{3}-\frac{N_g(N_c+1)}{3}-\frac16
=
\frac{43}{6}-\frac{N_g(N_c+1)}{3}.
\]
Asymptotic freedom therefore implies
\[
N_g(N_c+1)<\frac{43}{2}.
\]
Since the color-type requirement forces $N_c\ge 3$, this gives $N_g\le 5$. The realized MAR selector then chooses $N_g=3$.
\end{proof}

\begin{corollary}[Three Colors on the realized MAR-admissible branch]\label{cor:nc}
Under the hypotheses of Theorem~\ref{thm:sm}, with \(N_g=3\) from Corollary~\ref{cor:ng}, Witten's global \(\SU(2)\) anomaly together with the realized MAR selector gives
\[
N_c=3.
\]
\end{corollary}

\begin{proof}
Witten's global $\SU(2)$ anomaly requires the total number of left-handed weak doublets to be even:
\[
N_g(N_c+1)\equiv 0 \pmod 2.
\]
With $N_g=3$, this implies $N_c$ is odd. The case $N_c=1$ is inadmissible because it removes the intrinsically complex color sector demanded by Lemma~\ref{lem:mincontent}. Minimal admissibility therefore selects $N_c=3$ \cite{witten82}.
\end{proof}

\begin{proposition}[Global quotient from the realized matter spectrum]\label{prop:z6}
If the realized matter spectrum carries the hypercharges of Theorem~\ref{thm:hypercharge} with $N_c=3$, then the subgroup of $\SU(3)\times\SU(2)\times\U(1)$ acting trivially on all realized states is exactly $\mathbb Z_6$.
\end{proposition}

\begin{proof}
Write $q:=6Y\in\mathbb Z$ and $\eta:=e^{i\pi/3}$, so the $\U(1)$ factor acts by $\eta^q$ on charge-$q$ multiplets. The element
\[
g_0:=(\omega_3,-1,\eta)
\]
acts trivially on the realized matter and Higgs multiplets:
\[
Q:\ \omega_3(-1)\eta=1,\qquad
u^c:\ \omega_3^{-1}\eta^{-4}=1,\qquad
d^c:\ \omega_3^{-1}\eta^2=1,
\]
\[
L:\ (-1)\eta^{-3}=1,\qquad
e^c:\ \eta^6=1,\qquad
H:\ (-1)\eta^3=1.
\]
Its powers therefore generate a subgroup of order six acting trivially on all realized states. Conversely, any central element acting trivially on $e^c$ must have $\U(1)$ part $\eta^n$, and triviality on $Q$ then fixes the $\SU(3)$ and $\SU(2)$ center factors uniquely. Hence every trivial central element is a power of $g_0$, so the kernel is exactly $\mathbb Z_6$.
\end{proof}

\begin{corollary}[Standard Model Gauge Group]
Under Theorem~\ref{thm:sm}, Theorem~\ref{thm:hypercharge}, Corollaries~\ref{cor:ng} and \ref{cor:nc}, and Proposition~\ref{prop:z6},
\[
G_{\mathrm{phys}}=
\frac{\SU(3)\times\SU(2)\times\U(1)}{\mathbb Z_6}.
\]
\end{corollary}

\begin{proof}
Theorem~\ref{thm:sm} yields the product structure up to finite quotient, and Proposition~\ref{prop:z6} fixes that quotient to $\mathbb Z_6$ once the realized hypercharge lattice and $N_c=3$ are in place.
\end{proof}

\begin{corollary}[Product-Group Consequences]
The adjoint representation of the full connected gauge group contains only
\[
(8,1,0)\oplus(1,3,0)\oplus(1,1,0).
\]
Equivalently, the adjoint representation of the derived gauge group contains only
\[
(8,1,0)\oplus(1,3,0).
\]
There are no mixed $(3,2,\pm 5/6)$ gauge bosons. Hence gauge-mediated proton decay is absent. The unbroken gauge symmetries forbid elementary gauge-boson mass terms: the photon remains massless, while gluons remain massless gauge fields of the unbroken color symmetry.
\end{corollary}

\begin{remark}
The same product-group structure removes the usual grand-unification monopole sector. In the present framework, unification of couplings does not proceed by embedding the gauge group into a simple group that is later broken. It proceeds through a shared geometric origin of the edge heat-kernel parameters. Absence of gauge-mediated proton decay and unification therefore coexist without tension. Any confinement statement is a separate infrared dynamical issue rather than a corollary of the product-group adjoint argument alone.
\end{remark}

\subsection{Fundamental scales}

The quantitative D10 branch is a forward calibration sector, not part of the recovered-core theorem package. Its external continuous particle-physics input is
\[
P \equiv a_{\mathrm{cell}}/\ell_P^2,
\]
together with the realized discrete gauge data of the structural branch, in particular
\[
N_c=3,
\qquad
\beta_{\mathrm{EW}}=N_c+1=4.
\]
The logical direction is therefore
\[
P
\longmapsto
\bigl(M_U(P),E_{\mathrm{cell}}(P)\bigr)
\longmapsto
\alpha_U(P)
\longmapsto
\bigl(t_U(P),t_{\mathrm{tr}}(P)\bigr)
\longmapsto
\bigl(t_2(P),t_3(P),v(P)\bigr)
\longmapsto
\alpha_i(\mu_\ast;P).
\]
The edge-law input on this lane is split explicitly by claim tier. The fixed-cutoff
same-overlap thermal/Casimir theorem is carried on the separate microphysics surface. This
compact paper imports that closure only through the declared merge boundary: one fixed-cutoff
edge-law package below, then the compact-group / Peter--Weyl lift used in the D10 pixel
constraint above it. The remaining open burden is refinement-stable universality/uniqueness of
that lift and the scalar-root control of the printed D10 solve, not the existence of the
finite-cutoff Casimir branch itself.
Measured electroweak quantities are not used to infer the internal transmutation parameters. They enter only at the end as compare-only validation targets for the printed D10 running/matching package.
The forward transmutation certificate on the live D10 lane makes that logical order explicit: the source-only basis reconstructs the same \(\alpha_U(P)\), hence the same unified diffusion parameter
\[
t_U(P)=4\pi^2\alpha_U(P),
\]
and transmutation exponent
\[
t_{\mathrm{tr}}(P)=\frac{2\pi}{(N_c+1)\alpha_U(P)}
\]
as the pixel-closure solve itself. Here the paper-side transmutation factor is the same
\(\beta_{\mathrm{EW}}=N_c+1\) used below; older overloaded \(\beta\)-ratios are compare-only
diagnostic readouts and are not part of the theorem contract.

The scales fixed directly from \(P\) are
\[
M_U(P)=\frac{E_P}{e^{2\pi}}\,P^{1/6},
\qquad
E_{\mathrm{cell}}(P)=\frac{E_P}{\sqrt P}.
\]
The factor \(e^{-2\pi}\) is the same modular normalization used on the Lorentz/BW side, and \(P^{1/6}\) is the cell-area scaling relation carried by the present D10 implementation.

The one-dimensional internal variable solved on the D10 branch is the unified coupling \(\alpha_U\). For a trial value of \(\alpha_U\), define
\[
v(\alpha_U,P)
:=
E_{\mathrm{cell}}(P)\,
\exp\!\left(-\frac{2\pi}{\beta_{\mathrm{EW}}\alpha_U}\right).
\]
Run the one-loop D10 couplings from
\[
\alpha_1(M_U)=\alpha_2(M_U)=\alpha_3(M_U)=\alpha_U
\]
down to a scale \(\mu\) using
\[
\alpha_i^{-1}(\mu;\alpha_U,P)
=
\alpha_U^{-1}
+
\frac{b_i}{2\pi}\log\frac{M_U(P)}{\mu},
\]
with the printed D10 one-loop coefficients \(b_i\). Let \(\mu_\ast=\mu_\ast(\alpha_U,P)\) be the fixed point determined by the tree-level \(Z\)-mass relation
\[
\mu_\ast
=
\frac{v(\alpha_U,P)}{2}\,
\sqrt{g_2(\mu_\ast)^2+g_Y(\mu_\ast)^2},
\qquad
\alpha_Y=\frac35\,\alpha_1.
\]
At that fixed point define
\[
t_2(\alpha_U,P)=4\pi^2\alpha_2(\mu_\ast;\alpha_U,P),
\qquad
t_3(\alpha_U,P)=4\pi^2\alpha_3(\mu_\ast;\alpha_U,P).
\]

Let \(\bar\ell_{\SU(2)}(t)\) and \(\bar\ell_{\SU(3)}(t)\) denote the nonabelian edge-entropy functions appearing in the D10 pixel constraint. The pixel-closure functional is
\[
\mathcal F(\alpha_U;P)
:=
\bar\ell_{\SU(2)}\!\bigl(t_2(\alpha_U,P)\bigr)
+
\bar\ell_{\SU(3)}\!\bigl(t_3(\alpha_U,P)\bigr)
-\frac{P}{4}.
\]
The printed D10 package takes \(\alpha_U(P)\) to be the branch value selected by solving
\[
\mathcal F(\alpha_U;P)=0.
\]
Only after that forward solve are the internal transmutation parameters fixed:
\[
t_U(P):=4\pi^2\alpha_U(P),
\qquad
t_{\mathrm{tr}}(P):=\frac{2\pi}{(N_c+1)\alpha_U(P)},
\]
and with them the downstream scale data
\[
t_2(P):=t_2(\alpha_U(P),P),
\qquad
t_3(P):=t_3(\alpha_U(P),P),
\qquad
v(P):=v(\alpha_U(P),P).
\]
Thus the branch runs from OPH input \(P\) to the internal transmutation data
\((t_U(P),t_{\mathrm{tr}}(P))\), then to the downstream local scale data
\((t_2(P),t_3(P),v(P))\), and only then to the low-energy couplings. It does not use measured
\(\alpha_i(m_Z)\) to infer those internal \(t\)-parameters.

Every quoted D10 electroweak number is downstream of this forward map. In particular, \(\alpha_{\mathrm{em}}^{-1}(m_Z)\), \(\sin^2\theta_W(m_Z)\), \(m_W\), and the stage-3 \(m_Z\) readout are forward-emitted outputs of the chosen D10 running/matching package. When compared with observation, they serve as validation of that printed implementation in the calibration-sector sense recorded in the premise ledger. What remains outside theorem level is a general existence/uniqueness or monotonicity proof for the scalar root selected by the printed D10 package under arbitrary convention changes.

\section{Relation to Holography and Existing UV Frameworks}

OPH belongs to the holographic family of ideas, but its primitive data are different from those used in asymptotic-boundary constructions. The distinction matters because several of the central claims of this paper, including the treatment of $\Lambda$, the handling of factorization, and the emergence of gauge structure, depend on it.

\subsection{Static-patch screen versus asymptotic boundary}

In asymptotic-boundary holography one starts from a dual theory at infinity and reconstructs the bulk inward. Here one starts from a finite-capacity screen equipped with a net of local patch algebras and reconstructs the effective bulk from overlap consistency. The primitive object is therefore a family of observer patches with nontrivial overlaps, not a single global boundary theory.

This shift has two technical consequences. First, subsystem factorization is treated from the beginning as a gluing problem with centers and sector labels. Second, the physical role of the horizon is local and operational: every observer has direct access only to a patch algebra, and global law is whatever survives reconciliation on overlaps.

\subsection{Positive $\Lambda$ and finite capacity}

The null-modular reconstruction of Section~\ref{sec:gravity} determines the stress tensor only up to a metric term. This means that local null data do not fix $\Lambda$. The cosmological constant enters only when one supplies the global screen capacity
\[
N_{\mathrm{scr}}=\log\dim\mathcal H_{\mathrm{tot}}.
\]
The resulting relation
\[
\Lambda=\frac{3\pi}{G N_{\mathrm{scr}}}
\]
fits naturally with a de Sitter static patch of finite entropy via $N_{\mathrm{scr}}=A_{\mathrm{dS}}/(4G)$. The local null-modular reconstruction leaves the metric ambiguity, while screen capacity supplies the global datum that fixes \(\Lambda\). OPH is therefore formulated in a de Sitter-first language. The positive cosmological constant is tied to finite capacity rather than to a deformation of a negative-$\Lambda$ starting point.

\subsection{Factorization, gauge structure, and the string sector}

The factorization problem of gauge theory and gravity is often treated as a technical nuisance. Here it is part of the architecture. Edge-center completion turns the collar center into a first-class object, and the entropy split
\[
S_{\mathrm{gen}}=\langle L_C\rangle + S_{\mathrm{bulk}}
\]
is then a structural consequence rather than an added prescription.

In the bosonic EFT branch, compact gauge reconstruction follows only after the explicit transportable sector-colimit and the bosonic ordinary/central gauge premises of Assumption~\ref{ass:gaugeprem} are supplied. The genuinely noncentral fixed-cutoff branch is handled separately by crossed-module data rather than by the ordinary compact-group theorem. The Standard Model quotient is then selected by minimal admissibility. This differs sharply from approaches in which the gauge group is specified in advance.

The string-theory relation is similarly reversed. The edge partition function
\[
Z_{\mathrm{edge}}(t)=\sum_R d_R^2 e^{-t C_2(R)}
\]
is the closed partition function obtained after gluing the open-edge weights $p_R(t)\propto d_R e^{-tC_2(R)}$. Peter--Weyl identifies this sum with the compact-group heat kernel at the identity, \(Z_{\mathrm{edge}}(t)=K_t(1)\), and the Chapman--Kolmogorov law gives the corresponding collar-sewing rule. If a large-$N_{\mathrm{edge}}$ regime exists, with $N_{\mathrm{edge}}$ distinct from the physical color rank $N_c=3$, its free energy admits a genus expansion. Here this is read as a conditional effective reorganization of edge degrees of freedom, not as a derived critical-superstring completion. Any further lift to critical superstring structure would additionally require worldsheet supersymmetry, critical dimension, modular invariance, and GSO projection beyond the OPH axiom chain.

\begingroup
\footnotesize
\begin{tabular}{@{}>{\raggedright\arraybackslash}p{0.24\textwidth}>{\raggedright\arraybackslash}p{0.32\textwidth}>{\raggedright\arraybackslash}p{0.34\textwidth}@{}}
\toprule
Feature & Asymptotic-boundary holography & OPH \\
\midrule
Primitive data & Boundary theory at infinity & Finite-capacity screen with observer-patch net \\
Factorization across cuts & Subtle and often indirect & Explicit edge-center completion with central labels \\
Status of $\Lambda$ & External background datum & Fixed globally by $N_{\mathrm{scr}}$ after local null reconstruction \\
Gauge group & Usually specified as part of the model & Reconstructed from edge sectors, then filtered by admissibility \\
String sector & Fundamental dual description & Fixed-cutoff heat-kernel edge law on the declared compact-gauge branch; any large-$N_{\mathrm{edge}}$ worldsheet reorganization is conditional, and any critical-superstring lift is conjectural \\
\bottomrule
\end{tabular}
\endgroup

\section{Claim Tiers and Falsifiability}

The recovered core is
\[
(D1\text{--}D5)\cup(D7\text{--}D9),
\]
namely the relativity chain together with the realized Standard Model structural chain. D6 is a separate input-dependent capacity corollary, D10 is the integrated calibration sector, and D12 collects phenomenological continuations.

\subsection{Prediction audit by claim tier}

\begingroup
\footnotesize
\begin{longtable}{@{}>{\raggedright\arraybackslash}p{0.20\textwidth}>{\raggedright\arraybackslash}p{0.31\textwidth}>{\raggedright\arraybackslash}p{0.39\textwidth}@{}}
\toprule
Tier & Representative outputs & What would actually falsify the tier \\
\midrule
\endhead
\bottomrule
\endfoot
Phase I recovered core (Theorem~\ref{thm:master}; D1--D5, D7--D9) & Confluence of overlap repair, conditional Lorentz kinematics, the conditional Jacobson-type Einstein branch in the stated scaling regime, compact gauge reconstruction in the bosonic branch, the realized Standard Model quotient chain, the exact hypercharge lattice on the realized matter package, the realized counting chain \(N_g=3\) then \(N_c=3\), and the product-group consequence of no gauge-mediated proton decay & A mathematical failure in the derivation chain, or data that require a different realized gauge quotient, different realized hypercharges, a different color or generation count on the admissible branch, or gauge-mediated proton decay \\
Phase II input-dependent corollary (D6) & \(\Lambda=\frac{3\pi}{G N_{\mathrm{scr}}}\) once the screen-capacity identification is made, while any retained capacity-level neutrino side estimate on that input is legacy bookkeeping rather than the live neutrino theorem lane & Failure of the screen-capacity identification or of the specific global capacity relation. Such a failure would not by itself erase the recovered core \\
Phase II calibration sector (D10) & forward transmutation data \(\alpha_U(P)\), \(t_U(P)\), \(t_{\mathrm{tr}}(P)\), pixel-closure gauge-coupling consistency, \(\alpha_i(m_Z)\), \(\alpha_{\mathrm{em}}^{-1}(m_Z)\), \(\sin^2\theta_W(m_Z)\), and the target-free source-only electroweak quintet \((W,Z,\alpha_{\mathrm{em}}^{-1},\sin^2\theta_W,v)\) on the D10 calibration branch & Failure of the integrated calibration closure once \(P\) and the printed running/matching conventions are imposed \\
Phase III phenomenological continuations (D12 and beyond) & Flavor ans\"atze, charged-lepton continuation ans\"atze beyond exact centered readback, texture branches, the weighted-cycle / Majorana-holonomy neutrino branch above the legacy D6 side estimate, further neutrino mass/mixing refinements, dark-sector response laws, baryogenesis estimates, proton-spin bookkeeping, proton-lifetime estimates beyond the gauge-channel exclusion, black-hole spectroscopy templates, conditional string/worldsheet reorganizations, conjectural critical-superstring extensions, and other downstream phenomenology & Retraction of the corresponding continuation only. These branches are not part of the recovered core and their failure is not, by itself, a falsification of relativity-plus-Standard-Model recovery \\
\end{longtable}
\endgroup

Phase I is the recovered-core tier, Phase II is adjacent but non-core, and Phase III contains phenomenological continuations. Interpretive strange-loop discussions lie outside these tiers: the paper does not define an OPH closure map on an invariant observer-supporting sector or prove uniqueness or stability for such a map.

\subsection{Recovered-core discriminators}

Once the claim boundary is made explicit, the sharpest discriminators are the ones that do not depend on auxiliary flavor, dark-sector, baryogenesis, or string ans\"atze.

\subsubsection{Realized Standard Model gauge structure}

On the realized MAR-admissible branch, this paper fixes
\[
G_{\mathrm{phys}}=
\frac{\SU(3)\times\SU(2)\times\U(1)}{\mathbb Z_6},
\qquad
N_g=3,
\qquad
N_c=3,
\]
with the exact one-generation hypercharge lattice. Here the realized counting chain is sequential: the CP/UV admissibility window and minimal admissibility yield \(N_g=3\), and Witten's \(\SU(2)\) anomaly parity with that realized generation count yields \(N_c=3\).
\[
Y_Q=\frac16,\qquad
Y_L=-\frac12,\qquad
Y_u=-\frac23,\qquad
Y_d=\frac13,\qquad
Y_e=1,\qquad
Y_H=\frac12.
\]
Evidence that the realized low-energy gauge structure or realized hypercharges differ from these values would directly contradict the recovered core.

\subsubsection{Product-group corollary}

Because the realized gauge structure is a product group up to the finite central quotient fixed above, there are no mixed \((3,2,\pm 5/6)\) gauge bosons. Hence
\[
\tau_p^{(\mathrm{gauge})}=\infty
\]
for gauge-mediated proton decay. Observation of proton decay specifically attributable to gauge-boson exchange would therefore falsify the realized product-group branch.

\subsubsection{Relativity branch}

The relativity claim is conditional but still sharp about its scope: if the stated geometric-modular, null-bridge, and fixed-cap-stationarity premises are realized in the intended scaling regime, then this paper predicts Lorentz kinematics on the screen and a Jacobson-type Einstein branch. What is being tested here is the existence of that scaling regime and its premises, not a fixed-cutoff matrix-algebra identity.

\subsection{What does \emph{not} count as a contradiction of the recovered core}

The following do \emph{not} falsify Theorem~\ref{thm:master} if they fail in their present form:
\begin{enumerate}[leftmargin=2em]
\item the uniform $\mathbb Z_6$ center-label ensemble and the associated $\varepsilon=1/6$ flavor ansatz;
\item charged-lepton continuation ans\"atze beyond the exact centered-readback / common-shift frontier, texture exponents, or legacy capacity-level neutrino side estimates;
\item modular-anomaly dark-sector response ans\"atze, including MOND/RAR-style scaling claims;
\item baryogenesis estimates based on extra suppression-counting assumptions; a baryogenesis theorem would still need a concrete out-of-equilibrium mechanism, derived CP-odd source terms, justified defect/sphaleron counting, washout control, and a freeze-out computation of the final asymmetry;
\item proton-spin ans\"atze or proton-lifetime estimates beyond the structural exclusion of gauge-mediated decay;
\item discrete-horizon spectroscopy templates;
\item worldsheet-like reorganizations requiring a controlled large-$N_{\mathrm{edge}}$ regime distinct from the proved compact-gauge chain, or any further critical-superstring lift.
\end{enumerate}
These are all post-Phase-I, non-core branches. Some are Phase-II input-dependent estimates, while the rest are Phase-III continuations. None is part of this paper's recovered relativity-plus-Standard-Model core.

\section{Common Objections and Clarifications}

\subsection{\texorpdfstring{Is the use of $P$ circular?}{Is the use of P circular?}}

The pixel area \(P\) is the external scale input for the separate gauge-coupling sector used here, and the circular step that read the internal D10 transmutation parameter back from measured low-energy couplings is omitted from this paper's presentation. Within the printed D10 implementation, \(P\) first fixes \(M_U(P)\) and \(E_{\mathrm{cell}}(P)\), then a one-dimensional pixel-closure solve fixes \(\alpha_U(P)\), equivalently the internal transmutation data \(t_U(P)\) and \(t_{\mathrm{tr}}(P)\), and only then do \(t_2(P)\), \(t_3(P)\), \(v(P)\), and the running electroweak outputs appear. Quantities algebraically entangled with that calibration remain on the Phase-II implementation surface: they are forward-emitted there, and comparison with observation checks that printed implementation rather than enlarging the recovered-core claim set.

Operationally, changing \(P\) moves the calibration family \((\alpha_U,\alpha_i(m_Z),\sin^2\theta_W(m_Z))\) through that forward solve. It does not alter the parameter-free structural outputs such as the gauge quotient or hypercharge lattice. Downstream matter-sector continuations, including charged-lepton continuation ans\"atze beyond the exact centered-readback / common-shift frontier, remain outside this SM/GR derivation paper's recovered-core theorem package.

\subsection{Does a fixed UV cell size break Lorentz invariance?}

The UV regulator is placed on the screen algebra, not on an emergent bulk spacetime lattice. Lorentz kinematics arise from the continuum modular action on caps, summarized in Corollary~\ref{cor:lorentz}, rather than from microscopic invariance of a bulk discretization.

This is conceptually similar to other continuum limits in statistical mechanics and lattice field theory. A regulator may break a symmetry at finite cutoff while the infrared fixed point restores it exactly. Here the restoration is stronger than a generic RG expectation because the modular-flow theorem identifies the symmetry group itself:
\[
\Conf^+(S^2)\cong \mathrm{SO}^+(3,1).
\]
Residual Lorentz violation would therefore have to arise from a failure of the regulator premises or from a failure of the modular-flow theorem, not merely from the existence of a UV cell size.
The remaining UV-side burden on the Lorentz branch is therefore the explicit scaling-limit plus geometric-subnet extraction-and-rigidity package under \textbf{T2}, not the mere existence of a finite cell size.

\subsection{Why use type-I algebras if continuum QFT uses type III?}

Type I is a regulator statement. At every finite cutoff the patch and collar algebras are finite-dimensional matrix algebras, so entropy, recovery maps, and
\[
K=-\log\rho
\]
are literal finite-dimensional objects. The Lorentz branch is \emph{not} a claim that the continuum observer algebra stays in that class. The refinement picture is
\[
\text{finite type-I regulator net}
\;\longrightarrow\;
\text{scaling-limit observer net } \mathcal A_\infty,
\]
and the limit can leave the regulator class. Axiom~\ref{ax:maxent} controls the realized state-side branch across refinement; it does \emph{not} by itself prove that the refinement limit remains type I or that the emitted scaling-limit pair is a vacuum or canonical cap state.

On the geometric-subnet branch of Theorem~\ref{thm:bw}, the scaling-limit cap algebras may be non-type-I and, in the continuum-QFT case of interest, are expected to be type III. Then there is generally no cap density matrix \(\rho_C\in \mathcal A_\infty^{\mathrm{geo}}(C)\) and no bounded operator \(K_C=-\log\rho_C\) inside the limit algebra. The correct continuum object is the modular automorphism group of the pair \((\mathcal A_\infty^{\mathrm{geo}}(C),\omega_\infty^{\mathrm{geo},C})\), which on that branch acts geometrically and typically outerly.

There is therefore no contradiction between using type-I regulators and aiming at a type-III continuum. The finite regulator provides the collar decomposition, recovery control, and carried errors. What the fixed-cutoff MaxEnt package adds is control of the realized state-side branch through one common finite-dimensional multiplier family under refinement. What it still does \emph{not} prove is the algebraic type of the continuum limit or that the limit state is a vacuum or canonical BW cap state. The phrase ``BW/canonical cap phase,'' when used in this paper, means only the case in which the realized refinement branch emits a scaling-limit geometric cap pair with the standard geometric modular automorphism group. In the generic non-type-I case that is an outer geometric modular action, not an inner density-matrix identity.

The clean internal extension route remains a two-step chain, and the first object is the pressure point. First, one needs a canonical scaling-limit observer cap-pair realization on the geometric subnet from transported marginals. The theorem contract for that object is explicit: the realized transported cap-local system packages the cap-local test family on \(\mathcal A_n^{\mathrm{geo}}(C_n)\), the projectively compatible transported marginal family, and the asymptotic transport-equivalence certificate. The remaining emitted witness is a derived vanishing carried-collar schedule on fixed local collars, explicitly decomposed into the constructive-recovery remainder and the faithful modular-defect term, namely the convergence to zero of the carried quantity
\[
r_{\mathrm{FR}}(\varepsilon_{n,m,\delta})
+
4\lambda_\ast^{-1}\,
\delta^{\mathrm M}_{A_{n,m,\delta}:B_{n,m,\delta}:D_{n,m,\delta}}
\bigl(\varepsilon_{n,m,\delta}\bigr)
\]
after transport to each fixed local collar model. Only after that can local weak-\(*\) extraction and GNS gluing emit the realized scaling-limit geometric cap pair.

Second, one needs ordered cut-pair rigidity on that realized limit. That stage would collapse the remaining cap-preserving conformal freedom to the unique BW hyperbolic subgroup, after which the modular KMS condition fixes the \(2\pi\) normalization. Ordered cut-pair rigidity remains symbolic until the realized scaling-limit geometric cap pair exists. Until those two objects are internalized, every stronger BW / vacuum / canonical identification remains explicitly theorem-external shorthand for the geometric-subnet branch of Theorem~\ref{thm:bw}, not a derived OPH theorem.

This is analogous to ordinary lattice field theory: one computes at finite regulator in a type-I algebra and then asks which continuum algebraic phase the scaling limit realizes. OPH's additional structural point is that the realized state-side branch is controlled under refinement. The open problem is to identify which scaling-limit observer algebra/state pair the realized branch produces and then to prove BW rigidity on that pair.
\subsection{UV branch and scaling-limit scope}

At fixed cutoff, Axiom~\ref{ax:maxent} yields a quasi-local finite-range interacting branch, and Propositions~\ref{prop:gauge} and \ref{prop:ancillauv} show that the physical normal form is unique only modulo quotient-level OPH-stable equivalence, while the terminal expectation functionals on the declared fixed-point / quotient-local physical algebras are already schedule-independent on that carrier even when microscopic representatives differ by gauge labels globally or by sector labels on one regional quotient-local glued state. The invariant fixed by the axiom language is the class \([\mathfrak U]_{\mathrm{OPH}}\) of the physical branch modulo implementation hiding and inert ancillary stabilization, not a unique microscopic representative. The genuinely noncentral topological case is also closed at fixed cutoff by the higher-gauge package of Theorems~\ref{thm:hgdatum}--\ref{thm:hgtransport} and Proposition~\ref{prop:hginteracting}, so the fixed-cutoff topological UV surface is closed on the ordinary, central-defect, and crossed-module branches alike. Beyond fixed cutoff, the continuum modular/geometric lift remains the geometric-subnet extraction-plus-rigidity problem under \textbf{T2}, and the transportable bosonic gauge lift still carries Assumption~\ref{ass:gaugeprem} when invoked.

\subsection{Why is charge quantization possible without a simple GUT?}

The framework uses the global quotient and anomaly structure rather than a simple-group embedding. Once the realized matter content is fixed, the subgroup acting trivially on all states is exactly $\mathbb Z_6$, and the anomaly equations force the sixth-integer hypercharge lattice. Integer electric charge for color singlets is then a property of the realized quotient structure. A simple unified gauge group is not required.

\section{Discussion}

The framework is strongest where the outputs are discrete or structurally rigid:
\begin{enumerate}[leftmargin=2em]
\item the Lorentz branch on the extracted prime geometric cap pair once the theorem-local BW lift objects are supplied;
\item the conditional scaling-limit Einstein branch under fixed-cap stationarity and the null modular bridge together with the bounded-interval projective branch;
\item the Standard Model gauge quotient chain;
\item exact hypercharges on the realized one-generation matter package with one Higgs doublet;
\item the realized counting chain \(N_g=3\) then \(N_c=3\) on the realized MAR-admissible branch;
\item the compact two-input quantitative layer, kept explicit as a secondary layer rather than as part of Phase I.
\end{enumerate}

The local MaxEnt branch gives a finite-range interacting fixed-cutoff dynamics, Propositions~\ref{prop:gauge} and \ref{prop:ancillauv} fix uniqueness of the UV branch only modulo OPH-stable equivalence, and Theorems~\ref{thm:hgdatum}--\ref{thm:hgtransport} with Proposition~\ref{prop:hginteracting} close the genuinely noncentral topological branch at fixed cutoff. The open items on the UV side are the continuum BW/geometric lift on the extracted prime geometric subnet under \textbf{T2} and the fuller transportable compact-gauge lift under Assumption~\ref{ass:gaugeprem}. The smallest internal extension route to the BW/geometric lift is the two-step chain of a scaling-limit cap-pair extraction theorem that extracts \((\mathcal A_\infty^{\mathrm{geo}}(C),\omega_\infty^{\mathrm{geo},C})\) as local weak-\(*\) limits of transported cap marginals on \(\mathcal A_n^{\mathrm{geo}}(C_n)\) with vanishing carried collar errors, followed by an ordered null cut-pair rigidity theorem that collapses the residual cap-preserving conformal freedom on those extracted cap pairs to the unique BW hyperbolic subgroup. Beneath the derived carried-collar witness, the sharper raw datum is a transported fixed-local-collar Markov/faithfulness package: collarwise conditional mutual information vanishing together with an eventual common floor on the finite modular-transport family on each fixed local collar model. Within that lower datum, the single missing clause is the modular-transport common floor, which feeds the faithful modular-defect term. No separate exact-Markov-reference faithfulness datum is missing: on one fixed collar model, the comparison marginals inherit the same eventual floor once the exact-Markov modulus goes to zero. The corpus does not derive a refinement-uniform common floor on the finite modular-transport marginals; this single clause remains external to the emitted theorem chain. That two-step chain and its transported fixed-local-collar datum are not part of the proved theorem surface stated here.

\paragraph{Shared excitation dictionary and flavor continuation status.}
The D10 calibration sector gives integrated gauge-coupling closure on the displayed carrier. That carrier closes its own exact \(W/Z\) chart, and the promoted target-free source-only repair theorem emits the public electroweak quintet on the Phase II calibration branch; the freeze-once coherent repair surface serves only as compare-only validation. On that frozen authoritative repair surface the canonical \(W/Z\) reference pair is hit exactly, but it remains a compare-only sidecar below the target-free theorem. The separate D6 cosmological-capacity relation remains an input-dependent corollary. On the flavor side the manuscript treats one shared OPH excitation dictionary as the only proof-facing family datum; lane-specific integers are not introduced independently. For each realized same-label refinement arrow \(e\) on the three-generation bundle, let \(\omega_e\) be the same-label overlap holonomy scalar, let
\[
d_e:=1-\omega_e
\]
be the derived edge defect, let \(g_e\) be the same-label gap scalar, and define
\[
q_e:=\sqrt{g_e d_e},
\qquad
\eta_e:=\log q_e-\frac13\sum_f \log q_f,
\qquad
\mu_e:=\frac{e^{\eta_e}}{\frac13\sum_f e^{\eta_f}}.
\]
These are the common descendants of overlap holonomy and edge-defect data. They exist once the same-label gap and overlap witnesses are fixed, and they are the shared family-side input used by the charged, quark, and neutrino continuations.

To include the charge-side bookkeeping explicitly, let \(s_e\) denote the realized sector-charge step carried by the same-label arrow \(e\). For any lane-specific readout \(\gamma\) on the realized same-label transport graph, let
\[
m_{\gamma,e}\in\mathbb Z_{\ge 0}
\]
be the multiplicity with which \(\gamma\) traverses the elementary same-label excitation arrow \(e\). Equivalently, along the realized transport path one has
\[
s_\gamma=\sum_e m_{\gamma,e}s_e.
\]
The associated excitation weight and excitation action are
\[
A_\gamma
:=
-\sum_e m_{\gamma,e}\log q_e
=
-\log\!\Bigl(\prod_e q_e^{\,m_{\gamma,e}}\Bigr),
\qquad
Y_\gamma^{\mathrm{exc}}
:=
\prod_e q_e^{\,m_{\gamma,e}}.
\]
So overlap holonomy, edge defects, gap data, and sector-charge transport determine the suppression law before any base choice is made. What earlier D12 language called a ``texture exponent'' or ``defect count'' is therefore not an extra primitive beyond the structural branch. The primitive flavor-side integers are the transport multiplicities \(m_{\gamma,e}\). A one-number exponent in ``units of \(\log 6\)'' is only the compressed summary
\[
n_\gamma^{(6)}:=A_\gamma/\log 6,
\]
and it becomes a literal defect count only after the additional uniform sixfold center-label collapse compatible with the realized
\[
\frac{\SU(3)\times\SU(2)\times\U(1)}{\mathbb Z_6}
\]
quotient. Without that extra collapse, the derived OPH data are the multiplicity vector \(m_{\gamma,\bullet}\) and the excitation action \(A_\gamma\), not an independently postulated flavor integer.

The imported realized-branch premise here is the D9 counting chain \(N_g=3\) then \(N_c=3\).
This shared dictionary feeds the downstream flavor lanes. The intrinsic neutrino lane factors through the same-label scalar certificate \((g_e,\omega_e)\) and its centered descendants \((q_e,\eta_e,\mu_e)\); the weighted-cycle branch therefore uses no extra primitive flavor labels. The quark mass-side continuation compresses the same dictionary to
\[
(\Delta_{ud}^{\mathrm{overlap}},\eta_Q^{\mathrm{centered}})
\]
and, on the compact same-family branch, to the emitted ray \(\mathcal R_{D12}^{ud}\). On the CKM/CP side the same-label left-handed orbit collapses to the singleton \(\sigma_{\mathrm{ref}}\), so the selected D12 sheet is fixed and is the wrong CKM sheet. The quark frontier is therefore layered. On the broader D12 continuation surface, the smaller honest theorem-side primitive is the light-quark isospin-breaking / overlap-defect scalar \(\Delta_{ud}^{\mathrm{overlap}}\): once that scalar is emitted, odd-budget neutrality fixes the light-sector pure-\(B\) payload pair and the \(\tau\)-pair algebraically on the ordered three-point family. On the emitted same-family ray, that same scalar is equivalently the downstream one-scalar law \ophid{quark_d12_t1_value_law}, with \ophid{intrinsic_scale_law_D12} retained only as the derived wrapper. The charged sector is equally tied to the same backbone: once a charged source pair \((\eta_{\mathrm{ch}},\sigma_{\mathrm{ch}})\) exists on the ordered charged carrier \((-1,x_2,1)\), the centered charged logs are fixed algebraically by
\[
\begin{aligned}
e_{\log,\mathrm{centered}}
&=
-\frac{(3+x_2)\sigma_{\mathrm{ch}}-\eta_{\mathrm{ch}}}{6},\\
\mu_{\log,\mathrm{centered}}
&=
\frac{x_2\sigma_{\mathrm{ch}}-\eta_{\mathrm{ch}}}{3},\\
\tau_{\log,\mathrm{centered}}
&=
\frac{(3-x_2)\sigma_{\mathrm{ch}}+\eta_{\mathrm{ch}}}{6}.
\end{aligned}
\]
so no separate charged-family integers remain to be chosen there.

The remaining burdens sit strictly above the shared dictionary. For charged leptons, the centered common-shift quotient is already closed negatively, and the remaining theorem-grade waiting set is the promotion
\[
\widehat C_e^{\mathrm{cand}}\longrightarrow \widehat C_e
\]
and then the post-promotion lift whose descended scalar is \(\mu_{\mathrm{phys}}(Y_e)\), with the physical identity-mode equalizer as the exact smaller forcing object beneath that scalar. For quarks, the broader honest continuation frontier is the light-quark isospin-breaking / overlap-defect scalar \(\Delta_{ud}^{\mathrm{overlap}}\); on the emitted D12 mass ray this is equivalently the downstream one-scalar law \ophid{quark_d12_t1_value_law}, with \ophid{intrinsic_scale_law_D12} as the derived wrapper. For neutrinos, the same-label scalar certificate closes beneath the weighted-cycle branch, the bridge-rigidity theorem emits
\[
C_\nu
=
G^2Q S^{-1/2}
=
0.9994295999075177,
\]
with emitted proxy
\[
P_\nu=6.699825740519345,
\]
while the induced paper-facing amplitude parameterization is
\[
B_\nu=P_\nu C_\nu=6.696004159297337.
\]
Thus the shared excitation dictionary is not the missing theorem object. It is the common OPH family base, while old base-6 exponents survive only as compare-only compressions of its excitation action.

\paragraph{Exact non-hadron output lane.}
For reader-facing exact hits, the non-hadron surface is explicit and lane-split. The structural carrier chain
\[
\text{Axioms }1\text{--}5 \longrightarrow D7\text{--}D9 \longrightarrow
\bigl(m_\gamma,m_g,m_{\mathrm{grav}}\bigr)=(0,0,0)
\]
is theorem-grade structural exactness. The electroweak calibration chain
\[
\text{Axioms }1\text{--}5 \longrightarrow D7\text{--}D10
\]
has an exact frozen authoritative repair sidecar
\[
(M_W,M_Z)=
(80.377,\ 91.18797809193725)\,\mathrm{GeV}
\]
on the frozen authoritative repair surface, but that exact pair remains compare-only beneath the target-free source-only theorem. The D11 Higgs/top branch carries a closed one-scalar forward seed with
\[
m_H=125.218922~\mathrm{GeV},
\qquad
m_t=172.388646~\mathrm{GeV}.
\]
The same D11 Jacobian also has a compare-only exact inverse slice at the canonical Higgs/top reference pair
\[
(m_H,m_t)=
(125.1995304097179,\ 172.3523553288311)\,\mathrm{GeV}.
\]
The charged lane has an exact same-family witness
\[
(m_e,m_\mu,m_\tau)=
(0.00051099895,\ 0.1056583755,\ 1.7769324651340912)\,\mathrm{GeV}
\]
on the same ordered eigenvalue family, but it remains same-family-only beneath the exact centered-readback / closed-common-shift frontier and the open theorem chain \(\widehat C_e^{\mathrm{cand}}\to \widehat C_e\) together with the post-promotion lift whose descended scalar is \(\mu_{\mathrm{phys}}(Y_e)\). The quark lane has an exact same-family witness
\[
(u,d,s,c,b,t)=
(0.00216,\ 0.00470,\ 0.0935,\ 1.273,\ 4.183,\ 172.3523553288311)\,\mathrm{GeV},
\]
again same-family-only, because the selected D12 sheet is the wrong CKM branch, the broader honest frontier is the light-quark overlap-defect scalar \(\Delta_{ud}^{\mathrm{overlap}}\), and on the emitted ray this is equivalently the downstream one-scalar law \ophid{quark_d12_t1_value_law}. The neutrino lane has the weighted-cycle bridge-rigid absolute family
\[
(m_1,m_2,m_3)=
(0.017454720257976796,\ 0.019481987935919015,\ 0.05307522145074924)\,\mathrm{eV},
\]
\[
\begin{aligned}
\Delta m_{21}^2&=7.488059465106851\times10^{-5}\,\mathrm{eV}^2,\\
\Delta m_{31}^2&=2.5123118727618473\times10^{-3}\,\mathrm{eV}^2,\\
\Delta m_{32}^2&=2.4374312781107786\times10^{-3}\,\mathrm{eV}^2.
\end{aligned}
\]
The exact two-parameter positive-segment adapter remains on disk only as a diagnostic sidecar beneath the bridge-rigidity invariant \(C_\nu\) and its induced paper-facing amplitude \(B_\nu\). The exact non-hadron output lane is complete, on a mixed theorem/compare-only/same-family surface rather than one uniform theorem tier.

\paragraph{Charged-lepton continuation status.}
The charged-lepton lane is not a theorem-level mass closure, and Koide-style fitting language does not describe its theorem boundary. What is exact is narrower. On the ordered charged carrier
\[
(-1,x_2,1),
\qquad
x_2=-0.5175863354681689,
\]
every charged source pair \((\eta,\sigma)\) determines the centered charged logs by
\[
e_{\log,\mathrm{centered}}
=
-\frac{(3+x_2)\sigma-\eta}{6},
\qquad
\mu_{\log,\mathrm{centered}}
=
\frac{x_2\sigma-\eta}{3},
\qquad
\tau_{\log,\mathrm{centered}}
=
\frac{(3-x_2)\sigma+\eta}{6}.
\]
So the centered charged hierarchy closes once a charged source pair is emitted.

The absolute frontier is explicitly two-layered. First, the promotion
\[
\widehat C_e^{\mathrm{cand}}
\longrightarrow
\widehat C_e
\]
waits on the upstream package \ophid{oph_generation_bundle_branch_generator_splitting}. Its sharp remaining clause is the descended-commutator statement \ophid{compression_descendant_commutator_vanishes_or_is_uniformly_quadratic_small_after_central_split}.
The corpus proves neither exact vanishing nor uniform quadratic smallness of that descended commutator, so \(\widehat C_e^{\mathrm{cand}}\) is not promoted.

Second, the common-shift quotient surface is closed negatively: the equalizer/common-shift route does not determine an absolute charged scale from centered data alone. The post-promotion slot is the refinement-stable uncentered trace lift. Its descended scalar is \(\mu_{\mathrm{phys}}(Y_e)\), and the exact smaller forcing object beneath that scalar is \ophid{charged_physical_identity_mode_equalizer}, namely the physical identity-mode equalizer \(\delta(r,r')=0\) on common physical fibers. If a theorem-grade \(\widehat C_e\) and refinement-stable physical \(Y_e\) are obtained, then the uncentered lift, determinant-line section, and affine anchor follow canonically from that one scalar:
\[
\widetilde C_e(Y_e)=\widehat C_e(Y_e)+\mu_{\mathrm{phys}}(Y_e)\,\mathbf 1,
\qquad
s_{\det}(Y_e)=3\mu_{\mathrm{phys}}(Y_e),
\qquad
A_{\mathrm{ch}}(Y_e)=\mu_{\mathrm{phys}}(Y_e).
\]
Hence the public claim boundary is exact same-carrier readback plus the closed common-shift no-go, with the remaining burdens carried first by the splitting package \ophid{oph_generation_bundle_branch_generator_splitting} and then by the post-promotion refinement-stable uncentered trace lift, with \ophid{charged_physical_identity_mode_equalizer} beneath \(\mu_{\mathrm{phys}}(Y_e)\). No current-corpus theorem emits public charged masses, \(g_e\), or \(\Delta_e^{\mathrm{abs}}\). If a future universal charged-mass law from the external calibration input \(P\) exists, it cannot enter through the present centered same-carrier readback surface, because that surface is common-shift invariant; it would have to descend through the still-missing affine object \(A_{\mathrm{ch}}(Y_e)=\mu_{\mathrm{phys}}(Y_e)\), equivalently through a theorem-grade bridge from the D10 calibration family to the charged determinant line. A target-anchored same-family exact witness reproduces the charged reference triple exactly on the same ordered eigenvalue family, and within that same-family scope the ordered three-point quadratic readout is closed through \ophid{oph_lepton_current_family_quadratic_readout_theorem}. On that same fixed target-anchored witness, the scoped affine coordinate also closes through \ophid{oph_lepton_current_family_affine_anchor_theorem}, with
\[
A_{\mathrm{ch}}^{(\mathrm{sf})}
=
\tfrac13\log\det(Y_e^{(\mathrm{sf})}),
\qquad
g_e^{(\mathrm{sf})}
=
e^{A_{\mathrm{ch}}^{(\mathrm{sf})}};
\]
this remains same-family-only and does not promote the charged theorem lane. Any balanced-circulant / \(Q=2/3\) / \(\delta=2/9\) relation remains compare-only continuation and is not an emitted mass theorem.
This second layer is a descent contract rather than a second fitted invariant: if theorem-grade \(\widehat C_e\) and an admissible refinement-stable uncentered lift exist, the physical identity-mode equalizer forces the descended scalar \(\mu_{\mathrm{phys}}(Y_e)\).

\paragraph{Quark continuation status.}
The light-quark lane is a D12 continuation branch and not a recovered-core theorem. The emitted same-label left-handed solver surface closes to the singleton \(\sigma_{\mathrm{ref}}\) (\ophid{sigma_ref}). That closure is negative: \(\sigma_{\mathrm{ref}}\) is the wrong D12 CKM sheet, and the corpus emits no sector-attached same-label left-handed branch-changing transport from that selected sheet to the physical CKM shell. Same-sheet overlap scans, chirality-swapped basis diagnostics, and other non-sector-attached orbit improvements are compare-only and do not alter the selected-sheet statement.

The first missing theorem object is the mass-side scalar law
\[
\Theta_{ud}^{\mathrm{mass}}
:=
\text{\ophid{quark_d12_t1_value_law}}.
\]
On the minimal light branch
\[
y_u=c_u\,\varepsilon^6,
\qquad
y_d=c_d\,\varepsilon^6,
\qquad
\varepsilon=\frac16,
\]
so the light-quark isospin-breaking scalar is
\[
\Delta_{ud}^{\mathrm{overlap}}
=
\frac16\log\frac{c_d}{c_u}.
\]
The emitted same-family D12 mass object is the one-parameter ray
\[
D12_{ud}^{\mathrm{mass}}
\equiv
\mathcal R_{D12}^{ud}
\qquad
\text{(paper id \ophid{D12_ud_mass_ray})},
\]
and on that emitted ray the same scalar is equivalently the one-scalar law \(\Theta_{ud}^{\mathrm{mass}}\), identifying the ray modulus with \(t_1\). The exact scalar identities are
\[
\Delta_{ud}^{\mathrm{overlap}}=\frac{t_1}{5},
\qquad
\log\frac{c_d}{c_u}=\frac65\,t_1,
\qquad
t_1=5\,\Delta_{ud}^{\mathrm{overlap}}=\frac56\log\frac{c_d}{c_u}.
\]
The full D12 mass-side package is then functorial:
\[
\eta_Q^{\mathrm{centered}}=-\frac{1-x_2^2}{27}\,t_1,
\qquad
\kappa_Q=-\frac{t_1}{54},
\qquad
x_2=-0.5175863354681689.
\]
The odd source package is also forced:
\[
\begin{aligned}
\beta_{u,\mathrm{diag},B}^{\mathrm{source}}&=\frac{t_1}{10},\\
\beta_{d,\mathrm{diag},B}^{\mathrm{source}}&=-\frac{t_1}{10},
\end{aligned}
\]
\[
\begin{aligned}
\texttt{source\_readback\_u\_log\_per\_side}
&=
\left(-\frac{t_1}{10},\,0,\,+\frac{t_1}{10}\right),\\
\texttt{source\_readback\_d\_log\_per\_side}
&=
\left(+\frac{t_1}{10},\,0,\,-\frac{t_1}{10}\right),
\end{aligned}
\]
and with the emitted spread pair \((\sigma_u,\sigma_d)=(5.5905,3.3049)\),
\[
\tau_u=\frac{\sigma_d}{10(\sigma_u+\sigma_d)}\,t_1,
\qquad
\tau_d=\frac{\sigma_u}{10(\sigma_u+\sigma_d)}\,t_1.
\]
The active builder-facing burden beneath this theorem is the pure-\(B\) payload pair; \ophid{intrinsic_scale_law_D12} is the derived wrapper above \(\Theta_{ud}^{\mathrm{mass}}\).

The second missing theorem object is the physical-sheet lift
\[
\Theta_{ud}^{\mathrm{phys}}
\]
to the same-label left-handed physical carrier
\[
\Sigma_{ud}^{\mathrm{phys}}
:=
\left\{
(\sigma_{\mathrm{id}},\tau,U_{u,L},U_{d,L},V_{\mathrm{CKM}},I_{\mathrm{CKM}})
:
V_{\mathrm{CKM}}=U_{u,L}^\dagger U_{d,L}
\right\}/\!\sim,
\]
where
\[
(U_{u,L},U_{d,L},V)\sim
(U_{u,L}D_u,\ U_{d,L}D_d,\ D_u^\dagger V D_d)
\]
for diagonal \(D_u,D_d\in U(1)^3\). Equivalently, this is a sector-attached section of the common-refinement frame-overlap quotient into \(\Sigma_{ud}^{\mathrm{phys}}\). The negative selector closure to \(\sigma_{\mathrm{ref}}\) forces this theorem: same-sheet rephasing preserves CKM moduli and therefore cannot move the selected branch to the physical shell.

The third missing theorem object is the target-free absolute readout
\[
\Theta_{ud}^{\mathrm{abs}}
\]
on the physical sheet. Its formula package is fixed on the current-family selected sheet. Define
\[
\sigma_{\mathrm{seed}}^{ud}:=\frac{\sigma_u+\sigma_d}{2},
\qquad
\eta_{ud}:=\frac{\sigma_u-\sigma_d}{2},
\]
\[
A_{ud}:=\frac{1}{2(1+\rho_{\mathrm{ord}}-x_2^2)},
\qquad
B_{ud}:=\frac{1}{2\!\left(1-x_2^2-\frac{x_2^2}{1+\rho_{\mathrm{ord}}}\right)},
\qquad
\rho_{\mathrm{ord}}=1.294284936377706.
\]
Then
\[
g_u
=
g_{\mathrm{ch}}
\exp\!\bigl(-(A_{ud}\sigma_{\mathrm{seed}}^{ud}-B_{ud}\eta_{ud})\bigr),
\qquad
g_d
=
g_{\mathrm{ch}}
\exp\!\bigl(-(A_{ud}\sigma_{\mathrm{seed}}^{ud}+B_{ud}\eta_{ud})\bigr).
\]
Once \((g_u,g_d)\) are emitted on the physical sheet, the ordered three-point quadratic readout emits the sextet algebraically.

The exact closure contract on this quark lane is therefore three-step:
\[
\Theta_{ud}^{\mathrm{mass}}
\Longrightarrow
\bigl(
\Delta_{ud}^{\mathrm{overlap}},
\eta_Q^{\mathrm{centered}},
\kappa_Q,
\tau_u,
\tau_d
\bigr),
\]
\[
\Theta_{ud}^{\mathrm{phys}}
\Longrightarrow
\sigma_{ud}^{\mathrm{phys}}
=
\bigl(\sigma_{\mathrm{id}},\tau,U_{u,L},U_{d,L},V_{\mathrm{CKM}},I_{\mathrm{CKM}}\bigr),
\]
\[
\Theta_{ud}^{\mathrm{abs}}
\Longrightarrow
(g_u,g_d)
\Longrightarrow
(m_u,m_d,m_s,m_c,m_b,m_t).
\]

\paragraph{Maximal theorem-emitted quark package on the present premise set.}
Let
\[
\begin{aligned}
P:={}&p_a\\
&+\text{(Axioms 1--5)}\\
&+\text{Assumption 6}\\
&+\text{the emitted OPH nodes D1--D10}\\
&+\text{the current public D12 quark objects}.
\end{aligned}
\]
Then the present corpus theorem-emits exactly the following quark-side package:
\[
\textbf{(i)}\quad
D12_{ud}^{\mathrm{mass}}
=
\mathcal R_{D12}^{ud}
=
\left\{
\lambda\left(\frac15,-\frac{1-x_2^2}{27}\right):\lambda\ge0
\right\},
\qquad
\lambda=\mathrm{ray\_modulus}=t_1,
\]
\[
\textbf{(ii)}\quad
\sigma_{ud}=\sigma_{\mathrm{ref}}
\]
on the emitted same-label left-handed local solver surface, with canonical token
\[
\texttt{D12::same\_label\_left::reference\_sheet}.
\]
This is a negative closure because \(\sigma_{\mathrm{ref}}\) is the wrong D12 CKM sheet.
\[
\textbf{(iii)}\quad
g_{\mathrm{ch}}=0.9231656602589082
\qquad
(\texttt{shared\_budget\_only}),
\]
\[
\sigma_u=5.5905,\qquad
\sigma_d=3.3049,\qquad
\sigma_{\mathrm{seed}}^{ud}=4.4477,\qquad
\eta_{ud}=1.1428,
\]
\[
A_{ud}=0.2467442927388686,
\qquad
B_{ud}=0.8125616987157023,
\]
\[
g_u=0.7797392875757557,
\qquad
g_d=0.12172551081512113
\qquad
(\texttt{current\_family\_only}).
\]
No stronger physical closure follows from this theorem-emitted package:
\[
P\nvdash
\Theta_{ud}^{\mathrm{mass}},
\qquad
P\nvdash
\Theta_{ud}^{\mathrm{phys}},
\qquad
P\nvdash
\Theta_{ud}^{\mathrm{abs}}.
\]
On the mass side, the emitted ray carries no theorem-grade nonzero value of its modulus \(t_1\). On the mixing side, same-sheet rephasing preserves CKM moduli, so the emitted selector value \(\sigma_{\mathrm{ref}}\) carries no sector-attached same-label left-handed lift to the physical CKM shell. On the absolute side, the affine sector-mean split is emitted only on the restricted current-family/shared-budget surfaces, so it carries no target-free physical-sheet pair \((g_u,g_d)\). Hence the exact same-family sextet and the D12 internal backread scalar package are sidecar representatives on their declared surfaces, not \(P\)-emitted physical-sheet theorem outputs.

\paragraph{Minimal-extension closure theorem.}
The exact minimal extension triple above that theorem-emitted package is
\[
H_{\mathrm{mass}}
\;:\;
\ell_{ud}:=\log\frac{c_d}{c_u},
\]
\[
H_{\mathrm{phys}}
\;:\;
s_{ud}^{\mathrm{phys}}:\mathfrak M_{ud}^{\mathrm{CR,phys}}\to \Sigma_{ud}^{\mathrm{phys}},
\]
\[
H_{\mathrm{abs}}
\;:\;
A_q^{\mathrm{phys}}:\Sigma_{ud}^{\mathrm{phys}}\to\mathbb R.
\]
With these three objects adjoined, the physical quark lane closes formally. First,
\[
t_1=\frac56\,\ell_{ud},
\]
so \(H_{\mathrm{mass}}\) emits \(\Theta_{ud}^{\mathrm{mass}}\) and therefore fixes
\[
\Delta_{ud}^{\mathrm{overlap}}=\frac{t_1}{5},
\qquad
\eta_Q^{\mathrm{centered}}=-\frac{1-x_2^2}{27}\,t_1,
\qquad
\kappa_Q=-\frac{t_1}{54},
\]
together with the odd package and the pure-\(B\) payload pair. Second,
\[
\Sigma_{ud}^{\mathrm{phys}}
:=
s_{ud}^{\mathrm{phys}}([F_0^\dagger F_1])
\]
emits \(\Theta_{ud}^{\mathrm{phys}}\). Third, writing
\[
\log g_u
=
A_q^{\mathrm{phys}}-(A_{ud}\sigma_{\mathrm{seed}}^{ud}-B_{ud}\eta_{ud}),
\qquad
\log g_d
=
A_q^{\mathrm{phys}}-(A_{ud}\sigma_{\mathrm{seed}}^{ud}+B_{ud}\eta_{ud}),
\]
\(H_{\mathrm{abs}}\) emits \((g_u,g_d)\), hence \(\Theta_{ud}^{\mathrm{abs}}\), and the ordered three-point quadratic readout then emits
\[
(m_u,m_d,m_s,m_c,m_b,m_t).
\]
So the exact closure route is
\[
P+H_{\mathrm{mass}}+H_{\mathrm{phys}}+H_{\mathrm{abs}}
\Longrightarrow
\Theta_{ud}^{\mathrm{mass}},\Theta_{ud}^{\mathrm{phys}},\Theta_{ud}^{\mathrm{abs}}
\Longrightarrow
(m_u,m_d,m_s,m_c,m_b,m_t).
\]

A target-anchored same-family exact witness reproduces the six running quark reference masses exactly on the same ordered three-point family,
\[
(u,d,s,c,b,t)=(0.00216,0.00470,0.0935,1.273,4.183,172.3523553288311)\,\mathrm{GeV},
\]
and within that same-family scope the selected-sheet exact readout closes through \ophid{oph_quark_current_family_selected_sheet_exact_closure}, built on the ordered three-point quadratic readout theorem \ophid{oph_quark_current_family_quadratic_readout_theorem}. That exact-hit sidecar does not move the selected D12 sheet and does not emit \(\Delta_{ud}^{\mathrm{overlap}}\), \ophid{quark_d12_t1_value_law}, or a target-free physical-sector pair \((g_u,g_d)\). The builder-facing pure-\(B\) payload pair beneath the reported quark values is a separate open object.
A separate D12 internal backread sidecar is also on disk. On that continuation-only surface, feeding the emitted reference-free forward light-quark pair through the coefficient-ratio identification \(c_u/c_d=y_u/y_d=m_u/m_d\) and the D12 normalization \(\Delta_{ud}^{\mathrm{overlap}}=\tfrac16\log(c_d/c_u)\) fixes
\[
\Delta_{ud}^{\mathrm{overlap}}=0.13431740757663183,\qquad
t_1=0.6715870378831591,
\]
\[
\eta_Q^{\mathrm{centered}}=-0.018210067243314802,\qquad
\kappa_Q=-0.012436796997836279.
\]
That sidecar closes the D12 mass-side scalar package numerically on its own continuation surface, but it does not replace the public theorem frontier and does not remove the wrong-sheet CKM boundary.

\subsection{Neutrino and hadron continuation status}

The isotropic intrinsic neutrino branch is ruled out by the exact atmospheric-splitting cap, so the live continuation lane is the weighted-cycle branch. On that branch the standard-math-fixed balanced transport-load selector is the arithmetic midpoint on the positive segment between \(\chi=1+\varepsilon\) and \(1+\gamma/2\):
\[
D_\nu=\frac{\chi+1+\gamma/2}{2},
\qquad
p_\nu(\gamma,\varepsilon)=1+\gamma+\frac{\varepsilon}{D_\nu}.
\]
This weighted-cycle branch repairs the oscillation-shape output to the physical PMNS/hierarchy regime on this continuation surface. On that lane the theorem-grade bridge object is the reduced invariant \(C_\nu\) above the emitted proxy \(P_\nu\), while \(B_\nu=P_\nu C_\nu\) is retained only as the paper-facing amplitude parameterization and the older positive-segment adapters remain diagnostic-only. It emits
\[
\theta_{12}=34.225904631810025^\circ,\qquad
\theta_{23}=49.72282845058266^\circ,\qquad
\theta_{13}=8.686355527700156^\circ,
\]
\[
\begin{aligned}
\delta_{\mathrm{PMNS}}&=305.58061231449796^\circ,\\
J&=-0.02753115613565372,\\
\frac{\Delta m_{21}^2}{\Delta m_{32}^2}&=0.030721110097966534.
\end{aligned}
\]
So the continuation branch reaches the observed oscillation window, fixes the normal-hierarchy shape, and leaves no hidden extra discrete selector on that lane.

The normalized same-label overlap-defect weight section closes beneath that weighted-cycle shape, and the paper-facing amplitude parameterization is
\[
B_\nu:=\lambda_\nu q_{\mathrm{mean}}^{p_\nu}/m_{\star,\mathrm{eV}}.
\]
The emitted internal proxy is
\[
P_\nu:=I_\nu^{1/2}\,\widehat{\mathrm{ratio}}^{\,1/2}\,\mathrm{sum\_defect}^{-1}=6.699825740519345.
\]
Writing
\[
G:=\mathrm{sum\_gap},
\qquad
Q:=\mathrm{prod\_qbar},
\qquad
S:=\mathrm{solar\_response\_over\_mstar},
\]
the weighted-cycle bridge rigidity theorem identifies the physical reduced bridge invariant with the distinguished family-assisted optimizer
\[
C_\nu
=
G^2\,Q\,S^{-1/2}
=
\mathrm{sum\_gap}^2\,\mathrm{prod\_qbar}\,\mathrm{solar\_response\_over\_mstar}^{-1/2}
=
0.9994295999075177.
\]
Hence
\[
\begin{aligned}
B_\nu&=P_\nu\,C_\nu=6.696004159297337,\\
\lambda_\nu&=\frac{m_{\star,\mathrm{eV}}}{q_{\mathrm{mean}}^{p_\nu}}\,P_\nu\,C_\nu=1.7237014208357415.
\end{aligned}
\]
If \(\hat m_i\) and \(\widehat{\Delta m_{ij}^2}\) denote the branch-fixed scale-free eigenmasses and splittings, then
\[
m_i=\lambda_\nu\,\hat m_i,
\qquad
\Delta m_{ij}^2=\lambda_\nu^2\,\widehat{\Delta m_{ij}^2},
\]
with branch-fixed coefficients
\[
\hat m
=
(0.010126301485273374,\ 0.011302414501969324,\ 0.03079142408841061),
\]
\[
\begin{aligned}
\widehat{\Delta m_{21}^2}&=2.520259180367675\times 10^{-5}\,\mathrm{eV}^2,\\
\widehat{\Delta m_{31}^2}&=8.455698156217034\times 10^{-4}\,\mathrm{eV}^2,\\
\widehat{\Delta m_{32}^2}&=8.203672238180266\times 10^{-4}\,\mathrm{eV}^2.
\end{aligned}
\]
Therefore the emitted absolute family is
\[
(m_1,m_2,m_3)
=
(0.017454720257976796,\ 0.019481987935919015,\ 0.05307522145074924)\,\mathrm{eV},
\]
\[
\begin{aligned}
\Delta m_{21}^2&=7.488059465106851\times10^{-5}\,\mathrm{eV}^2,\\
\Delta m_{31}^2&=2.5123118727618473\times10^{-3}\,\mathrm{eV}^2,\\
\Delta m_{32}^2&=2.4374312781107786\times10^{-3}\,\mathrm{eV}^2.
\end{aligned}
\]
The defect-weighted same-label edge family
\[
q_e=\sqrt{g_e d_e},
\qquad
\eta_e=\log q_e-\frac13\sum_f \log q_f,
\qquad
\mu_e=\mu_\nu\,\frac{e^{\eta_e}}{\frac13\sum_f e^{\eta_f}},
\]
remains the best closed constructive local object beneath the bridge, while the exact two-parameter positive-segment adapter, the bridge corridor, and the reduced-correction candidate audit are retained only as diagnostic surfaces and do not feed back into theorem state.
Taken together with the exact \(W/Z\) frozen repair pair, the closed Higgs/top forward seed together with its compare-only exact inverse slice, the exact charged/quark same-family witnesses, and the weighted-cycle bridge-rigid absolute neutrino family, the exact non-hadron output lane is complete, even though the theorem-grade charged and quark lanes remain open at their sharper remaining objects.

\paragraph{Local unification surface and exact-release frontier.}
The local quantitative bridge can be stated without inflating the recovered-core claim. On the
bosonic side, the declared pixel input \(P\) fixes the D10/D11 trunk
\[
P\longmapsto \alpha_U(P) \longmapsto \bigl(t_U(P),t_{\mathrm{tr}}(P)\bigr)\longmapsto v(P),
\]
after which the D10 source basis emits the public electroweak pair and the D11 forward seed
\[
\sigma_{D11,\mathrm{HT}}
=
\alpha_U(P)\cos(2\theta_{W0})/\sqrt{\pi}
\]
with no inverse readback of the internal transmutation data from measured low-energy couplings.
emits the Higgs/top branch. On the gravity side, the same D10 pixel law packages
\[
\ellbar_{\mathrm{SU(2)}}(t_{2,\mathrm{run}})
\;+\;
\ellbar_{\mathrm{SU(3)}}(t_{3,\mathrm{run}})
=
P/4,
\qquad
G=\frac{a_{\mathrm{cell}}}{4\,\ellbar(t)}.
\]
The exact local release burden consists of five explicit objects: one theorem/convention split for
the familiar-unit readout, one branch-preserving D10/gravity
edge-entropy bridge beneath the scalar identity
\[
\ellbar_{\mathrm{shared}}
=
\ellbar_{\mathrm{SU(2)}}(t_{2,\mathrm{run}})
\;+\;
\ellbar_{\mathrm{SU(3)}}(t_{3,\mathrm{run}}),
\]
one strict classical-regime clause above the BW/Lorentz honesty package, one target-free D10
chart-identity theorem for exact live \(W/Z\), and one live-forward D11 Higgs-promotion theorem.
On that declared extension surface the local release values are
\[
c=299792458\,\mathrm{m/s},
\qquad
G=6.674299995910528\times10^{-11}\,\mathrm{m^3\,kg^{-1}\,s^{-2}},
\]
\[
(M_W,M_Z,M_H)=
(80.377,\ 91.18797809193725,\ 125.1995304097179)\,\mathrm{GeV}.
\]
This is extension-level rather than current-corpus theorem status: \(c\) is structural from the
Lorentz branch rather than fitted from \(P\), and the displayed \(G\) row is an exact emitted
branch value rather than a literal zero-difference identity against the rounded benchmark
\(6.6743\times10^{-11}\).

The hadron lane remains separate. It depends on one production backend export bundle with publication-complete manifest provenance, real stable-channel correlator arrays, and later nonperturbative runtime/systematics completion rather than on a closed theorem-level mass surface. The current particle roadmap does not include hadron-mass closure, because that remaining nonperturbative execution burden is not a realistic target within the current project compute scope, so the neutrino weighted-cycle improvement does not promote hadron observables.

The paper does not address simulation narratives, observer-continuation mechanisms, or complete closure of all low-energy observables.\footnote{Questions outside this theorem package do not alter the specific claim set established here. Even a separate interpretive habitat theorem for OPH state-and-law data would not enlarge this SM/GR derivation paper's Phase-I scope unless the actual closure map, invariant admissible sector, and stability hypotheses were also proved at the same tier.}
\section{Conclusion}

Observer-Patch Holography starts from overlap consistency of observer patches and yields, within one framework, several effective structures usually treated separately in modern fundamental physics. This SM/GR derivation paper supports a unique schedule-independent normal form for overlap repair and hence a unique physical UV branch only modulo boundary redundancy, implementation hiding, and inert ancillary stabilization, a conditional Lorentz branch on the extracted prime geometric cap pair, a conditional Jacobson-type Einstein branch under fixed-cap stationarity and the null modular bridge together with the bounded-interval projective branch, the Standard Model gauge quotient with exact hypercharges on the realized matter package under the MAR admissibility package in the bosonic internal-gauge branch, the realized counting chain \(N_g=3\) then \(N_c=3\) on that branch, a separate integrated gauge-calibration sector, and a separate input-dependent cosmological-capacity corollary. Phenomenological continuations are discussed separately from the recovered core.

The framework therefore presents general relativity and the Standard Model as effective descriptions arising once the stated scaling-limit, geometric-subnet, and categorical premises are realized. The local quantitative claim surface is equally explicit: on a five-object extension surface, the same declared pixel input \(P\) controls the electroweak/Higgs mass trunk and the classical-regime gravity coupling bridge, while the invariant causal speed \(c\) is the structural Lorentz output with explicit SI readout. Open items are sharper quantitative control, full derivation of the refinement-limit bosonic ordinary/central gauge package of Assumption~\ref{ass:gaugeprem}, a fuller fermionic gauge branch, explicit microscopic realizations that emit the prime geometric cap pair and support ordered cut-pair rigidity, and proof of those five local exact-release objects on the live corpus.

\appendix
\input{tex_fragments/COMPACT_TECHNICAL_SUPPLEMENT_PORT.tex}

\begin{thebibliography}{99}
\input{tex_fragments/FOUNDATIONAL_REFERENCES.tex}
\end{thebibliography}

\end{document}
